\documentclass[11pt]{article}

\usepackage[margin=1in]{geometry}
\usepackage{amsmath,amssymb,mathtools,booktabs,microtype,xcolor}
\usepackage{array,enumitem}
\usepackage[colorlinks=true,linkcolor=blue!50!black,urlcolor=blue!50!black]{hyperref}

\title{Rotational Diagonal Kalman Linear Attention:\\
Model, Lessons from Mamba-2/Mamba-3, and Production Implementation}
\author{}
\date{September 4, 2026}

\newcommand{\R}{\mathbb{R}}
\newcommand{\ve}[1]{\boldsymbol{#1}}
\newcommand{\tr}{\mathsf{T}}
\newcommand{\sg}{\operatorname{sigmoid}}
\newcommand{\splus}{\operatorname{softplus}}
\newcommand{\detach}{\operatorname{stopgrad}}
\newcommand{\mean}{\operatorname{mean}}
\newcommand{\diag}{\operatorname{diag}}
\let\pathcode\path

\setlist{itemsep=2pt,topsep=4pt}
\setlength{\emergencystretch}{3em}

\begin{document}
\maketitle

\begin{abstract}
Mamba-2 supplies a strong selective state-space baseline, but its SISO state transition is a
scalar radial contraction: it cannot rotate a two-dimensional latent subspace.  The installed
Mamba-3 SISO implementation adds rotation economically by accumulating a bounded angular rate and
applying the same rotary transform to queries and keys.  Its most useful lesson is therefore not a
particular state-space write rule, but a representation: pairwise rotation can be moved out of the
recurrent state update and into a phase-rotated query/key basis.

Rotational diagonal Kalman linear attention (RDiag KLA, E2 clock) adopts that representation while
retaining KLA's diagonal Kalman covariance and innovation write.  It pair-ties radial transition
and process-noise parameters, uses a detached live radial statistic to set a bounded token/head
clock, scans phase in a documented negative gauge, and isolates the phase VJP from radial dynamics.
This note defines the model, the precise finite-precision implementation contract, and the tests
used for promotion.  The sealed E2 baseline implementation is commit
\texttt{c3346feaecae1486cdbc71f82b4edeacc4198950}; its immutable production snapshot passed the
complete one-H200 correctness and performance gate (117 core RDiag tests plus 54 E2 tests).
Distributed training-system checks and the multi-node soak/attestation remain before the planned
750M DiagKLA production submission.  The one-H200 result concerns implementation correctness and
measured kernel/layer performance, not the eventual language-model quality of a long training run.
A newer shared-DT investigation, including its unresolved production-qualification status, is
reported in Section~\ref{sec:shared-dt-status}.
\end{abstract}

\section{Scope and notation}

Let $B,T,H,K,V,D$ denote batch size, sequence length, number of heads, key dimension per
head, value dimension per head, and model width.  The rotary prefix has even width $K_r=2P$,
where $P$ is its number of adjacent pairs.  Tensors use the following shapes:
\begin{align*}
  \ve{x}&\in\R^{B\times T\times D},
  &\ve{q},\ve{k},\ve{\alpha},\ve{\omega}
    &\in\R^{B\times T\times H\times K},\\
  \ve{v}&\in\R^{B\times T\times H\times V},
  &\ve{e}&\in\R^{B\times T\times P},\\
  \ve{\tau}&\in\R^{B\times T\times H},
  &\ve{\delta},\ve{\phi}&\in\R^{B\times T\times H\times P}.
\end{align*}
The raw angle $e_{btj}$ is shared across heads; the clock $\tau_{bth}$ makes the resulting
increment head-specific.  For an adjacent pair $(2j,2j+1)$, define
\begin{equation}
  R(\vartheta)=
  \begin{bmatrix}
    \cos\vartheta&-\sin\vartheta\\
    \sin\vartheta& \cos\vartheta
  \end{bmatrix},
  \qquad R(\vartheta)^{\tr}R(\psi)=R(\psi-\vartheta).
  \label{eq:rotation}
\end{equation}
All equations below are per batch element unless a batch index is useful.  ``SISO'' follows the
Mamba implementation's terminology; a head can still carry multiple value channels that share
the same scalar state-space coefficients.

\section{What Mamba-2 has---and what it lacks}
\label{sec:mamba2}

For head $h$, value channel $p$, and SSM state coordinate $n$, the Mamba-2 selective SISO update
can be written
\begin{align}
  d_{th} &= \splus(\widetilde d_{th}+b^{d}_h),
  & a_{th}&=\exp(A_h d_{th}),\qquad A_h<0,\label{eq:m2-dt}\\
  X_{t,h,p,n}
    &=a_{th}X_{t-1,h,p,n}+d_{th}B_{t,n}u_{t,h,p},
  & y_{t,h,p}&=\sum_n C_{t,n}X_{t,h,p,n}+D_hu_{t,h,p}.
  \label{eq:m2-state}
\end{align}
The input projection is ordered as $[z,x,B,C,dt]$.  It emits one $dt$ row per token and head;
$B$ and $C$ are token-dependent state vectors.  The factor $a_{th}$ is a scalar radial contraction
over the head state, and $d_{th}B_tu_t$ is its discretized injection.  There is no angle projection,
phase state, adjacent-pair structure, sine/cosine operation, or $2\times2$ rotation in this path.
Mamba-2 therefore supplies selective decay and injection, but not oscillatory transport.  A pair
of real coordinates within a head shares the same scalar decay, but that factor is $a_{th}I$: it
has no skew component, phase state, or mechanism for rotating one coordinate into the other.

\section{The installed Mamba-3 SISO rotational path}
\label{sec:mamba3}

This section describes the code shipped in the installed \texttt{mamba-ssm 2.3.2.post1} package,
not an inferred variant.  Its input projection is ordered
\begin{equation}
  [z,x,B,C,dd_{dt},dd_A,\mathrm{trap},\mathrm{angle}].
\end{equation}
The $dd_{dt}$, $dd_A$, and trapezoid fields each have $H$ rows.  In contrast, the angle field has
only $P_M$ rows, where $P_M$ is half of the selected half- or full-rotary state prefix.  The raw
angle rows are expanded across heads.  Thus Mamba-3 separates a token/head clock from token/pair
angular content:
\begin{align}
  DT_{th} &= \splus(dd_{dt,th}+b^d_h),\\
  A_{th} &= \min\{-\splus(dd_{A,th}),-A_{\rm floor}\},
  & \rho_{th}&=\exp(A_{th}DT_{th}),\\
  \nu_{tj} &= \pi\tanh(e_{tj}),\\
  \Delta\theta_{thj} &= DT_{th}\nu_{tj},\\
  \theta_{thj} &= \theta_{0,hj}+\sum_{s=1}^{t}\Delta\theta_{shj}pmod{2\pi}.
  \label{eq:m3-phase}
\end{align}
The cumulative phase is positive.  Its Triton loop accumulates in FP32, but the combined training
wrapper casts raw angles to BF16 and materializes phase and final carry in that input dtype; this is
not an FP32-storage guarantee.  For every selected adjacent pair, the SISO kernel applies the
same $R(+\theta_{thj})$ to $Q_t$ and $K_t$.  Consequently the same-token dot product is unchanged,
whereas a cross-token term obtains the relative rotation
\begin{equation}
  (R(\theta_t)q_t)^{\tr}R(\theta_s)k_s
  =q_t^{\tr}R(\theta_s-\theta_t)k_s.
  \label{eq:m3-relative}
\end{equation}
If $S_t$ denotes a state in a fixed physical frame, this is equivalent to the moving-frame
recurrence
\begin{equation}
  S_t^- = \rho_t R(-\Delta\theta_t)S_{t-1},
  \label{eq:m3-moving-frame}
\end{equation}
pair by pair, with scalar radial decay $\rho_t$.  The identity follows from
$R(\theta_t)^{\tr}R(\theta_{t-1})=R(-\Delta\theta_t)$.  It is the key computational idea: no dense
rotational state transition is needed when radial coefficients are pair-isotropic.

Mamba-3 additionally uses a trapezoidal SSM write.  With
\begin{equation}
  \gamma_t=DT_t\,\sg(\mathrm{trap}_t),
  \qquad
  \eta_{t+1}=DT_{t+1}\bigl(1-\sg(\mathrm{trap}_{t+1})\bigr),
\end{equation}
the offline kernel scales token $t$'s rotated key by
$\gamma_t+\eta_{t+1}$, while its same-token contribution uses $\gamma_t$.  At a streaming boundary,
the cached previous key/value contributes the missing $\eta_t$ term.  This is a two-endpoint
discretization of the Mamba write; it is not the KLA innovation update defined below.

The angle clock is also coupled in backward.  The VJP of the positive cumulative phase is a
reverse cumulative sum.  It produces both
\begin{align}
  \overline e_{tj}
    &=\sum_h\overline{\Delta\theta}_{thj}\,DT_{th}\,
      \pi\bigl(1-\tanh^2(e_{tj})\bigr),\\
  \overline{DT}^{\mathrm{angle}}_{th}
    &=\sum_j\overline{\Delta\theta}_{thj}\,\pi\tanh(e_{tj}).
\end{align}
Here the sum over $h$ is the backward of the head expansion of the shared raw angle rows.  The
combined implementation explicitly adds $\overline{DT}^{\mathrm{angle}}$ to the direct
trapezoidal-write $DT$ cotangent; ordinary autograd subsequently adds the
$A\,\overline{ADT}$ contribution through $ADT=A\,DT$.  This is a coherent choice for a shared
physical discretization clock, but it is intentionally not the RDiag E2 gradient contract.
The Mamba-2/Mamba-3 comparison here is architectural, not a matched causal rotation ablation:
Mamba-3 also removes Mamba-2's convolutional path, makes $A$ token-dependent, normalizes and
biases $B/C$, and changes the write discretization.

\section{RDiag KLA with the E2 bounded clock}
\label{sec:rdiag}

\subsection{Pair-tied radial dynamics}

Let the radial preactivation be $f_{bthk}$ with bias $b^d_{hk}$, and let
\begin{equation}
  d_{bthk}=\splus(f_{bthk}+b^d_{hk}).
\end{equation}
For rotary pair $j$, define the pair step and tied process noise
\begin{align}
  \bar d_{bthj}
    &=\tfrac12(d_{bth,2j}+d_{bth,2j+1}),\\
  \bar\omega_{bthj}
    &=\tfrac12(\omega^{\rm in}_{bth,2j}+\omega^{\rm in}_{bth,2j+1}),\\
  \alpha_{bth,2j}=\alpha_{bth,2j+1}
    &=\exp\!\left[-\exp(A^{\log}_h)\bar d_{bthj}\right],\\
  \omega_{bth,2j}=\omega_{bth,2j+1}&=\bar\omega_{bthj}.
  \label{eq:pair-radial}
\end{align}
Only the selected prefix is pair-tied; tail lanes retain the diagonal-KLA frontend's original
$\alpha$ and $\omega$.  Pair tying is not cosmetic.  On each pair it makes
$\diag(\alpha_t)=\alpha_{tj}I_2$ and the diagonal process covariance
$\diag(\omega_t)=\omega_{tj}I_2$, so the \emph{prediction} terms commute with
$R(\vartheta)$.  It does not force posterior information or the Kalman gain to be pair-isotropic;
an observation generally splits them.  The diagonal posterior approximation is defined in the
current moving frame, as made explicit below.

\subsection{Detached live radial clock and persistent anchor}

Set the open upper clock bound to $C=0.1$.  At construction, snapshot one FP32 anchor per head,
\begin{equation}
  a_h=\mean_k\splus(b^d_{hk}),\qquad 0<a_h<C.
  \label{eq:anchor}
\end{equation}
The anchor is a persistent buffer, not a parameter.  At each token compute the detached live radial
mean
\begin{equation}
  m_{bth}=\mean_k\splus\!\left(\detach(f_{bthk}+b^d_{hk})\right),
  \label{eq:live-mean}
\end{equation}
and a learnable token/head residual
\begin{equation}
  z^{\rm clk}_{bth}=\frac{\langle w_h,x_{bt}\rangle}{\sqrt D}+\beta_h,
\end{equation}
evaluated in FP32.
The E2 clock is
\begin{equation}
  \boxed{
  \tau_{bth}=C\,\sg\!\left[
    \log m_{bth}-\log(C-a_h)+\tanh(z^{\rm clk}_{bth})
  \right]}
  \qquad(C=0.1).
  \label{eq:e2-clock}
\end{equation}
The bounded residual is strictly in $(-1,1)$ and the sigmoid keeps a finite valid clock strictly in
$(0,C)$.  The implementation does not clamp or use a straight-through estimator: invalid,
nonfinite, endpoint, or dead-derivative conditions are made visible to health checks, with invalid
clock outputs replaced by NaNs.

Detachment is a model choice, not merely an optimization.  The current radial timescale determines
the forward phase scale, while the phase loss cannot change $f$, $b^d$, $A^{\log}$, $\alpha$, or
$\omega$ through this clock.  The persistent anchor fixes the initialization reference across
save/resume; recomputing it from a trained $b^d$ would change the model after restart.

\subsection{Angle, sign convention, and rotating basis}

One projection shared across heads produces $e_{btj}$.  The head-specific increment is
\begin{equation}
  \delta_{bthj}=\pi\tanh(e_{btj})\tau_{bth},
  \qquad |\delta_{bthj}|<0.1\pi.
  \label{eq:rdiag-delta}
\end{equation}
RDiag uses the negative inclusive phase scan
\begin{equation}
  \phi_{bthj}=\phi_{b0,hj}-\sum_{s=1}^{t}\delta_{bshj},
  \label{eq:negative-phase}
\end{equation}
implemented in FP32 and reduced to $[0,2\pi)$ at the bounded-clock consumer.  The same
$R(+\phi_{thj})$ is applied to the corresponding pairs of $q_t$ and $k_t$.  Hence, for $s<t$,
\begin{equation}
  (R(\phi_t)q_t)^{\tr}R(\phi_s)k_s
  =q_t^{\tr}R\!\left(\sum_{u=s+1}^{t}\delta_u\right)k_s.
  \label{eq:rdiag-relative}
\end{equation}
Under the column-vector convention of Eq.~\eqref{eq:rotation}, the negative scan therefore
represents $R(+\delta_t)$ transport.  Mamba-3's positive scan represents $R(-\Delta\theta_t)$.
This sign is a gauge choice because replacing the zero-initialized angle projection by its negative
maps one convention to the other.  It is nevertheless an implementation invariant: forward,
streaming carry, and the negative suffix-scan VJP must use the same sign.

More explicitly, let ${\cal R}_t$ be block diagonal with the pair rotations $R(\phi_{thj})$ and an
identity tail.  The corresponding physical transition and process covariance are
\begin{equation}
  D_t={\cal R}_t^{\tr}\diag(\alpha_t){\cal R}_{t-1},
  \qquad
  \Omega_t^{\rm phys}={\cal R}_t^{\tr}\diag(\omega_t){\cal R}_t.
  \label{eq:rdiag-physical}
\end{equation}
On a tied pair, these reduce to $D_{tj}=\alpha_{tj}R(+\delta_{tj})$ and
$\Omega_{tj}^{\rm phys}=\omega_{tj}I_2$.  With moving-frame state
$\bar S_t={\cal R}_t S_t$ and addresses
$\bar k_t={\cal R}_tk_t$, $\bar q_t={\cal R}_tq_t$, the mean-state recurrence becomes purely
radial and can reuse the diagonal-KLA memory kernel.

\subsection{Diagonal Kalman gain and innovation memory}

Let the phase-rotated addresses be $\widetilde q_t=\bar q_t$ and
$\widetilde k_t=\bar k_t$.  With posterior diagonal
information $c_{t-1,hk}>0$, observation noise $r_{th}>0$, and information scale $s_h>0$, RDiag KLA
uses the moving-frame state (written as $S_t$ below for brevity)
\begin{align}
  p^-_{thk}
    &=\frac{\alpha_{thk}^{2}}{c_{t-1,hk}}+\omega_{thk},
    \label{eq:predict-cov}\\
  \kappa_{thk}
    &=\frac{p^-_{thk}\widetilde k_{thk}}
      {r_{th}+\sum_i p^-_{thi}\widetilde k_{thi}^{2}},
    \label{eq:gain}\\
  c_{thk}
    &=\frac{1}{p^-_{thk}}+
      s_h\frac{\widetilde k_{thk}^{2}}{r_{th}}.
  \label{eq:posterior-info}
\end{align}
For associative state $S_{th}\in\R^{K\times V}$,
\begin{align}
  S^-_{th}&=\diag(\alpha_{th})S_{t-1,h},\\
  S_{th}&=S^-_{th}+\ve{\kappa}_{th}
    \left(\ve{v}_{th}-\widetilde{\ve{k}}_{th}^{\tr}S^-_{th}\right)^{\tr},
    \label{eq:kla-state}\\
  o_{th}&=S_{th}^{\tr}\left(K^{-1/2}\widetilde{\ve{q}}_{th}\right).
  \label{eq:kla-read}
\end{align}
This is the KLA estimator.  It is not a Mamba SSM write: $\kappa_t$ is an adaptive anisotropic
Kalman gain and the write is the current innovation.  The diagonal information recurrence can be
computed by its existing associative M\"obius scan, followed by the existing chunked affine memory
recurrence.  In physical coordinates, $c_t$ represents the model's diagonal projection of
${\cal R}_tP_t^{-1}{\cal R}_t^{\tr}$; it is not a claim that rotating an arbitrary anisotropic
posterior covariance remains diagonal.  In particular, $c$, $\kappa$, $k$, and $q$ are never
pair-tied.

\subsection{VJP isolation}

The computational graph has two intentionally separate branches:
\begin{align*}
  (f,b^d,A^{\log},\omega^{\rm in})
    &\longrightarrow (\alpha,\omega)\longrightarrow\text{Kalman dynamics},\\
  \detach(f+b^d),\ a,\ (W,\beta),\ e
    &\longrightarrow \tau,\delta,\phi\longrightarrow\text{rotated }q,k.
\end{align*}
In its radial-only specialization, the direct-pair backward has no angle data dependency or load.
Conversely, the clock's $m$ and $a$ arguments are detached, so the phase VJP returns no gradient to
the radial gate, radial bias, or $A^{\log}$.
Unlike Mamba-3, RDiag does not add an angular-clock contribution to the dynamical radial-step
gradient.  Among the new phase-specific parameters, gradients reach the angle projection and clock
weight/bias; they also return to the shared hidden states through both projections.  The isolation
claim is specifically that no clock edge reaches the radial parameters.

\section{What to take from Mamba-3}
\label{sec:takeaways}

The useful transfer from Mamba-3 is structural: encode rotation as a relative query/key phase with
a small shared angle projection and a token/head clock.  The estimator-specific pieces should
remain KLA-specific.

\begin{center}
\small
\begin{tabular}{@{}p{0.16\linewidth}p{0.23\linewidth}p{0.25\linewidth}p{0.27\linewidth}@{}}
\toprule
Aspect & Mamba-2 SISO & Mamba-3 SISO & RDiag KLA E2 \\
\midrule
Transition geometry
  & Scalar radial decay; no rotation
  & Radial SSM plus relative pair rotation
  & Pair-tied radial KLA plus relative pair rotation \\
Clock
  & $DT_{th}>0$
  & $DT_{th}>0$ shared by decay and angle
  & $0<\tau_{th}<0.1$, based on detached live radial mean plus residual \\
Angle rows
  & None
  & Token/pair rows shared across heads
  & Token/pair rows shared across heads; head-specific $\tau$ \\
Phase gauge
  & None
  & Positive cumulative phase
  & Negative cumulative phase; sign documented by Eq.~\eqref{eq:rdiag-relative} \\
Write
  & Euler-like $DT\,B\,x$
  & Trapezoidal two-endpoint SSM write
  & Kalman innovation with anisotropic gain \\
Clock gradient
  & Only SSM dynamics
  & Angle $dDT$ is added to dynamical $dDT$
  & Phase/radial VJPs intentionally isolated \\
Covariance
  & No Kalman posterior covariance
  & No KLA diagonal information update
  & Diagonal information in moving frame; pair-isotropic prediction parameters \\
\bottomrule
\end{tabular}
\end{center}

Concretely, RDiag should take the following from Mamba-3:
\begin{enumerate}[leftmargin=1.8em]
  \item represent rotation by the same adjacent-pair transform on queries and keys, avoiding a
        dense recurrent rotation;
  \item use a bounded angular rate $\pi\tanh(e)$ and one token/head clock broadcast over pair modes;
  \item share the compact angle projection across heads, while permitting head-specific timing;
  \item use FP32 phase arithmetic and, unlike the installed training wrapper, retain the
        materialized phase and streaming carry in FP32 under a specified modular reduction
        contract;
  \item treat every dependency of a shared clock explicitly in backward and test cancellation-heavy
        reductions against the executable reference; and
  \item retain half- or full-prefix rotary layouts rather than forcing every key lane to rotate.
\end{enumerate}

It should intentionally \emph{not} copy the following:
\begin{enumerate}[leftmargin=1.8em]
  \item Mamba-2's scalar-only transition as a purported rotational mechanism;
  \item Mamba-3's trapezoidal SSM write, because KLA's innovation update defines a different
        estimator;
  \item Mamba-3's angle-to-$DT$ gradient coupling, because E2 deliberately uses a detached radial
        timescale;
  \item approximate softplus, trigonometric, or reassociated reduction sequences without proving
        the required finite-precision parity;
  \item independent radial coefficients inside a rotated pair while claiming a literal
        scalar-damped rotation; they instead define a more general rotated anisotropic
        contraction; or
  \item hard clock clamps or straight-through boundary gradients, which conceal saturation instead
        of making it observable.
\end{enumerate}

\section{Implementation contract}
\label{sec:implementation}

\subsection{Layout and forward pipeline}

The production-oriented route is restricted to contiguous CUDA tensors, BF16 projected
$q,k,v,f_{\rm raw},z_{\rm raw},y_{\rm raw}$, FP32 frontend parameters, equal key/value head counts,
and $K\in\{64,128\}$.  The clock preactivation $z^{\rm clk}$ is a separate FP32 quantity.
The rotary fraction is $1/2$ or $1$, so $K_r$ is positive and even.  Adjacent memory lanes
$(0,1),(2,3),\ldots,(K_r-2,K_r-1)$ form the rotary pairs; no split-half pairing is permitted.

The forward pipeline is:
\begin{enumerate}[leftmargin=1.8em]
  \item project $q,k,v$ and the diagonal-KLA radial/noise quantities;
  \item compute the standard frontend's frozen $\alpha^{\rm in}$ and $\omega^{\rm in}$;
  \item call the unchanged direct-pair operator with a detached angle, retain its pair-tied
        $\alpha$, $\omega$, and log-decay prefix, and discard its legacy phase output;
  \item independently materialize the precise detached FP32 $[B,T,H,K]$ clock steps and form the
        exact native-framework radial mean $m$ described below;
  \item keep $Wx$ as the existing FP32 GEMM, then fuse the pointwise clock chain and
        Eq.~\eqref{eq:rdiag-delta}, using the exact CUDA-eager sigmoid quotient;
  \item perform the bounded negative phase scan in chunks of 64, retaining an FP32 carry;
  \item normalize $q,k$ in FP32, round the memory-facing copies as specified, and rotate adjacent
        pairs by the same phase; and
  \item run the diagonal information/gain kernel and the S3 chunked KLA memory recurrence.
\end{enumerate}
The gain-facing key remains FP32.  The memory-facing query/key route is rounded through BF16 in the
same place as its executable reference before rotation/final rounding; changing that boundary is a
numerical model change for parity purposes.

\subsection{The precise detached-clock mean}
\label{sec:precise-mean}

The actual radial dynamics retain the FLA softplus used by the qualified baseline.  The detached
clock uses a separate precise FP32 copy
\begin{equation}
  \splus_{20}(x)=
  \begin{cases}
    x,&x>20,\\
    \log1p(\exp x),&x\leq20,
  \end{cases}
  \label{eq:precise-softplus}
\end{equation}
evaluated with the precise device intrinsics.  To reproduce the eager reference, the optimized
path must materialize
\begin{equation}
  \mathrm{clock\_dt}\in\R^{B\times T\times H\times K}
\end{equation}
in FP32 and then invoke the framework reduction
\begin{equation}
  \mathrm{radial\_clock}
  =\texttt{torch.mean(clock\_dt, dim=-1, out=radial\_clock)}.
  \label{eq:torch-mean}
\end{equation}
A mathematically equivalent Triton lane sum is not an exact substitute.  It changes addition order
and was shown to leave small residual differences that can be amplified by long phase scans and
cancellation-heavy angle gradients.  The precise clock copy is detached; the FLA-softplus radial
outputs and their VJP remain untouched.
The temporary costs exactly $4BTHK$ bytes.  At $B=4,T=4096,H=16$, that is 64 MiB for $K=64$ and
128 MiB for $K=128$; at $H=18,K=128$ it is 144 MiB.  This cost belongs in the end-to-end memory and
speed gate rather than being hidden by a numerically different reduction.

\subsection{Reusing the direct-pair radial path safely}

The production wrapper does not maintain a second implementation of Eq.~\eqref{eq:pair-radial}.
It calls the established direct-pair transform with \texttt{angle\_raw.detach()}, ignores the
returned legacy delta, and reuses only $g$, $\alpha$, and $\omega$.  This keeps the radial forward
arithmetic and its VJP identical to the already-qualified direct-pair path.

Discarding the legacy delta makes its incoming cotangent absent.  The direct-pair backward turns
that fact into the compile-time specialization \texttt{HAS\_DDELTA=false}.  In this specialization
the generated kernel does not load the detached angle, evaluate its tanh, or form the phase term
$\pi\tanh(e)\,\overline\delta$.  The guard is semantically important: merely setting
$\overline\delta=0$ would not isolate a nonfinite discarded angle, because IEEE arithmetic gives
$0\times\mathrm{NaN}=\mathrm{NaN}$.  The phase-enabled
\texttt{HAS\_DDELTA=true} algebra is unchanged, while dedicated NaN/$\pm\infty$ tests establish
that a discarded phase cannot contaminate radial outputs or gradients.

\subsection{Exact sigmoid and bounded-delta forward}

Let
\begin{equation}
  \ell_{bth}=\log m_{bth}-\log(C-a_h)+\tanh(z^{\rm clk}_{bth}).
\end{equation}
For FP32 CUDA inputs, PyTorch's native sigmoid kernel evaluates
$1/(1+\exp(-\ell))$.  The generic Triton \texttt{tl.sigmoid} uses a different approximate
lowering and produced one-ULP clock differences even when $\ell$ was bitwise identical.  The
production kernel therefore spells out the eager expression as
\begin{equation}
  \tau^{\rm raw}_{bth}
  =0.1\,\texttt{tl.div\_rn}
    \left(1,\,1+\texttt{libdevice.exp}(-\ell_{bth})\right),
  \label{eq:exact-sigmoid}
\end{equation}
then applies the fail-visible validity predicate to obtain $\tau$.  The angle nonlinearity is
computed once by native ATen as
\texttt{torch.tanh(angle\_raw.float()).contiguous()} and passed to the fused clock kernel, which
forms $\delta=\pi\tanh(e)\tau$.  This boundary makes the actual-logit $\tau$ and $\delta$ bitwise
equal to the eager oracle, rather than merely close in real arithmetic.

\subsection{Phase scan and backward order}

The phase forward is an inclusive FP32 scan.  Each chunk computes
$\phi=\mathrm{carry}-\operatorname{cumsum}(\delta)$, wraps valid values to $[0,2\pi)$, and passes the
wrapped final value as the next carry.  Before reduction, $|\phi|$ must not exceed $65{,}536$;
nonfinite, out-of-contract, or noncanonical residues become sticky NaNs.  Its VJP is the global
negative suffix scan
\begin{equation}
  \overline\delta_t=-\sum_{u=t}^{T}\overline\phi_u,
  \label{eq:phase-vjp}
\end{equation}
with the final-carry cotangent subtracted from every token when present.

For the fused clock/delta backward, operation order is part of the contract.  Let
$g_{thj}=\overline\delta_{thj}$.  The eager-compatible order is
\begin{align}
  \overline\tau_{th}
    &=\sum_j\left[g_{thj}\bigl(\pi\tanh(e_{tj})\bigr)\right],
    \label{eq:dtau-order}\\
  \overline z^{\rm clk}_{th}
    &=\operatorname{mask}_{th}(\overline\tau_{th})
      \,\tau^{\rm raw}_{th}\left(1-\frac{\tau^{\rm raw}_{th}}{C}\right)
      \left(1-\tanh^2 z^{\rm clk}_{th}\right),\\
  \overline e_{tj}
    &=\left[\left(\sum_h g_{thj}\tau_{th}\right)\pi\right]
      \left(1-\tanh^2 e_{tj}\right).
  \label{eq:dangle-order}
\end{align}
Here $\operatorname{mask}_{th}(g)$ is $g$ when the forward validity predicate selected
$\tau^{\rm raw}_{th}$ and zero when it selected the detached NaN branch.  The subsequent
multiplication order is retained, so a NaN local derivative remains fail-visible even after a zero
cotangent mask.
Equation~\eqref{eq:dtau-order} must be evaluated as the displayed broadcast multiply followed by
the framework's sum over $P$; pulling $\pi$ outside the sum is not accepted as ``equivalent.''
For Eq.~\eqref{eq:dangle-order}, multiply by $\tau$, reduce over $H$, multiply by $\pi$, and only
then apply the tanh derivative.  Concretely, production saves the native-ATen tanh output from the
forward and evaluates
\begin{center}
\texttt{torch.ops.aten.tanh\_backward(}\linebreak[0]
\texttt{  (d\_delta * tau[..., None]).sum(dim=2) * pi, angle)}.
\end{center}
The result is cast to \texttt{angle\_raw}'s dtype (BF16 on the production route) before the angle
projection VJP.  No Triton angle-partial tensor or head-reduction kernel remains.  The validity
mask is applied to the incoming clock cotangent before the sigmoid/tanh clock chain, matching the
eager \texttt{where(valid, raw\_tau, detached\_nan)} VJP.  Weight and hidden-state gradients then
flow through the ordinary linear operations.

\subsection{Persistent state and FSDP}

The anchor $a_h$ must satisfy all of the following:
\begin{enumerate}[leftmargin=1.8em]
  \item FP32, contiguous, shape $[H]$, finite, and strictly inside $(0,0.1)$;
  \item registered as a persistent buffer and present in the model state dictionary;
  \item absent from named parameters and optimizer groups;
  \item identical across distributed ranks after construction, FSDP wrapping, save, and resume;
  \item restored from the checkpoint rather than regenerated from the current radial bias.
\end{enumerate}
Production enforces this with \texttt{bounded\_rdiag\_fsdp\_mixed\_precision}: it retains the
installed BF16-mixed parameter/reduction policy but overrides
\texttt{buffer\_dtype=torch.float32}.  Resume tests compare the anchor bytes before and after round
trips and verify that fresh construction cannot silently overwrite a checkpointed anchor.

\subsection{Initialization and health signals}

The angle projection weight and bias, and the bounded-clock residual weight and bias, are initialized
to zero.  Rotation therefore starts at the \emph{pair-tied} diagonal-KLA function: $e=0$ gives
$\delta=0$ even when $\tau$ is nonzero, while the pair averaging in Eq.~\eqref{eq:pair-radial}
remains active.  This is preferable to initializing the clock at zero, which would place it on a
boundary and kill useful clock gradients once angles become live.

The audited bounded RDiag named configurations use $r_{\min}=0.01$, process-noise floor
$\omega_{\min}=q_{\min}=0$, and \texttt{gain\_init=iso}.  The last setting zeroes the output map of
the process-noise projection, so $\omega$ starts uniform across key channels; it does \emph{not}
initialize $\omega$ itself to zero, because the positive softplus remains.  These are recipe and
initialization choices rather than changes to Eqs.~\eqref{eq:predict-cov}--\eqref{eq:kla-read}.

At every checked optimizer boundary, diagnostics should aggregate across microbatches, layers, and
ranks:
\begin{itemize}[leftmargin=1.8em]
  \item nonfinite counts for the radial mean, base logit, residual, logit, raw clock, valid clock,
        derivative, phase, parameters, and gradients;
  \item exact endpoint and near-boundary counts, using $10^{-4}$ and $0.1-10^{-4}$ as the observed
        lower and upper near-boundary bands;
  \item minimum positive clock Jacobian and a near-dead count below $10^{-8}$;
  \item token variation of the residual projection after bootstrap;
  \item phase magnitude, delta magnitude, and canonical wrapped-range checks; and
  \item schema/count agreement across ranks, including exactly one observation per expected
        microbatch/layer pair.
\end{itemize}
The kernel itself fails visibly at invalid endpoints.  Promotion screens may impose stricter
distributional thresholds than the continuous health controller; both classes of gate are needed.

\section{Reference pseudocode}

The following pseudocode fixes the mathematical order without prescribing launch geometry.
\begin{enumerate}[leftmargin=1.8em]
  \item \textbf{Radial and clock preparation.}
  \begin{align*}
    x^{d}&=f+\texttt{dt\_bias},\\
    d^{\rm radial}&=\texttt{FLA\_softplus}(x^d),\\
    \texttt{clock\_dt}&=\splus_{20}(\detach(x^d))\quad\text{materialized FP32 }[B,T,H,K],\\
    m&=\texttt{torch.mean(clock\_dt, dim=-1, out=radial\_clock)}.
  \end{align*}
  Invoke the unchanged direct-pair transform with a detached angle, pair-average
  $d^{\rm radial}$ and $\omega$ on the rotary prefix, and construct $\alpha$ from
  Eq.~\eqref{eq:pair-radial}.  Discard its legacy delta and do not use
  \texttt{clock\_dt} for radial dynamics.

  \item \textbf{Bounded clock and phase.}
  \begin{align*}
    z^{\rm clk}&=\texttt{linear}(x,W)/\sqrt D+\beta,\\
    \ell&=\log m-\log(0.1-a)+\tanh z^{\rm clk},\\
    \tau&=0.1\,\texttt{div\_rn}(1,1+\texttt{libdevice.exp}(-\ell)),\\
    \widehat e&=\texttt{torch.tanh}(e.\texttt{float()}),\\
    \delta&=\pi\widehat e\,\tau,\\
    \phi_t&=\operatorname{wrap}_{2\pi}(\phi_{t-1}-\delta_t).
  \end{align*}
  Reject invalid open-clock or phase-contract values by producing observable NaNs.

  \item \textbf{Addresses and Kalman recurrence.}
  Rotate the same adjacent prefix of normalized $q$ and $k$ by $R(\phi)$, preserve tail lanes,
  then apply Eqs.~\eqref{eq:predict-cov}--\eqref{eq:kla-read}.  The chunked implementation must
  match this token recurrence under the documented BF16/FP32 boundary.

  \item \textbf{Backward.}
  Run Eq.~\eqref{eq:phase-vjp}; use Eqs.~\eqref{eq:dtau-order}--\eqref{eq:dangle-order} exactly;
  evaluate the angle derivative with native \texttt{aten.tanh\_backward} on the saved
  $\widehat e$; compile the discarded direct-pair phase branch with
  \texttt{HAS\_DDELTA=false}; return no clock gradient to $m$ or $a$; and run the independent
  radial/Kalman VJPs.
\end{enumerate}

\section{Lessons from the E-screen and performance/numerical debugging}
\label{sec:lessons}

\paragraph{An open interval is not a distributional guarantee.}
An earlier bounded-clock E screen completed 519 optimizer updates and produced no exact endpoints or
nonfinite clocks, yet the step-520 evaluator rejected it.  In layer 0, the lower-near-boundary
fraction was $0.00205326$, above the preregistered $10^{-3}$ threshold; the observed clock range was
approximately $[8.65\times10^{-5},\,0.09813]$.  Thus a smooth sigmoid bound alone did not prevent
probability mass from accumulating near a boundary.  The design response was the detached live
radial base plus a bounded token/head residual, not a harder clamp.

The redesigned E2 screen reached its step-520 evaluation after 519 completed optimizer updates
(520 measured gradients), with zero lower and upper near-boundary counts in the reported gate and
an observed clock range of approximately
$[0.01060,\,0.08181]$.  That result supports the clock design and justified integration work.  It
does not prove long-run training quality, kernel parity under every shape, distributed resume
correctness, or a production speedup.

\paragraph{Correctness-first composition exposed the actual performance target.}
At $B=4,T=4096,H=16,K=V=128,D=2048$, the unfused rotational-v1 wrapper measured 24.706 ms for
combined mixer forward/backward versus 13.735 ms for its frozen control, a $1.799\times$ ratio
(79.9\% overhead).  The identical-input frozen S3 stage was $0.999\times$ its control, localizing
the cost to the eager second-$dt$, pair tying, phase/trigonometry, three rotations, and prefix
rebuild.  This motivated the direct radial kernel, fused normalization/rotation, and fused
clock/delta path while retaining the established gain and S3 kernels.  These are historical
isolation numbers for the v1 geometry, not a speed claim for the final E2 composition.  The later
final A/B measurements, including the exact-mean FP32 scratch, are reported in
Table~\ref{tab:final-perf}.

\paragraph{A tiny local error can become a large gradient discrepancy.}
The first optimized clock path reused the FLA radial softplus.  Against eager PyTorch, the initial
radial-mean discrepancy was only about $6.33\times10^{-8}$ in maximum absolute value and
$1.54\times10^{-6}$ in relative $L_2$.  After the logit, long phase scan, BF16 rotation boundary,
and cancellation across heads, the angle-weight gradient showed about $3.81\times10^{-3}$ relative
$L_2$ error in the focused diagnostic.  Near-zero reference entries also made elementwise relative
metrics misleading, so absolute, normwise, and cancellation-stress gates are all required.

Using Eq.~\eqref{eq:precise-softplus} reduced the clock error substantially, but a Triton reduction
still did not reproduce eager bitwise.  Materializing the FP32 $[B,T,H,K]$ precise values and using
the exact framework mean was bitwise identical in the focused $K=64$ and $K=128$ experiment while
leaving the FLA radial $\alpha$, log-prefix, and $\omega$ outputs bitwise unchanged.  Likewise,
algebraically reassociated all-Triton backward reductions failed the focused VJP gate; the explicit
framework reduction of Eq.~\eqref{eq:dtau-order} and the factor order of
Eq.~\eqref{eq:dangle-order} passed it.  The lesson is narrow but important: the executable eager
graph, including reduction order, is the reference for this mixed-precision path.

\paragraph{Elementary functions are also part of the executable reference.}
After the exact mean fix, a full $K=64$ trace found bitwise-identical radial clock, base logit,
bounded residual, and final logit, yet the fused clock still differed from eager at 61,655 of
262,144 FP32 $\tau$ entries (maximum absolute error $7.45\times10^{-9}$).  Those one-ULP clock
differences propagated to 929,205 of 4,194,304 delta entries and were amplified by the length-4096
phase scan.  The cause was \texttt{tl.sigmoid}, not the model equations.  Replacing it with
Eq.~\eqref{eq:exact-sigmoid}, matching PyTorch's CUDA source expression and rounded division,
removed the actual-logit discrepancy.  Computing the angle tanh and its backward with native ATen
then removed the remaining reduction/transcendental-order drift.  This is why the final path uses
selective eager operations at two small boundaries rather than insisting that every operation be
fused into Triton.

\section{Invariants and qualification status}

The following are invariants, not tunable approximations:
\begin{enumerate}[leftmargin=1.8em]
  \item rotary prefix width is positive and even; pairing is adjacent;
  \item $\alpha$ and $\omega$ are exactly equal within each selected pair;
  \item $q$ and $k$ receive the same phase and same rotation matrix;
  \item the phase scan is negative and inclusive, and its backward is a negative suffix scan;
  \item $m$ and $a$ are detached FP32 quantities; phase gradients do not enter radial dynamics;
  \item the valid clock and its local derivative are finite and strictly positive, with
        $0<\tau<0.1$;
  \item the anchor is persistent FP32 checkpoint state, not optimizer state;
  \item precise clock softplus plus \texttt{torch.mean} is separate from FLA radial softplus;
  \item the bounded sigmoid uses Eq.~\eqref{eq:exact-sigmoid}, not \texttt{tl.sigmoid};
  \item the displayed forward and backward reduction orders are preserved, with the angle tanh
        and its VJP evaluated by native ATen;
  \item a discarded direct-pair delta selects \texttt{HAS\_DDELTA=false}, eliminating every angle
        access and phase product from the radial backward; and
  \item the unrotated tail is an exact pass-through of diagonal KLA.
\end{enumerate}

The one-H200 kernel/layer checklist is closed by the cumulative qualification campaign:
\begin{enumerate}[leftmargin=1.8em]
  \item pure-PyTorch token recurrence versus optimized forward for zero, half, and full rotation;
  \item first-order VJPs for every input/parameter, including adversarial cancellation and partial
        chunks, at $K=64$ and $K=128$;
  \item exact checks for radial outputs and detached clock mean where exactness is part of the
        contract, plus justified tolerances elsewhere;
  \item segmented/streaming equivalence, phase carry, wrap-boundary, and sticky-NaN tests;
  \item zero-angle equivalence to pair-tied diagonal KLA and angle-sign/gauge unit tests;
  \item multi-step health checks covering bootstrap-zero angle, subsequently live clock gradients,
        token variation, and no nonfinite/endpoints;
  \item exact full-layer forward/VJP comparison with the sealed E2 candidate at $K=64$, $K=128$,
        and the 18-head 1.3B-style geometry; and
  \item interleaved A/B timing against both the sealed E2 candidate and the legacy pair clock.
\end{enumerate}

The authoritative one-H200 job 1777777 completed with exit code 0 and reported 117 passing core
RDiag tests followed by 54 passing E2 tests.  Table~\ref{tab:final-perf} gives the paired runtime
ratios for an exact RDiag-layer forward plus backward; values below one are faster than the named
reference.  Every upper endpoint of the bootstrap 95\% confidence interval is below the 1.05
promotion limit.

\begin{table}[ht]
  \centering
  \small
  \caption{Final H200 paired A/B timing from job 1777777, with 40 measured rounds after 20 warmups.
  ``Accepted'' is the sealed eager E2 candidate and ``legacy'' is the production pair-$dt$ path.}
  \label{tab:final-perf}
  \begin{tabular}{@{}lrrrr@{}}
    \toprule
    Geometry $(H,D,K)$ & E2/accepted & 95\% CI high & E2/legacy & 95\% CI high \\
    \midrule
    $(16,1024,64)$   & 0.969512 & 0.978247 & 1.014135 & 1.020172 \\
    $(16,2048,128)$  & 0.950407 & 0.956546 & 1.008590 & 1.013666 \\
    $(18,2304,128)$  & 0.960111 & 0.964929 & 1.013622 & 1.014984 \\
    \bottomrule
  \end{tabular}
\end{table}

Thus the optimized E2 implementation is bitwise faithful where exact parity is required, faster
than the sealed accepted implementation in all three measured geometries, and within roughly
1--1.4\% of the simpler legacy clock.  This closes the E2 one-H200 kernel/layer gate, but not a full
production-submission gate.  Train-smoke, distributed checkpoint/resume, and multi-node
soak/attestation remained for that baseline at the time of sealing.  The next section records the
later shared-DT candidate and its separate, still-open qualification.

\section{Shared-DT engineering status as of September 6, 2026}
\label{sec:shared-dt-status}

\noindent\textit{Status updated 2026-09-06 18:53 UTC.}

This section records a newer production candidate and does not revise the closed E2 result above.
The candidate starts from commit
\texttt{6472a0711eb61177c3ef91845855fffa82202290} and replaces E2's detached bounded clock with a
shared-DT transition.  Its qualification is \emph{not complete}.  Four gates passed in the first
campaign, whose gradient gate exposed both a defective acceptance rule and real rotational-gradient
amplification; that soak was cancelled before allocation.  The subsequent reviewed head-scaled
commit passed all five non-soak gates, but its 3300-update soak failed closed with an aggregate
count at boundary 269.  A divergent diagnostic replay independently exposed rejection of the
finite-input exact-zero decay endpoint.  The narrow lower-endpoint correction then passed all five
non-soak gates and 675 soak updates before failing closed at boundary 676 with an unrecorded key.
A second divergent replay independently exposed rejection of the finite-input exact-one endpoint;
the reviewed upper-endpoint correction passed its fresh numerical H200 suite, but its evidence
could not be sealed because a pre-existing asynchronous log-finalization race omitted the terminal
PASS marker from the retained log.  The narrow source-gate logging correction was reviewed,
committed, and published.  Its fresh campaign passed and sealed the H200 gate and all four short
distributed gates, then its soak crossed both old boundaries and failed closed at update 790 on a
rank-equal pre-clip global norm that triggered the combined nonfinite, negative, or greater-than-10
rejection predicate.  The failed run did not retain which branch fired.  No soak seal or aggregate
authorization exists; the diagnostic replay and threshold-provenance analysis are recorded below.

\subsection{Relation to Mamba-2 and Mamba-3}

The source-level comparison in Sections~\ref{sec:mamba2} and~\ref{sec:mamba3} found no discrepant
transition equation or kernel sign bug in the shared-DT path.  Its radial factor is the Mamba-2
form
\begin{equation}
  \alpha_{th}=\exp\!\left[-\exp(A_{\log,h})DT_{th}\right],
\end{equation}
and its bias-free pair-angle projection follows the installed Mamba-3 construction with an
additional versioned head scale.  For the b37 v2 candidate,
\begin{equation}
  e_{tj}=(W_{\rm angle}x_t)_j,\qquad
  e^{\rm eff}_{tj}=0.5H^{-1/2}e_{tj},\qquad
  \delta_{thj}=\pi\tanh(e^{\rm eff}_{tj})DT_{th},
  \qquad
  \phi_{thj}=-\sum_{s\leq t}\delta_{shj},
  \label{eq:shared-dt-phase}
\end{equation}
with one angle value per token/pair broadcast over heads, a headwise $DT$, a cumulative phase scan,
and the corresponding negative reverse-suffix backward.  This is the literal RDiag source
convention; installed Mamba-3 uses the positive scan.  The two parameterized model families are
related by the gauge transform $W_{\rm angle}\mapsto-W_{\rm angle}$; exact continuation of an Adam
checkpoint also negates that parameter's first moment.  This rules out a structural sign error, but
flipping only the scan sign at one fixed, untransformed checkpoint can change outputs and gradient
norms.
The relevant difference from the E2 model is instead that the angle and $DT$ paths remain live and
the same shared $DT$ controls radial decay and rotation.
Writing $s_H=0.5H^{-1/2}$ for v2 and bars for cotangents, the exact phase branch contributes
\begin{align}
  \overline{\delta}_{thj}
    &=-\sum_{u\geq t}\overline{\phi}_{uhj},\\
  \overline e_{tj}
    &=s_H\pi\operatorname{sech}^2(s_He_{tj})
      \sum_h DT_{th}\overline{\delta}_{thj},\\
  \overline{DT}^{\,\mathrm{phase}}_{th}
    &=\pi\sum_j\tanh(s_He_{tj})\overline{\delta}_{thj}.
  \label{eq:shared-dt-vjp}
\end{align}
The first expression exposes the reverse suffix accumulation, the second the head reduction, and
the third the phase contribution added to the radial $DT$ gradient.  Existing aggregate category
norms do not separately observe those three effects.

\subsection{Issue, diagnosis, and completed work}

\paragraph{Two distinct problems were found.}
First, the original per-microbatch global-gradient limit of 10 was not a valid gate: all 24 measured
model/batch observations exceeded 10, including stock DiagKLA at 14.001--14.232.  Stock DiagKLA and
the bounded-clock layer were also architecture-mismatched controls.  Second, matched interventions
localize genuine amplification to the live shared-angle cumulative-rotation branch.  Its output is
broadcast over recurrent heads and its VJP contains both a head reduction and a reverse temporal
suffix sum, but the retained evidence does not yet separately measure head coherence, suffix
growth, or the phase contribution to the shared-DT gradient.  This is evidence for a
model-parameterization or optimization issue, not for a faulty radial or phase kernel, while the
exact causal decomposition remains open.

On the same eight real FineWeb batches, the unmodified shared-DT model produced global norms
34.596--45.300 (median 38.758), versus a DiagKLA median of 14.096.  Disabling rotation returned the
median to 14.180, while detaching only the angle path gave 18.035 and zero angle initialization
still gave 25.593.  A phase-only $DT$ cap of 0.1 was insufficient.  Scaling the effective angle
preactivation by $1/\sqrt{16}=0.25$ gave median 16.765 and shared/Diag median and 95th-percentile
nearest-rank ratios 1.191 and 1.261, with live angle and $DT$ gradients.  These interventions
localize amplification to head-shared cumulative rotation, but do not alone identify which of its
head sum, suffix accumulation, angle VJP, or phase-to-$DT$ VJP is sufficient for the late
instability.

\paragraph{Completed engineering work.}
The Mamba-2/Mamba-3 source audit, exact H200 measurement, rotation/gradient ablations, matched-
backbone control, and scale/cap sweep are complete.  The replacement design and implementation
plan received an independent review with no remaining Important findings.  The layer-side
$H^{-1/2}$ change and its nonpersistent-state host contracts are now implemented in the isolated
checkout below: 41 host tests passed.  Development H200 job 1827839 passed the focused H=18
all-parameter oracle/VJP case and H=16, T=4096 BF16 boundary/backward case; job 1827905 then passed
the full 109-test production-H200 suite.  These runs exercised the mutable development checkout,
so they are correctness evidence rather than immutable qualification seals.

The same checkout now also contains the complete gradient-v3 control closure: both arms construct
the same 220M configuration, strictly load one full reference state, prove matching configuration,
parameter-inventory, and full-state digests, and reset CPU/CUDA randomness symmetrically.  The
no-rotation arm installs connected-zero hooks only on the twelve angle projections, checks all
twelve angle-weight gradients are present, finite, and exactly zero for every batch, checks the
shared radial gradient norms remain finite and positive, and proves hook removal.  The v3 recipe,
catalog digest, result attestor, stale-schema rejection, and aggregate-v3 recursive requirement are
integrated.  The focused gate-runner host suite passed 38 tests; the complete stable host sweep is
also complete: the qualification-control suite passed 71 tests with one filesystem-capability skip
(the skipped case passed separately on local temporary storage), and the broader thirteen-file
shared-DT host sweep passed 418 tests with three expected skips.

In campaign
\texttt{9870d362bcbb4ec3b9b3cafb13459ef8}, production-H200 (job 1824028), train-smoke (1824365),
checkpoint-fault (1824367), and authenticated-resume (1824368) produced recursively valid seals
with SHA-256 prefixes \texttt{ecc03259}, \texttt{9ca6354e}, \texttt{06c1f60a}, and
\texttt{656c9412}, respectively.  Gradient job 1824370 failed the invalid absolute-10 rule and
produced no seal; the soak was cancelled before allocation.  These seals characterize the
unmodified base commit only and do not qualify the proposed head-scaled change.

\paragraph{Qualification outcome.}
The implementation commit passed separate exact-SHA specification/math and code-quality reviews
without a finding and was pushed to the branch named below.  Its immutable snapshot contains 4227
files and has source-manifest SHA-256
\pathcode{5b6b80bf8c38f4c959ac1f908a10fc01a25ff6986e77f0ee1c0a35de6843fbd9}.
Fresh campaign \pathcode{a65918f1771446e7acfbffcc3e5c7dc3} is closed incomplete.  Production-H200 gate job
1828250 passed all 109 tests and attestor job 1828252 published seal SHA-256
prefix \texttt{a05c99d0544a721a}.
The other four short gates and their attestors passed.  Gradient A/B job 1828811 measured
shared/no-rotation median and 95th-percentile ratios 1.179 and 1.262, inside the 1.25/1.50 limits;
attestor 1828820 published seal SHA-256
prefix \texttt{0eabab83a02ac397}.
Checkpoint-fault jobs 1828813/1828823, authenticated-resume jobs 1828814/1828821, and train-smoke
jobs 1828815/1828822 published seals with SHA-256 values
prefixes \texttt{a3558a17796d427c}, \texttt{1ea7a7da40576d20}, and
\texttt{c2e19737cff63305}, respectively.  They are valid partial evidence but cannot authorize the
commit without a valid soak seal.

The production soak did \emph{not} pass.  Job 1828812 trained normally through completed update
268 (loss 5.356, global gradient norm 1.101, and zero reported health failures), then failed closed
at the pre-optimizer boundary for update 269 with aggregate health-failure count 158720.  It
therefore emitted neither a result nor a seal, and dependent attestor 1828824 was cancelled.  The
exception predated the detailed failed-counter breakdown.  Its stage excluded
observation/parameter/gradient coverage, gradient-nonfinite, and bootstrap-liveness failures, but
could not distinguish the remaining forward-domain predicates from parameter non-finiteness or
collective-structure disagreement.  That ambiguity motivated the fail-closed diagnostic replays
below.

The first diagnostic replay identified the candidate predicate.  Non-authorizing job 1829026 ran
on a different H200 node pair and observed its first exact FP32 zero earlier, at boundary 241, so
the exact-boundary harness correctly exited nonzero rather than labeling the run a reproduction.
Every one of the 16 ranks nevertheless reported the same sole failed key,
\texttt{alpha\_nonpositive\_count=128}: one scalar $[B,T,H]$ decay site materialized over the
$K=128$ key width.  All other failure counters were zero; the aggregate diagnostics included
$\min\alpha=0$, $DT_{\max}=7.3498993$, and $\max|\tanh(e^{\rm eff})|=0.9377124$.
The retained scheduler log is
\pathcode{/checkpoints/ngocbh/hdn/ckpts/logs/rdiag_sdt_soak269_diag_1829026.log}, with SHA-256
{\ttfamily\scriptsize 5dc12cd3ea0ea52765495763b28bd0185b038c306bad6c333fabb537eb67e6ee}.
This is evidence for the same candidate failure mode, not an exact reproduction of the original
boundary-269 trajectory.  The original aggregate count $158720=1240\mathbin{\times}128$ is
consistent with the same width replication, but its exception did not retain a per-key breakdown.
Intermediate job 1829217, from commit
\pathcode{28531cf959304ab26b22789f2ea4ce082adbed13}, remained healthy through boundary 269 on yet
another node pair and correctly failed closed without a summary when its search condition was not
met.  Its source-manifest and scheduler-log SHA-256 values are respectively
{\ttfamily\scriptsize 0dee86b8f758f14aabd78047e2d0d530475be2d29f99d27587c256233644c73f} and
{\ttfamily\scriptsize 142120433da0eddb643d6d6160974c756397ed730fe802c009fbd1eaed1b3924}.

Final non-authorizing search job 1829305 ran from diagnostic commit
\begin{quote}
\pathcode{06eaf237f5be9f79ee8d00075b79dcc0e7a9e96a}.
\end{quote}
It completed 277 optimizer updates and stopped before clipping or update 278 with status
\texttt{SAME\_FAILURE\_MODE} and
\texttt{authorizing=false}.  All 16 ranks agreed on the sole production-predicate count
\texttt{alpha\_nonpositive\_count=128}; removing the repeated key-width dimension leaves one
unique scalar decay site.  Every recorded $A_{\log}$, $\exp(A_{\log})$, and
$\exp(A_{\log})DT$ value was finite.  Their global ranges were respectively
$[0.0472397,2.7735384]$, $[1.0483732,16.0152016]$, and
$[0.000105659,104.639008]$.  The largest exponent exceeds the FP32 round-to-nearest zero threshold
of approximately $103.972$; $87.3365$ is only the normal-to-subnormal onset.  Thus finite inputs to
$\exp[-\exp(A_{\log})DT]$ can produce the observed stored zero.

The retained summary is
\begin{quote}
\pathcode{diagnostic_runs/soak269_replay_06eaf237f5be_1829305/}\\
\pathcode{artifacts/replay-summary.json},
\end{quote}
with SHA-256
{\ttfamily\scriptsize 558a1dde9c842b95d6afaafb16c54994f8e75aae68b7fed0bfef11cfa4718fb5};
its trajectory has SHA-256
{\ttfamily\scriptsize eeb607f58f40a156e4fb42df899ef27f724a17c3ede5f85906395a4917692614}.
The scheduler log has SHA-256
{\ttfamily\scriptsize c46c802ac2b73fa875a40699ce2ac96b2f32bff61c1e128a54625e8916264b03}.
The replay compared 268 original updates, retained nine additional updates, and did not exactly
match the original prefix; it therefore establishes the floating-point endpoint and candidate
failure mode, not an exact boundary replay.  The harness temporarily exempted only
\texttt{alpha\_nonpositive\_count} so that it could retain this evidence and still stopped before
the update; its logged zero health-failure count is not a production pass.

The evidence classifies the reproduced zero-alpha endpoint as a health-policy false positive rather
than a model or kernel math defect; because the original exception lacks a per-key breakdown, it
does not prove that job 1828812 had identical causality.  At original update 268, 583512 distinct
scalar decay sites were already below the diagnostic $10^{-6}$ floor while loss and gradients
remained finite.  The S3 prefix clamps this range to $10^{-6}$ and assigns the below-floor branch
zero cotangent; the Kappa recurrence is also defined at raw $\alpha=0$.  The narrow correction keeps
\texttt{alpha\_nonpositive\_count} as telemetry, adds \texttt{alpha\_negative\_count} as the fatal
domain predicate, and leaves the model's alpha and both recurrences unchanged.  At commit
\pathcode{83118d2...}, alpha nonfinites and $\alpha\geq1$ remained fatal.

One login-namespace preparation attempt failed closed on the intentionally strict live-data
identity check, and the first compute-side preparation attempt (job 1828245) failed before manifest
publication because its three mutable roots had not been pre-created.  Preparation was rerun in the
Slurm user namespace with those private roots fixed at mode 0700; neither failed attempt executed a
numerical gate.

\paragraph{Diagnostic evidence.}
The primary metric artifacts are
\begin{itemize}[leftmargin=1.8em]
  \item \pathcode{diagnostic_runs/gradient_norm_6472a071_20260902T170943Z/metrics.json},
        SHA-256 \texttt{9c6a8a96a297799a};
  \item \pathcode{diagnostic_runs/gradient_rotation_ablation_6472a071_20260902T183817Z/metrics.json},
        SHA-256 \texttt{274abd1902e7c7d3};
  \item \pathcode{diagnostic_runs/gradient_matched_control_6472a071_20260902T185000Z/metrics.json},
        SHA-256 \texttt{d908637092fdfc3e}; and
  \item \pathcode{diagnostic_runs/gradient_angle_scale_6472a071_20260902T185100Z/metrics.json},
        SHA-256 \texttt{e503ab386a9de79d}.
\end{itemize}
All are rooted at
\pathcode{/lustre-storage/checkpoints/ngocbh/hdn/diagnostics/rdiag_shared_dt_v1}.

\subsection{Approved implementation and remaining qualification}

\paragraph{Chosen change.}
For $H$ recurrent heads, the development layer derives the nonpersistent scalar $H^{-1/2}$ and forms
\begin{equation}
  e^{\rm eff}_{tj}=H^{-1/2}(W_{\rm angle}x_t)_j,
  \qquad
  \delta_{thj}=\pi\tanh(e^{\rm eff}_{tj})DT_{th}.
  \label{eq:shared-dt-headscale}
\end{equation}
The projection remains bias-free and its stored weight is unchanged.  In the BF16 production path,
the layer performs the multiply on the BF16 projection output before the frontend's FP32 cast and
FP32 \texttt{tanh}; the frontend must not scale a second time.  The scale is neither a parameter nor
a buffer, so old state dictionaries remain strictly loadable.  This preserves the transition
family and both learning paths while normalizing the initial head-summed VJP.

The replacement gradient gate is an exact same-model/same-state two-arm intervention.  The
\texttt{shared\_dt} arm runs normally; \texttt{shared\_dt\_no\_rotation} replaces only each angle-
projection output by the connected zero \texttt{output * 0.0}.  Config, full state, parameter
inventory, and radial shared-DT path must be identical.  Acceptance requires shared/no-rotation
nearest-rank median $\leq1.25$ and 95th percentile $\leq1.50$, both inclusive.  The microbatch
runaway sentinel becomes an inclusive 100; the late-training soak limit of 10 is unchanged.  Only
the gradient sample/result/recipe/seal closure advances to schema v3.  Aggregate-v3 retains its
schema name but must recursively require the v3 gradient seal.

\paragraph{Head-scale checkout and plans.}
Implementation is isolated from the main checkout in this worktree:
\begin{quote}
\pathcode{/lustre-storage/checkpoints/ngocbh/hdn/diagnostics/rdiag_shared_dt_v1/source_worktrees/rdiag_shared_dt_headscale_20260902}.
\end{quote}
The observed checkout is branch
\pathcode{codex/rdiag-shared-dt-angle-headscale-20260902}, based on the commit above.  The clean
reviewed and published implementation snapshot is
\begin{quote}
\pathcode{a76dc7b930ae4d85157391d444c40ef7db14ef86},
\end{quote}
with tree \pathcode{4cdb6287080365089349f56b14f1d782202d6157}.
The reviewed design/plan snapshot is commit
\begin{quote}
\pathcode{9e621b5ede53529df50619804f0c9d8d32707284}.
\end{quote}
Its documents are
\begin{itemize}[leftmargin=1.8em]
  \item \pathcode{docs/plans/2026-09-02-rdiag-shared-dt-angle-headscale-design.md}; and
  \item \pathcode{docs/plans/2026-09-02-rdiag-shared-dt-angle-headscale-plan.md}.
\end{itemize}
The primary implementation and qualification paths are
\begin{itemize}[leftmargin=1.8em]
  \item \pathcode{lit_gpt/rdiag_shared_dt_kla.py};
  \item \pathcode{lit_gpt/kla_ops/rdiag_shared_dt.py} and
        \pathcode{lit_gpt/kla_ops/rdiag_bounded_phase.py};
  \item \pathcode{scripts/tests/rdiag/test_rdiag_shared_dt_production_host.py} and
        \pathcode{scripts/tests/rdiag/test_rdiag_shared_dt_production_h200_v1.py};
  \item \pathcode{scripts/tests/rdiag/test_rdiag_shared_dt_training_health_host.py};
  \item \pathcode{scripts/tests/rdiag/rdiag_shared_dt_gradient_ab.py} and
        \pathcode{scripts/tests/rdiag/rdiag_shared_dt_gradient_ab_run.py};
  \item \pathcode{scripts/train/rdiag_kla/rdiag_recipe.py},
        \pathcode{scripts/tests/rdiag/attest_rdiag_shared_dt_gate.py}, and
        \pathcode{scripts/train/rdiag_kla/rdiag_shared_dt_qualification.py};
  \item \pathcode{scripts/tests/rdiag/prepare_rdiag_shared_dt_source_gate.py} and
        \pathcode{scripts/tests/rdiag/prepare_rdiag_shared_dt_qualification.py};
  \item \pathcode{scripts/tests/rdiag/rdiag_shared_dt_production_soak.py},
        \pathcode{scripts/tests/rdiag/run_rdiag_shared_dt_production_soak.sbatch},
        \pathcode{scripts/tests/rdiag/attest_rdiag_shared_dt_production_soak.py}, and
        \pathcode{scripts/tests/rdiag/attest_rdiag_shared_dt_qualification.py}; and
  \item \pathcode{scripts/train/rdiag_kla/submit_rdiag_shared_dt.py},
        \pathcode{scripts/train/rdiag_kla/rdiag_job_control.py}, and
        \pathcode{scripts/train/rdiag_kla/rdiag_batch_supervisor.py} for the separately
        snapshotted production run and its mandatory evaluation tail.
\end{itemize}

The immutable qualification source is
\begin{quote}
\pathcode{/lustre-storage/checkpoints/ngocbh/hdn/diagnostics/rdiag_shared_dt_v1/}\\
\pathcode{source_snapshots/}\\
\pathcode{rdiag_shared_dt_auth_v3_qualification_}\\
\pathcode{a76dc7b930ae4d85157391d444c40ef7db14ef86},
\end{quote}
and the campaign root is
\begin{quote}
\pathcode{/lustre-storage/checkpoints/ngocbh/hdn/diagnostics/rdiag_shared_dt_v1/}\\
\pathcode{campaigns/}\\
\pathcode{rdiag_shared_dt_auth_v3_a76dc7b930ae4d85157391d444c40ef7db14ef86}.
\end{quote}
The failed-soak investigation is isolated in sibling worktree
\begin{quote}
\pathcode{/lustre-storage/checkpoints/ngocbh/hdn/diagnostics/rdiag_shared_dt_v1/source_worktrees/rdiag_shared_dt_soak269_diag_20260903}
\end{quote}
on branch \pathcode{codex/rdiag-shared-dt-soak269-diagnostic-20260903}, based exactly on
\pathcode{a76dc7b930ae4d85157391d444c40ef7db14ef86}.  The principal diagnostic code is
\pathcode{lit_gpt/rdiag_shared_dt_training_health.py}; the retained journal is under
\pathcode{production_soak/output/evidence/qual_rdiag_shared_dt_soak3300_a65918f1771446e7acfbffcc3e5c7dc3},
and the job logs are under \pathcode{production_soak/logs} within the campaign root.
The completed replay plan retained in that diagnostic branch is
\begin{quote}
\pathcode{docs/plans/2026-09-03-rdiag-shared-dt-soak269-diagnostic-plan.md}.
\end{quote}
Its replay entry points are \pathcode{scripts/analyses/replay_rdiag_shared_dt_soak269.py} and
\pathcode{scripts/analyses/run_rdiag_shared_dt_soak269_replay.sbatch}.

The diagnostic line progressed through initial and early-mode commits
\begin{quote}
\pathcode{f82fa86f52c0f661d0ddbc74197d1dbc82a9e8f2},\\
\pathcode{28531cf959304ab26b22789f2ea4ce082adbed13},
\end{quote}
before the final search commit and tree
\begin{quote}
\pathcode{06eaf237f5be9f79ee8d00075b79dcc0e7a9e96a},\\
\pathcode{f1ecf305edc04824ce14042fe75dbbba106f21c0}.
\end{quote}
The final immutable replay source is
\begin{quote}
\pathcode{/lustre-storage/checkpoints/ngocbh/hdn/diagnostics/rdiag_shared_dt_v1/}\\
\pathcode{source_snapshots/}\\
\pathcode{rdiag_shared_dt_alpha_zero_search_}\\
\pathcode{06eaf237f5be9f79ee8d00075b79dcc0e7a9e96a},
\end{quote}
with source-manifest SHA-256
\begin{quote}
{\ttfamily\scriptsize 81f40cb7ccd697483d7b2757e0c728da18c4ee3a99b85b54be2f63e9af51cffe}.
\end{quote}
Its replay job 1829305 wrote only below
\pathcode{diagnostic_runs/soak269_replay_06eaf237f5be_1829305} and produced neither a
qualification result nor a seal.

The then-current qualification checkout, containing the minimal lower-endpoint production-policy
correction, is isolated from both the main checkout and diagnostic harness in worktree
\begin{quote}
\pathcode{/lustre-storage/checkpoints/ngocbh/hdn/diagnostics/rdiag_shared_dt_v1/source_worktrees/rdiag_shared_dt_alpha_zero_fix_20260903}
\end{quote}
on branch \pathcode{codex/rdiag-shared-dt-alpha-zero-fix-20260903}.  It is based exactly on the
reviewed head-scaled commit.  The production-policy change itself is commit and tree
\begin{quote}
\pathcode{46ba8ec3105177a9caba2f2c153e9b71c8b62cac},\\
\pathcode{e38fdab82c68300a361613e4ee3a8c784b7c1a12}.
\end{quote}
The executable plan is
\pathcode{docs/plans/2026-09-03-rdiag-shared-dt-alpha-zero-policy-plan.md}.  Aside from adding that
plan, commit \pathcode{46ba8ec...} changes only
\pathcode{lit_gpt/rdiag_shared_dt_training_health.py}, its focused host test
\pathcode{scripts/tests/rdiag/test_rdiag_shared_dt_training_health_host.py}, and one existing case
in \pathcode{scripts/tests/rdiag/test_rdiag_shared_dt_production_h200_v1.py}; it does not change
model or kernel math.

The first qualification attempt exposed a test-construction issue rather than a kernel failure.
Its new subcase forced $DT=120$ at all 33 tokens, roughly 10.9 times the reproduced maximum, and
therefore also imposed phase stress outside the replay-observed range.  Immutable job 1829420
passed 108 of 109 tests;
the modified full-layer case passed its output comparison and finite elementwise/relative-$L_2$
hidden-gradient checks, but its hidden-gradient cosine was $0.996079<0.999$.  The gate failed
closed, result SHA-256 was
{\ttfamily\scriptsize d64821ee25f44c06aad4b6fb8d27b1433af328381be9a61d08e6e244e0a117bc},
and no seal was created.  This failed snapshot and campaign are retained as evidence and are not
reused.

Correction commit and tree
\begin{quote}
\pathcode{83118d2d96325c9b9b555d60ad4316c950c5a12d},\\
\pathcode{8c05ae8f2b5b205a17ba64042c4fad337a1d4d45}
\end{quote}
restore the ordinary strict output/all-VJP oracle case unchanged.  A separate subrun within the
same catalog item uses replay-bounded $DT=6.53125$ and rate 16, giving exponent 104.5 and exact FP32
$\alpha=0$.  It proves that Kappa receives raw zero, S3 receives the finite BT32 prefix of
$\log_2(10^{-6})$, the full forward matches the literal oracle, all model gradients remain finite,
and live Kappa/S3 cotangents are finite; no tolerance or 109-test inventory entry changed.  The
related host sweep passed 131 tests.  Development H200 job 1829576 passed the corrected item in
117.26 seconds; its stdout SHA-256 is
{\ttfamily\scriptsize 271d1ceb10772e4516cb9fe897e67cb7989e9af07939e329120a3d394f36f3a2}.
Two independent exact-SHA reviews approved this final commit without blockers, and the branch is
published.

The superseded 4228-file qualification snapshot for \pathcode{46ba8ec...} is
\begin{quote}
\pathcode{/lustre-storage/checkpoints/ngocbh/hdn/diagnostics/rdiag_shared_dt_v1/}\\
\pathcode{source_snapshots/}\\
\pathcode{rdiag_shared_dt_auth_v3_qualification_}\\
\pathcode{46ba8ec3105177a9caba2f2c153e9b71c8b62cac},
\end{quote}
with source-manifest SHA-256
\begin{quote}
{\ttfamily\scriptsize 2507dede155e9e50df876a493170d8ccfc2523413bab115e7a3f0e3369f21b93}.
\end{quote}
Failed campaign \pathcode{2f02ccb07eac4837bd6416546b17112c} is rooted at
\begin{quote}
\pathcode{/lustre-storage/checkpoints/ngocbh/hdn/diagnostics/rdiag_shared_dt_v1/}\\
\pathcode{campaigns/}\\
\pathcode{rdiag_shared_dt_auth_v3_}\\
\pathcode{46ba8ec3105177a9caba2f2c153e9b71c8b62cac}.
\end{quote}
Compute-side preparation job 1829417 passed after the login namespace correctly rejected a distinct
live-data inode identity.  Job 1829420 ran from execution-manifest SHA-256
{\ttfamily\scriptsize b23ea908ac550bc14c76267ff8e295e06f620d21412e37c766787f8ecfd72eb2}.

The now-superseded immutable qualification manifest for the corrected lower-endpoint test covers
4228 source files at
\begin{quote}
\pathcode{/lustre-storage/checkpoints/ngocbh/hdn/diagnostics/rdiag_shared_dt_v1/}\\
\pathcode{source_snapshots/}\\
\pathcode{rdiag_shared_dt_auth_v3_qualification_}\\
\pathcode{83118d2d96325c9b9b555d60ad4316c950c5a12d},
\end{quote}
with source-manifest SHA-256
\begin{quote}
{\ttfamily\scriptsize 52a82f97849f400730dfa2496472017916ca84e36dd785a7bac9e2da629936df}.
\end{quote}
That now-superseded campaign, \pathcode{a262363313d141e9bda423b0ce243940}, is rooted at
\begin{quote}
\pathcode{/lustre-storage/checkpoints/ngocbh/hdn/diagnostics/rdiag_shared_dt_v1/}\\
\pathcode{campaigns/}\\
\pathcode{rdiag_shared_dt_auth_v3_}\\
\pathcode{83118d2d96325c9b9b555d60ad4316c950c5a12d}.
\end{quote}
Compute-side preparation job 1829624 passed.  Production-H200 gate job 1829629 then completed on
H200 node \pathcode{h200-028-028} with scheduler exit \pathcode{0:0} and all 109 tests passing.  It
ran from execution-manifest SHA-256
{\ttfamily\scriptsize 98488ffdb2db208f015e34df678ef57a4003cc8ac8c29337aa7648755509e0df}.
Its retained result has SHA-256
{\ttfamily\scriptsize 32c5b2eed517f2add1e327d841c43f012858454edcf975bb5681403ec45e8675};
the JUnit report and run log have SHA-256 values
{\ttfamily\scriptsize 507bfa8de347b1967d573b41578c3c3760a23b42e8b412f4076fe7b443dcba99} and
{\ttfamily\scriptsize c6357f6d46f284dda75280ae94c53d3668e0e78bdce92abd6b5c051e83fed30d},
respectively.  Compute-side attestor job 1829796 completed successfully and published the read-only
H200 bootstrap seal with SHA-256
{\ttfamily\scriptsize 6f900869b3c13e6ef41e8920a52f9115bc31f7e69f7f8169a8b4014dc55b05f8}.
Five compute-side preparers then authenticated that seal and emitted read-only execution manifests
for train-smoke, checkpoint-fault, authenticated-resume, gradient A/B, and the 3300-update soak,
with SHA-256 values respectively
\begin{quote}
{\ttfamily\scriptsize 0d1e63001f57615c21a494e11eacfda8f04f464d258eb6d3c10207557f70fdc8},\\
{\ttfamily\scriptsize 9287308502d9e2bbffc6ff0bb2e71506c4ea9d14758a703edbef8519914157d2},\\
{\ttfamily\scriptsize 34280cd404e1680b3bc13a94bea26678cd4b24fd37bc49c4d64ad5be817ed281},\\
{\ttfamily\scriptsize 8f4fd8e74f56c3770c0464a385faac255561933b510d741d3820ddcae4ef7da7},\\
{\ttfamily\scriptsize 5f3e7a352e529666aebfbbd841c06413310b7db0a00525f45184880772e0f8ed}.
\end{quote}
Only their emitted commands were submitted, as jobs 1829974--1829978.  Train-smoke, checkpoint-
fault, authenticated-resume, and gradient A/B jobs 1829974--1829977 passed and were independently
attested.  Their seal SHA-256 values are respectively
\begin{quote}
{\ttfamily\scriptsize 472ebcf70fa24a50870b68b25024c74df296d6f35f1e218eeeb36965458dd13a},\\
{\ttfamily\scriptsize b4f44aa38f2e11397f300f88f50b03881008ebe4d19ea7afa688ce9565ea416f},\\
{\ttfamily\scriptsize 153c83d8be341bf7600990adeea69bb35e792b441459159e6fdc928316036a12},\\
{\ttfamily\scriptsize 164138bf3a3c2809d8ed4fb7dd93a7a829fa42714fd123ddd8a634b6e3bc2d11}.
\end{quote}
Gradient A/B again passed with shared/no-rotation median and 95th-percentile ratios 1.179 and 1.260.
Production-soak job 1829978 ran on nodes \pathcode{h200-207-042} and
\pathcode{h200-211-043}; its immutable-source and runtime-identity stages completed successfully.
It crossed the original failure point: updates 268, 269, and 270 all completed with zero health
failures.  At update 269, loss was 5.3710, global gradient norm was 1.2287, and 11211776 alpha
elements were below the S3 floor but remained diagnostic-only.  The run then completed 675 healthy
updates.  Its final retained record has loss 4.0055, global gradient norm 1.4537, and zero health
failures.

The next pre-optimizer check, boundary 676, failed closed with aggregate failure count 256 on all
16 ranks; update 676 was not applied.  The Python step exited 1 and the surrounding Slurm job
terminated with exit \pathcode{15:0}.  The ordinary soak evidence deliberately does not retain the
failed-key breakdown, so at that stage the responsible predicate was not identified.  No result or seal was
emitted.  The 675-line update journal, stdout, and stderr have SHA-256 values respectively
\begin{quote}
{\ttfamily\scriptsize d26a63a89430a738c689d3e1fddc2d554ce30f3c400b1258c63c2a3fd2ef31cc},\\
{\ttfamily\scriptsize 958939cf35074c8f9ca8611a19ec3dabd2830bd3090ee2e2bb104a12627dfd93},\\
{\ttfamily\scriptsize 202e1bf22ac711421f006ef01c4c9820ac1522726cb4569480f635e7217c5a28}.
\end{quote}

\paragraph{Boundary-676 diagnostic replay.}
The non-authorizing investigation is isolated in worktree
\begin{quote}
\pathcode{/lustre-storage/checkpoints/ngocbh/hdn/diagnostics/rdiag_shared_dt_v1/}\\
\pathcode{source_worktrees/rdiag_shared_dt_boundary676_diag_20260903}
\end{quote}
on branch \pathcode{codex/rdiag-shared-dt-boundary676-diagnostic-20260903}.  Its executable plan is
\pathcode{docs/plans/2026-09-03-rdiag-shared-dt-boundary676-diagnostic-plan.md}.  Diagnostic-only
health rendering was introduced at commit
\pathcode{ef4ee3c053bc1d34243f4086edd88e62b85f73d4}; the reviewed replay harness is commit and tree
\begin{quote}
\pathcode{a9ce0db738dc65c3d72e4bcd85b955815952b163},\\
\pathcode{bfb397412c1977acaef4742850738f69e948fc87}.
\end{quote}
The relevant code paths are
\pathcode{lit_gpt/rdiag_shared_dt_training_health.py},
\pathcode{scripts/analyses/replay_rdiag_shared_dt_boundary676.py},
\pathcode{scripts/analyses/run_rdiag_shared_dt_boundary676_replay.sbatch}, and
\pathcode{scripts/tests/rdiag/test_rdiag_shared_dt_boundary676_replay_host.py}.  The focused replay
suite passed 29 tests, the broad shared-DT host sweep passed 522 tests with four skips, and three
independent reviews approved the exact staged tree that became the replay commit.

The immutable 4232-file diagnostic source is
\begin{quote}
\pathcode{/lustre-storage/checkpoints/ngocbh/hdn/diagnostics/rdiag_shared_dt_v1/}\\
\pathcode{source_snapshots/rdiag_shared_dt_boundary676_diagnostic_}\\
\pathcode{a9ce0db738dc65c3d72e4bcd85b955815952b163},
\end{quote}
with source-manifest SHA-256
\begin{quote}
{\ttfamily\scriptsize 6f7e52aa691c684543be51aeda0b6ab3b8d5d43c90b753b306070af68209da45}.
\end{quote}
Replay job 1830570 ran from that snapshot with the pinned failed-run data manifest and the exact
original nodes \pathcode{h200-207-042,h200-211-043}.  It completed with scheduler exit
\pathcode{0:0}; every retained record carries \pathcode{authorizing: false}, and no qualification
result or seal was emitted.  The replay stopped naturally before optimizer boundary 649 after 648
completed updates.  Its sole fatal key was
\pathcode{alpha_at_or_above_one_count=640}, exactly five shared scalar sites expanded over width
128.  Global \pathcode{alpha_max} was exactly one.  All $DT$, $A_{\log}$, $\exp(A_{\log})$, and
decay-exponent values were finite; $DT_{\min}=1.1172518910\times10^{-8}$ and the smallest positive
decay exponent was $3.4472314780\times10^{-8}$.  Parameter and gradient nonfinite counts were zero.
Thus the captured failure is rejection of the finite FP32 identity-decay endpoint, and no monitored
nonfinite predicate fired.  The replay alone is not a Kappa/S3 oracle check; the later targeted
$K=128,T=33$ H200 subrun supplies that forward/VJP evidence.  The observed exponent is slightly
above the ideal round-to-nearest-even crossover
$-\log(1-2^{-25})=2.9802322832\times10^{-8}$; CUDA's transcendental implementation need not be
correctly rounded, so the retained output value and bitwise endpoint test, rather than the ideal
crossover alone, govern the implementation claim.

The replay is deliberately classified \pathcode{DIFFERENT_FAILURE}: it did not reproduce the old
trajectory, boundary, or count; instead it failed at boundary 649 with
\pathcode{alpha_at_or_above_one_count=640}, while the original boundary-676 key remains unknown.
Across the 648 comparable updates, maximum absolute loss and gradient-norm differences were 0.1493
and 1.3539.  This evidence does not identify the old boundary-676 key retroactively, but it
independently demonstrates that the replayed \pathcode{83118d2...} upper-endpoint policy rejected a valid
production-shaped state.  The retained artifact root is
\begin{quote}
\pathcode{/lustre-storage/checkpoints/ngocbh/hdn/diagnostics/rdiag_shared_dt_v1/}\\
\pathcode{diagnostic_runs/boundary676_replay_a9ce0db738dc_1830570}.
\end{quote}
The summary, trajectory, invocation, inner stdout, inner stderr, and artifact-inventory SHA-256
values are respectively
\begin{quote}
{\ttfamily\scriptsize 9a8c357b37731924aed61e02b7bf8ea599b943f7508cb93f91fb7bd572c89157},\\
{\ttfamily\scriptsize 7eb3ab52f083079f85059522c7b0a219e22cacdd3923c3b55df135d59b851ac7},\\
{\ttfamily\scriptsize 82848cb2fb25f4b448cebb2e87dddf197d320c7c962dbc6946e654ce6a531fa7},\\
{\ttfamily\scriptsize 859d3c15b6223acafabddecc491afd896ebaf22368b5fc54823af481a05c18f4},\\
{\ttfamily\scriptsize 645326b62f8847a96dcc47cb8ddf9b39ed214dde65d6b731131c3f8b90e8b01c},\\
{\ttfamily\scriptsize af23c4fef6684ae97102f3d81bb4cb6edc46e116623cc5e9f23fca88da78e370}.
\end{quote}
All 24 inventoried artifact hashes independently revalidated.  The terminal outer scheduler log
and empty scheduler error file have SHA-256 values respectively
\begin{quote}
{\ttfamily\scriptsize 2034750abdc9f1baea2074a0a8a28f2e6196838a1beac4fa84ad984d9ddbc151},\\
{\ttfamily\scriptsize e3b0c44298fc1c149afbf4c8996fb92427ae41e4649b934ca495991b7852b855}.
\end{quote}

The resulting narrow correction is isolated from both qualification and replay evidence in
worktree
\begin{quote}
\pathcode{/lustre-storage/checkpoints/ngocbh/hdn/diagnostics/rdiag_shared_dt_v1/}\\
\pathcode{source_worktrees/rdiag_shared_dt_alpha_one_fix_20260903}
\end{quote}
on branch \pathcode{codex/rdiag-shared-dt-alpha-one-fix-20260903}, based exactly on
\pathcode{83118d2...}.  Its plan is
\pathcode{docs/plans/2026-09-03-rdiag-shared-dt-alpha-one-policy-plan.md}.  The reviewed and
published correction commit and tree are
\begin{quote}
\pathcode{e5d07ad0f49f74e90d89a0b49ab8279b39152546},\\
\pathcode{2573a97d33244e9125ba94557b428868816360f4}.
\end{quote}
That commit changes exactly the plan and three implementation/test files
\begin{itemize}[leftmargin=1.8em]
  \item \pathcode{docs/plans/2026-09-03-rdiag-shared-dt-alpha-one-policy-plan.md};
  \item \pathcode{lit_gpt/rdiag_shared_dt_training_health.py};
  \item \pathcode{scripts/tests/rdiag/test_rdiag_shared_dt_training_health_host.py}; and
  \item \pathcode{scripts/tests/rdiag/test_rdiag_shared_dt_production_h200_v1.py}.
\end{itemize}
It retains \pathcode{alpha_at_or_above_one_count} as telemetry, adds
\pathcode{alpha_above_one_count}, and makes only strict $\alpha>1$ fatal; negative and nonfinite
checks are unchanged.  The 109-case H200 catalog gains no test node: an existing full-layer case
now includes a positive-exponent exact-one subrun covering raw FP32 bits, Kappa, the BT32 S3
prefix, literal recurrence agreement, and live finite VJPs.  Host policy tests passed 36/36, the
broad host sweep on the final pre-commit worktree passed 492 tests with four skips, and targeted H200 job 1830711
passed in 114.16 seconds.  This was non-authorizing mutable-worktree development evidence, run
before the final commit; the immutable full H200 gate is job 1830943 below.  The targeted job's
stdout and stderr SHA-256 values are
\begin{quote}
{\ttfamily\scriptsize 689c9451b67e5e9593fb15bc858723f4d4e7e89d3b47fdb90ebbce5044501442},\\
{\ttfamily\scriptsize 57ac2d0201a2657267a97ef3371038e486fe85eac5d21983843138591659f847}.
\end{quote}
Three independent exact-SHA reviews approved the correction without blockers.

The superseded \pathcode{e5d07ad...} source snapshot manifest covers 4229 source files; the
snapshot directory also contains \pathcode{SOURCE_MANIFEST.json}.  That snapshot is
\begin{quote}
\pathcode{/lustre-storage/checkpoints/ngocbh/hdn/diagnostics/rdiag_shared_dt_v1/}\\
\pathcode{source_snapshots/rdiag_shared_dt_auth_v3_qualification_}\\
\pathcode{e5d07ad0f49f74e90d89a0b49ab8279b39152546},
\end{quote}
with source-manifest SHA-256
\begin{quote}
{\ttfamily\scriptsize 3f6a8ae9176ddb9e9aee323febf6cb5d9dc22bc88d8933540be727c23a604832}.
\end{quote}
Its corresponding campaign \pathcode{7a24d88d497d4ad28b572b4ee885ac90} is closed and incomplete;
it is rooted at
\begin{quote}
\pathcode{/lustre-storage/checkpoints/ngocbh/hdn/diagnostics/rdiag_shared_dt_v1/}\\
\pathcode{campaigns/rdiag_shared_dt_auth_v3_}\\
\pathcode{e5d07ad0f49f74e90d89a0b49ab8279b39152546}.
\end{quote}
Compute-side preparation job 1830934 passed and emitted production-H200 execution-manifest
SHA-256
{\ttfamily\scriptsize 0d9277ab08b69f8d9e673e96cdd2d91ce8c753ce1d495879850edf0355377426}.
Its exact control path is
\begin{quote}
\pathcode{production_h200/control/}\\
\pathcode{production_h200_7a24d88d497d4ad28b572b4ee885ac90.manifest.json}
\end{quote}
relative to the campaign root above.  The corresponding scheduler logs are under
\pathcode{production_h200/logs/}, and retained gate artifacts are written below
\pathcode{production_h200/artifacts/production_h200_7a24d88d497d4ad28b572b4ee885ac90/}.
The resulting H200 gate, job 1830943 on \pathcode{h200-004-175}, completed with scheduler exit
\pathcode{0:0} and 32:21 wall time; pytest reported all 109 cases passed in 405.96 seconds after the
source/runtime preflight.  Its result, JUnit report, and run log have SHA-256 values respectively
\begin{quote}
{\ttfamily\scriptsize 6ac6a4ce2775230cfcfd929aa77638c968ca6841b4a863a85824e479019f67da},\\
{\ttfamily\scriptsize b6132c80118ec46578178ab567cae9a76361d30d8476d403b8071a3c154b65be},\\
{\ttfamily\scriptsize 5301d528c8b215fddca189209c17f8f3efefe7cbd831204e141ec29aed067c41}.
\end{quote}
Its outer stdout has SHA-256
{\ttfamily\scriptsize e33b84600aa1d7896723fcac68d40d342c5264914c73a052b1c608863d303f2f};
the empty stderr has SHA-256
{\ttfamily\scriptsize e3b0c44298fc1c149afbf4c8996fb92427ae41e4649b934ca495991b7852b855}.

The numerical result cannot authorize downstream work because its retained \pathcode{run.log} is
the exact 2706-byte prefix of the 2728-byte scheduler stdout: only the final 22-byte
\pathcode{PASS: production_h200} line is absent.  At \pathcode{e5d07ad...}, both source-gate
runners redirected through an asynchronous process-substitution \pathcode{tee}, froze
\pathcode{run.log}, printed PASS, and exited without waiting for that writer.  The race affected
both production-H200 and train-smoke; the attestor's requirement that both closed logs contain
PASS was correct and remains unchanged.

The first attestor allocation, job 1830944, had an operationally insufficient 20-minute wall limit.
After 5:43 it had read 34496.23 MiB of the runtime closure; the corresponding prior successful
attestor read 270926.75 MiB in 42:47, projecting this attempt well beyond its limit.  It was
therefore manually cancelled before publishing a seal; dependent preparer 1830946 never allocated
and was also cancelled.  Replacement attestor job 1831239 used the byte-identical command and
evidence pins with a one-hour limit.  It completed the full closure read and then correctly failed
\pathcode{2:0} after 57:28 on the missing retained-log marker; it emitted no seal.  Dependent
preparer 1831240 never allocated.  A redundant fail-safe attestor, job 1831513, was cancelled after
the deterministic artifact defect was identified rather than spending another full validation
cycle.

The logging correction is isolated in worktree
\begin{quote}
\pathcode{/lustre-storage/checkpoints/ngocbh/hdn/diagnostics/rdiag_shared_dt_v1/}\\
\pathcode{source_worktrees/rdiag_shared_dt_source_gate_log_fix_20260903}
\end{quote}
on branch \pathcode{codex/rdiag-shared-dt-source-gate-log-fix-20260903}, based exactly on
\pathcode{e5d07ad...}.  Its executable plan is
\begin{quote}
\pathcode{docs/plans/}\\
\pathcode{2026-09-03-rdiag-shared-dt-source-gate-log-finalization-plan.md}.
\end{quote}
\paragraph{Working-location index.}
The worktree above is the historical qualification-source checkout, on branch
\pathcode{codex/rdiag-shared-dt-source-gate-log-fix-20260903} at the exact reviewed commit
\pathcode{1fce1b9...}; current diagnostic and control development is isolated in the
generation-zero, forward-determinism, RMSNorm-control, trigger-capture, and matched-intervention
worktrees named below.  The main checkout is not the source of any qualification job.  Within the
historical source worktree, the shortest paths into
the implementation and its controls are:
\begin{itemize}[leftmargin=1.8em]
  \item \pathcode{lit_gpt/rdiag_shared_dt_kla.py} for the layer parameterization and the
        $H^{-1/2}$ angle head scale;
  \item \pathcode{lit_gpt/kla_ops/rdiag_shared_dt.py} and
        \pathcode{lit_gpt/kla_ops/rdiag_shared_dt_reference.py} for the composed GPU path and
        reference recurrence, with phase/clock helpers beside them;
  \item \pathcode{lit_gpt/rdiag_shared_dt_training_health.py} for the alpha endpoint telemetry and
        fatal-domain policy;
  \item \pathcode{scripts/tests/rdiag/} for H200 cases, distributed smoke/soak/fault gates, runners,
        and attestors; and
  \item \pathcode{scripts/train/rdiag_kla/} for recipe pinning, campaign/job control,
        authorization, and the production submitter.
\end{itemize}
The design and implementation plans for the numerical correction are, relative to that worktree,
\begin{itemize}[leftmargin=1.8em]
  \item \pathcode{docs/plans/}\\
        \pathcode{2026-09-02-rdiag-shared-dt-angle-headscale-design.md}; and
  \item \pathcode{docs/plans/}\\
        \pathcode{2026-09-02-rdiag-shared-dt-angle-headscale-plan.md}.
\end{itemize}
The two subsequently discovered finite-precision endpoint policies are recorded in the
corresponding \pathcode{alpha-zero-policy} and \pathcode{alpha-one-policy} plans dated
2026-09-03.  The boundary-269 and boundary-676 plans live only in their respective historical
diagnostic worktrees.  Those filesystem-writable worktrees are retained for investigation but are
not qualification sources; the current immutable source snapshot and campaign root are named
below.

The only intended runtime changes are
\pathcode{scripts/tests/rdiag/run_rdiag_shared_dt_production_h200_v1.sbatch} and
\pathcode{scripts/tests/rdiag/run_rdiag_shared_dt_train_smoke.sbatch}; regression coverage is in
\pathcode{scripts/tests/rdiag/test_rdiag_shared_dt_source_gate_host.py}.  The corrected lifecycle is
PASS emission, closure of both pipe writers, an explicit wait for the captured \pathcode{tee} PID,
and only then read-only log publication.  It preserves any earlier nonzero status and makes a
logger or chmod failure fatal.  No result, seal, aggregate, or attestor schema is weakened.

The reviewed and published logging correction is commit/tree
\begin{quote}
\pathcode{1fce1b9a8e4a80d4da82a689dfb69983109aad0c},\\
\pathcode{297b4b1f3f4ba7e4ac64ac35df0510c9bdea69be}.
\end{quote}
It changes exactly the two runners and the source-gate host test named above, plus the tracked plan.
A focused selection over \pathcode{test_rdiag_shared_dt_source_gate_host.py},
\pathcode{test_rdiag_shared_dt_production_h200_v1_host.py},
\pathcode{test_rdiag_shared_dt_train_smoke_host.py}, and
\pathcode{test_rdiag_shared_dt_gate_runners_host.py} passed 134 tests.  The full 22-file RDiag host
sweep passed all 823 tests, and independent collection still reports exactly 109 H200 cases.  Two earlier broad
attempts were deliberately excluded: concurrent test artifacts filled the 512 MiB
\pathcode{/var/tmp} filesystem in one, while placing temporary trees on Lustre changed the
filesystem-specific test branch in the other.  The final sweep used a spacious local filesystem
and had no failures.  Three independent exact-SHA reviews approved the commit without blockers.

The replacement immutable source snapshot contains 4230 tracked regular files and excludes only
the approved \pathcode{paper} gitlink.  It is
\begin{quote}
\pathcode{/lustre-storage/checkpoints/ngocbh/hdn/diagnostics/rdiag_shared_dt_v1/}\\
\pathcode{source_snapshots/rdiag_shared_dt_auth_v3_qualification_}\\
\pathcode{1fce1b9a8e4a80d4da82a689dfb69983109aad0c},
\end{quote}
with source-manifest SHA-256
\begin{quote}
{\ttfamily\scriptsize 429348ceb16350782bc0eb2b4a34affabc5eb73f1d3ac997e5da42a2d1949067}.
\end{quote}
Fresh campaign \pathcode{ae2041e3d1fd4b1ca7c1d609184c2642} is rooted at
\begin{quote}
\pathcode{/lustre-storage/checkpoints/ngocbh/hdn/diagnostics/rdiag_shared_dt_v1/}\\
\pathcode{campaigns/rdiag_shared_dt_auth_v3_}\\
\pathcode{1fce1b9a8e4a80d4da82a689dfb69983109aad0c}.
\end{quote}
Compute-side preparation job 1832038 completed with execution-manifest SHA-256
\begin{quote}
{\ttfamily\scriptsize a1759ea8f8b058b4ab49b235d7ce0d083d7760cc41447d0d9447e01515ee925e}.
\end{quote}
The manifest is
\begin{quote}
\pathcode{production_h200/control/}\\
\pathcode{production_h200_ae2041e3d1fd4b1ca7c1d609184c2642.manifest.json}
\end{quote}
relative to that campaign root.  Its authenticated emitted command was checked as an exact argv
sequence inside H200 allocation 1832039 and submitted as job 1832040.  That gate completed
\pathcode{0:0} on \pathcode{h200-052-125} in 31:54; all 109 cases passed in 423.33 seconds.  The
result, JUnit report, and retained log have SHA-256 values respectively
\begin{quote}
{\ttfamily\scriptsize e39ebe47ce0e1bd43f23c312c75269a5a6300b6f367f6f43df7906568bc8d057},\\
{\ttfamily\scriptsize a654c778be9062d4cc9f41b62b0f779732f09bf0c5e9c27521538384c1c9a732},\\
{\ttfamily\scriptsize 5335a009a8407fcf2965cc67eb09f02124d562d378c98019e05119cf57e1401d}.
\end{quote}
The mode-0400 retained log is byte-identical to scheduler stdout, includes exactly one terminal
\pathcode{PASS: production_h200} marker, and the scheduler error file is empty (SHA-256
{\ttfamily\scriptsize e3b0c44298fc1c149afbf4c8996fb92427ae41e4649b934ca495991b7852b855}).
This directly exercises and resolves the former publication race.  The two-hour external attestor
is job 1832042; it completed \pathcode{0:0} in 52:16 and published a mode-0400 bootstrap seal with
SHA-256
\begin{quote}
{\ttfamily\scriptsize a6c63dae54416ad4085631de5e7d27ee97423ad881a5d1f75596d68d4dc35cff}.
\end{quote}
Five manifest-only preparation jobs were dependency-bound to that attestor: production soak
1832100, train smoke 1832101, checkpoint fault 1832102, authenticated resume smoke 1832103, and
gradient A/B 1832104.  All five completed \pathcode{0:0}; their stderr logs are empty and their
mode-0400 execution-manifest SHA-256 values, in that order, are
\begin{quote}
{\ttfamily\scriptsize 2039a4b702586e5549c6a92bcfb0d0b3054151df8af290086e153439674c9873},\\
{\ttfamily\scriptsize 36c6b0a2597dd6ee31d95b048d9434eb6a021f2f8a65a26736fa4fca6822e0df},\\
{\ttfamily\scriptsize f2ae3776688b6a3c06e6fd7b2f6e1e71b303e50663fb2b296bebf799797fb63f},\\
{\ttfamily\scriptsize 5ca194e4c54c7342c498689000d377acbf283e1e5a0be9075e7b9116c6e9a15c},\\
{\ttfamily\scriptsize db126e731f9bae7cdd1172fa4815a966dc49084ecb21310ca0026dfa30e5d119}.
\end{quote}
Control allocation 1832771 independently rehashed each manifest, loaded it through the pinned
control code, rebuilt the scheduler argv inside an H200 allocation, required exact equality with
the preparer-emitted argv, and found no existing job with any of the five unique scheduler
comments.  This compute-side step was required because the login node intentionally has a
different data-root device identity.  It submitted the exact workloads and dependency-bound
external attestors as the following pairs:
\begin{quote}
\pathcode{production-soak 1832774 -> 1832775},\\
\pathcode{train-smoke 1832776 -> 1832777},\\
\pathcode{checkpoint-fault 1832778 -> 1832779},\\
\pathcode{authenticated-resume 1832780 -> 1832781},\\
\pathcode{gradient-A/B 1832782 -> 1832783}.
\end{quote}
The four short workloads and attestors then completed successfully:
\begin{itemize}[leftmargin=1.8em]
  \item train-smoke jobs \pathcode{1832776 -> 1832777}, seal SHA-256\\
        {\ttfamily\scriptsize d04cbec0b679496a8d4e05c22435c919ce1dbce455ea1e08dd48bbc04abad556};
  \item checkpoint-fault jobs \pathcode{1832778 -> 1832779}, seal SHA-256\\
        {\ttfamily\scriptsize d382d2ef31dbb49588316a7850f8c022ee3796681ee6aa418e10a53f0c0902aa};
  \item authenticated-resume jobs \pathcode{1832780 -> 1832781}, seal SHA-256\\
        {\ttfamily\scriptsize dc1ab93029eca699c27d4d1a8ac23a946cc6ff2c76c95d30de77edd82d897e23}; and
  \item gradient-A/B jobs \pathcode{1832782 -> 1832783}, seal SHA-256\\
        {\ttfamily\scriptsize a17fb2ecc2f116c41c7b333853571a8665b4811a23685b8530fea071b1e12e2b}.
\end{itemize}
All eight scheduler exits are \pathcode{0:0}; every external-attestor stderr is empty, and each
independently rehashed seal is an owner-held, single-link, mode-0400 regular file.  The new
train-smoke retained log is byte-identical to scheduler stdout and contains exactly one terminal
PASS marker, providing a second end-to-end exercise of the log-finalization fix.  Checkpoint-fault
and authenticated-resume emitted the same framework initialization, FSDP deprecation, and
process-group shutdown warnings seen in the earlier successful campaign; neither emitted a Python
traceback, and the external attestors accepted their exact results.

The two-node soak allocation 1832774 started at 2026-09-03 20:15:33 UTC on
\pathcode{h200-232-233} and \pathcode{h200-233-046}.  Slurm emitted two bootstrap warnings because
it probed the exported job-private \pathcode{TMPDIR} before the first per-node bootstrap step had
created that directory, and therefore used \pathcode{/tmp} for that step; the runner then created
and used its intended mode-0700 cache trees.  This is retained scheduler stderr, not the training
failure.

The journal crossed both historical failure boundaries with zero aggregate health failures.  Its
exact update-269 row records loss 5.395138889551163, global gradient norm 1.213560938835144, and
26,092,800 alpha values below the diagnostic S3 floor out of 17,179,869,184.  Its update-676 row
records loss 4.053754523396492, global gradient norm 1.0578112602233887, and 102,754,304
below-floor values out of the same denominator.  The workload then failed closed at the
pre-optimizer boundary for update 790.  All 16 ranks reported
\pathcode{SoakEvidenceError: global gradient exceeds the soak limit}; job 1832774 ended
\pathcode{FAILED 1:0}, and dependency-bound attestor 1832775 was cancelled without allocation.
No soak result, final checkpoint, seal, or aggregate exists.

The retained journal contains exactly 789 contiguous accepted rows and has SHA-256
\pathcode{4814e6c0cedbdc4efdceefb4b944ca2e8b1df357d4f7a49118857a64f8da286f}.
Its median and lower-empirical 95th and 99th percentile pre-clip norms are respectively 1.14646,
2.27364, and 3.35581.  The maximum is 7.970221519470215 at update 789; excluding that row, the
maximum is 4.06596 at update 24.  Updates 783--789 measure 3.39357, 1.82029, 1.61033, 2.27467, 2.26284,
3.04553, and 7.97022, followed by the unretained update-790 rejection.  Loss at the rejected
boundary was finite, and every accepted row has zero health failures.

The name \pathcode{after_clip} is temporal rather than semantic.  Lightning delegates this call
to FSDP, which returns the accumulated distributed norm \emph{before} multiplying gradients in
place by approximately \pathcode{min(1, 1/(norm+1e-6))}.  Thus update 789 was subsequently clipped
to approximately one; update 790 had a rank-equal pre-clip norm that was nonfinite, negative, or
above 10 and was rejected before \pathcode{optimizer.step}.  Its exact scalar and predicate branch
are irrecoverable from this run: the harness checks the combined predicate before constructing or
appending a record, and named-gradient sampling begins only at update 3000.  The run retained only
its generation-zero checkpoint.

This is a genuine failure of the frozen qualification contract, but not yet proof that the 750M
model is unstable.  The limit 10 was copied from the older bounded-E2 soak.  That accepted E2
screen observed a sampled maximum 1.888 over updates 300--500, but it did not calibrate the tail of
a completed 3300-update shared-DT run.  Conversely, the separate 220M gradient-v3 gate uses 100
only as a gross fail-fast ceiling around a matched shared/no-rotation ratio experiment; it cannot
justify changing the 750M soak threshold to 100.  The late rise makes attribution mandatory before
either a model change or a versioned gate-policy change.

The non-authorizing investigation is isolated in worktree
\begin{quote}
\pathcode{/lustre-storage/checkpoints/ngocbh/hdn/diagnostics/rdiag_shared_dt_v1/}\\
\pathcode{source_worktrees/rdiag_shared_dt_soak790_diag_20260903}
\end{quote}
on branch \pathcode{codex/rdiag-shared-dt-soak790-diagnostic-20260903}, branched from
\pathcode{1fce1b9...}.  The reviewed and published execution commit is
\pathcode{47bb7812da1bd5809d974ea9c5504209f3a357f3}, tree
\pathcode{2fc10ee5320a418a72960a86145083916a3d739c}.  Post-execution evidence hardening is the current
branch head:
\begin{quote}
commit \pathcode{4091b2572f19f8cefaab8932b2d035969f663ec0},\\
tree \pathcode{bf8abd733df27ba03331fadf35560107fadf1e8b}.
\end{quote}
It does not alter the executed source.  Relative
to that worktree, the completed replay plan is
\pathcode{docs/plans/}\\
\pathcode{2026-09-03-rdiag-shared-dt-soak790-diagnostic-plan.md}; its entry points are
\pathcode{scripts/analyses/replay_rdiag_shared_dt_soak790.py} and
\pathcode{scripts/analyses/run_rdiag_shared_dt_soak790_replay.sbatch}.  The replay keeps the
production recipe, ordered data, seed, optimizer schedule, 2-by-8 topology, FP32 FSDP buffers,
clipping, and strict limit under measurement-only instrumentation.  This is a same-recipe replay,
not bit-exact production reproduction: read-only scans, copies, and reductions can perturb GPU
scheduling, and ``same boundary'' denotes only the boundary number.  It revalidates the failed
run's complete runtime fingerprint, computes per-rank target-token hashes and pre-clip category and
top-parameter diagnostics at every boundary, and retains the hashes for a gradient-rejection
boundary.  It stops non-authorizingly on a natural production-health failure, a finite norm above
10, a nonfinite returned norm or post-health named gradient, or exhaustion through update 820.
Nonfinite values use canonical JSON-safe tokens, finite-only norms are labeled as such, and every
offending parameter remains separately listed.
The validator requires the exact 32-field PRE-health inventory, the full-update alpha denominator
17,179,869,184, and name-to-category agreement for every retained contributor.  Physical-family
nonfinites remain owned by the earlier production-health failure path.

The final pre-execution replay, health, and pretrain-integration host suite passed 87 tests.  Independent
runner, schema, and mathematical-contract reviews approved execution commit \pathcode{47bb7812...}.
A Git-object
snapshot was materialized at
\begin{quote}
\pathcode{/lustre-storage/checkpoints/ngocbh/hdn/diagnostics/rdiag_shared_dt_v1/}\\
\pathcode{source_snapshots/}\\
\pathcode{rdiag_shared_dt_soak790_diag_47bb7812da1bd5809d974ea9c5504209f3a357f3}
\end{quote}
with source-manifest SHA-256
{\ttfamily\scriptsize a948eb8a0540a2ff7625f344372a2d04389865bb5e3a977ad7ceff123e4a742f}.
The manifest closes 4,234 tracked regular files; its
\pathcode{SOURCE_MANIFEST.json} is stored at the snapshot root.
Two-node H200 replay job \pathcode{1833926} ran on \pathcode{h200-101-019} and
\pathcode{h200-208-076} from 2026-09-03 22:40:57 through 23:02:14 UTC.  It and all Slurm steps
completed \pathcode{0:0}.  The rank-complete result is
\pathcode{NO_GRADIENT_LIMIT_EXCEEDANCE_THROUGH_820}: all 820 returned norms were finite,
nonnegative, and no larger than 10; all named-gradient nonfinite counts and all health-failure
counts were zero.  Update 790 measured norm 1.0757828950881958 and loss 4.024734705686569, update
820 measured norm 4.972261428833008, and the replay maximum was 7.394214630126953 at update 817.
Thus the replay did not reproduce the original update-790 rejection.

The replay is not a matched trajectory.  Of the 789 comparable records, 673 differ in at least one
shared field; floating-point norms differ at update 1, alpha-below-floor counts first differ at
update 3, and every update from 139 through 789 differs.  Original versus replay update 789 is
respectively 7.970221519470215 versus 1.4658762216567993.  Gradient norms over 1--789 have Pearson
correlation 0.5998 and first-difference correlation 0.0389; over 1--99 those values are 0.999565 and
0.99191, while update 100 starts the first five-update run with absolute norm differences above
0.25.  Loss remains tightly same-step coupled (Pearson 0.999511), showing the same broad recipe but
not the same numerical state.  No original update-790 input hashes exist, and this no-rejection
outcome retained no replay batch hashes, so data equality at that boundary is unproved.

The late replay excursion is broad rather than an isolated radial failure.  Across updates
783--820, squared gradient-norm mass is 45.57\% ordinary/backbone, 32.91\% angle, 21.50\% DT
weight, 0.026\% DT bias, and 0.00085\% \pathcode{A_log}.  At the maximum, update 817, category norms
are 5.2138 ordinary/backbone, 4.5117 angle, 2.6695 DT weight, 0.0909 DT bias, and 0.0271
\pathcode{A_log}.  This localizes co-amplification but is not a causal decomposition; it supports
neither removing rotation nor changing the frozen limit.

The finalized artifacts are under
\begin{quote}
\pathcode{/lustre-storage/checkpoints/ngocbh/hdn/diagnostics/rdiag_shared_dt_v1/}\\
\pathcode{diagnostic_runs/soak790_replay_47bb7812da1b_1833926/artifacts}.
\end{quote}
The retained evidence digests are:
\begin{quote}
summary: {\ttfamily\scriptsize 9620387f55d38bfde32e5763b1ecf72c1c4fe23223ea6d01d051acb246a718f5},\\
trajectory: {\ttfamily\scriptsize d8fbb3c1fdea2bc1f63f089cb457b9ee832372cd828465fbee2070327801a7fc},\\
invocation: {\ttfamily\scriptsize ff2a6a83950942e2d7ba12d8ac1bac163264e628f6d34a97e010ab27ad644634}.
\end{quote}
An independent audit found one packaging defect: the invocation named Slurm's ephemeral spool copy
of the runner, which disappeared after cleanup.  Its recorded SHA-256 and size exactly match the
immutable snapshot runner, and substituting that path regenerates the summary byte-for-byte, so
the result is corroborated but is not self-validating as originally packaged.  Commit
\pathcode{4091b2572f19f8cefaab8932b2d035969f663ec0} fixes future packages by retaining a mode-0400 artifact-local
\pathcode{runner.sbatch}, binding it to the retained source-manifest entry, and rejecting external
or writable runner records.  Its focused suite passes 48 tests.  There is no evidential reason to
rerun job 1833926 merely to repair this historical packaging defect.

The shortest code paths for this boundary are \pathcode{pretrain.py} (data fetch, health, clip,
optimizer, and checkpoint lifecycle), \pathcode{data.py} (stateful worker/buffer order),
\pathcode{lit_gpt/rdiag_checkpoint.py} and \pathcode{lit_gpt/fsdp_strategy.py} (full model/optimizer
and per-rank sidecars), and
\pathcode{scripts/tests/rdiag/rdiag_shared_dt_gradient_ab_run.py} (the existing connected-zero
rotation intervention).  Replay-specific instrumentation remains confined to the helper and runner
named above.

\paragraph{Remaining work.}
The head scale and gradient-v3 closure remain frozen and published at
\pathcode{a76dc7b930ae4d85157391d444c40ef7db14ef86}; the alpha-zero policy source is frozen and
published at \pathcode{83118d2d96325c9b9b555d60ad4316c950c5a12d}.  That source crossed the former
boundary-269 failure but did not pass the full soak.  The alpha-one correction is frozen at
\pathcode{e5d07ad0f49f74e90d89a0b49ab8279b39152546}; the failed source head is its
log-finalization child \pathcode{1fce1b9a8e4a80d4da82a689dfb69983109aad0c}.  Every existing seal
is source-bound historical evidence; the five successful non-soak seals cannot authorize a
source or policy change.

The formerly proposed production root for \pathcode{1fce1b9...} remains deliberately absent:
\begin{quote}
\pathcode{/lustre-storage/checkpoints/ngocbh/hdn/production/}\\
\pathcode{rdiag_shared_dt_auth_v2_1fce1b9a8e4a80d4da82a689dfb69983109aad0c}.
\end{quote}
It will not be created from a failed qualification.  A future production root, if authorized from
a later source commit, will have fixed children \pathcode{source_snapshot}, \pathcode{manifests},
\pathcode{control}, \pathcode{output}, and \pathcode{logs}.  The separate production snapshot must
be created atomically by the immutable launcher's \pathcode{prepare-fresh} mode; the qualification
snapshot is not reused as the training snapshot.  The four recipe digests listed below are
historical pins for \pathcode{1fce1b9...}, not reusable authorization for changed source:
\begin{itemize}[leftmargin=1.8em]
  \item \pathcode{rdiag_kla_shared_dt_v1_220M}:\\
        {\ttfamily\scriptsize 8f4d2d21f38ee93bda045e74d80bf6c5c1dd1feef9102e3ab413583b1542d9e1};
  \item \pathcode{rdiag_kla_shared_dt_v1_750M}:\\
        {\ttfamily\scriptsize 7a07b7af54a9df20736fc85c2eb615900c969c8a977ca12836cee210f301b2c4};
  \item \pathcode{rdiag_kla_shared_dt_v1_1.3B}:\\
        {\ttfamily\scriptsize f11e601a6a60fe67dd069ec5fc72954c2f32396bda4d910a32b13ed34f377834}; and
  \item \pathcode{swa_rdiag_kla_shared_dt_v1_1.3B}:\\
        {\ttfamily\scriptsize 225f1cd89233f0a290c527ad5c59ce4b71f9d884dfcdb1a6569044a93bd96b3f}.
\end{itemize}
The launcher credential prerequisite has been materialized through the no-argument/no-stdout
credential path at \pathcode{/storage/home/ngocbh/.wandb-key}; it is an owner-controlled,
single-link regular file with mode 0400.  No separate production-training snapshot, execution
manifest, output directory, or training job has been created.
Because update 790 was not reproduced, the next investigation first locates the earliest divergence
from the exact original generation-zero checkpoint.  It is isolated in worktree
\begin{quote}
\pathcode{/lustre-storage/checkpoints/ngocbh/hdn/diagnostics/rdiag_shared_dt_v1/}\\
\pathcode{source_worktrees/rdiag_shared_dt_gen0_repro_diag_20260903}
\end{quote}
on branch \pathcode{codex/rdiag-shared-dt-gen0-repro-diagnostic-20260903}.  Its executable plan is
\pathcode{docs/plans/}\\
\pathcode{2026-09-03-rdiag-shared-dt-gen0-reproducibility-plan.md}.  It pins the original 2.9-GB
generation-zero model with SHA-256
\begin{quote}
{\ttfamily\scriptsize 8a37e667678d2743c5bdcfb3b808c909f980c50b906e2ff9e97e29b02074b581}
\end{quote}
and all 16 rank sidecars.  It runs three six-microbatch loader controls: no snapshot (the 1833926
path), snapshot-and-continue on the same loader (the failed fresh-production path), and new-loader
sidecar restore.  Four common-anchor training arms
run in counterbalanced \pathcode{minimal-A1/instrumented-B1/instrumented-B2/minimal-A2} order for
three updates on the same allocation and GPU mapping.  Complete local parameter/optimizer Merkle
inventories and exact input hashes are prerequisites.  Historical process RNG was not checkpointed,
so this test localizes only the first observed divergence among captured stages; it cannot alone
identify the unique historical cause.  This experiment is diagnostic only.
The generation manifest and its sidecar/ordered-file composites are pinned respectively to
\begin{quote}
{\ttfamily\scriptsize 6649a1a68f83b80df9f5600864dc52e51639d19b42a827b78c9eeacef9e9f8f6},\\
{\ttfamily\scriptsize 162ff05e5940a427504c0075c82002aa8894091a8bd4f1c22c5bb5e0cb69172f},\\
{\ttfamily\scriptsize 6a184a861f445a0eb9ed383e029d726f9cec78f97285df3c350fda95bc52fd86}.
\end{quote}
The validator requires exact zero counters, modes, identities, and the 17-file model-plus-sidecar
closure before any arm runs.

Implementation and review are complete in the worktree above.  The reviewed execution identity is
\begin{quote}
commit \pathcode{4804be477a09f19adbd88d5ca773be721608d9f3},\\
tree \pathcode{b3d882ee859038366f6d18cbd612ad91aa2d31d3}.
\end{quote}
It is based on diagnostic-evidence head
\pathcode{4091b2572f19f8cefaab8932b2d035969f663ec0}.  Relative to the worktree, its executable entry
points and host contract are
\begin{itemize}[leftmargin=1.8em]
  \item \pathcode{scripts/analyses/reproduce_rdiag_shared_dt_gen0.py};
  \item \pathcode{scripts/analyses/run_rdiag_shared_dt_gen0_repro.sbatch}; and
  \item \pathcode{scripts/tests/rdiag/test_rdiag_shared_dt_gen0_repro_host.py}.
\end{itemize}
Pre-execution review found and corrected four material harness defects: the post-run snapshot
verifier omitted its required private-copy argument; the terminal completion marker preceded
verified node-local cleanup; an instrumented-arm wrapper discarded ordinary loss capture; and the
classifier could ignore some global-loss, learning-rate, named-gradient, rank-agreement, and
loader-to-training input discrepancies.  The corrected helper now rederives every field used for
classification, binds the restored L2 input matrix to all training arms, requires finite rank-equal
pre-clip norms, distinguishes minimal- and instrumented-pair collective differences, and validates
an exact 214-file artifact closure before completion.  The exact reviewed helper, runner, and test
SHA-256 values are respectively
\begin{quote}
{\ttfamily\scriptsize 37b0e2061a145b2b1935a74a362235a1d462457119ed42e4e1efd3d36b7be7bc},\\
{\ttfamily\scriptsize f9d5c8a04f7384b360eb04435c32d0a2abeb177e6f6534e91f4910e5b7c1bb14},\\
{\ttfamily\scriptsize 750d0bae60ca24c7ae7fecffc0a52c2decba915854b5f282910b038bd088bfdc}.
\end{quote}
The focused generation-zero suite plus the retained soak-replay suite passes 94 tests; Python
compilation, batch-shell parsing, staged whitespace checks, and a fresh exact validation of all 17
generation-zero inputs also pass.  Independent artifact/runner, mathematical-contract, and
Lightning/FSDP runtime reviews found no remaining Critical or Important issue on these exact
bytes.  The broader RDiag host sweep first reached 898 passes before 19 late fixture writes failed
because the login node's 1-GB \pathcode{/tmp} filled; these were \pathcode{ENOSPC} errors rather
than assertion failures.  Re-running both affected files with a diagnostic-Lustre temporary root
passed 82/82, closing all 917 distinct broad-suite tests.  The two disposable test-owned temporary
trees were then removed.

Commit \pathcode{4804be477...} is pushed on the branch named above.  Its immutable Git-object
snapshot is
\begin{quote}
\pathcode{/lustre-storage/checkpoints/ngocbh/hdn/diagnostics/rdiag_shared_dt_v1/}\\
\pathcode{source_snapshots/}\\
\pathcode{rdiag_shared_dt_gen0_repro_diag_4804be477a09f19adbd88d5ca773be721608d9f3}.
\end{quote}
The snapshot contains 4,238 tracked regular files, excludes only the external
\pathcode{paper} gitlink, and was independently revalidated against source-manifest SHA-256
\begin{quote}
{\ttfamily\scriptsize 1bbf00a65b30ffd26bd01926e09b7c306c2836054f40d83aab9963904e16c893}.
\end{quote}
Two-node H200 diagnostic job \pathcode{1834859} ran directly from that snapshot on
\pathcode{h200-091-173} and \pathcode{h200-101-011} from 2026-09-04 01:23:02 through 01:28:07
UTC and failed \pathcode{15:0} in the first loader probe.  Source staging, both runtime identities,
invocation capture, and the complete generation-zero input preflight succeeded.  Before any
training update or rank result was written, all ranks rejected the first batch at
\pathcode{batch_record}: the live Transformers collator returns a \pathcode{BatchEncoding}, which
implements \pathcode{Mapping} and contains \pathcode{input_ids}, \pathcode{attention_mask}, and
\pathcode{labels}, but the diagnostic required the concrete type \pathcode{dict}.  No summary,
artifact inventory, completion marker, or scientific result exists.  The retained scheduler log
and stderr SHA-256 values are respectively
\begin{quote}
{\ttfamily\scriptsize 29939733705f0f3e49dd2fe2e1210484f68a3d6b90f6fafe361fb992cc7718a2},\\
{\ttfamily\scriptsize 4f5e6fb8ab374b4bce06213cab824ed588b5ee4ae5c6ad818ce4b2f10ee34326}.
\end{quote}
The failure is a fail-closed diagnostic type-contract bug, not model evidence.

The corrected execution identity is
\begin{quote}
commit \pathcode{f1cfd305cfe76039cf8a0d62881e84ec5faee6d7},\\
tree \pathcode{5c558d3e470b3c5f47aa7912f8ffed71bcbd5255}.
\end{quote}
It replaces that concrete-type check with
\pathcode{isinstance(batch, Mapping)}, documents that the retained identity is the two-field
training-input projection, and tests an actual three-field \pathcode{BatchEncoding} plus negative
non-mapping/missing-field cases.  The focused generation-zero plus replay suite now passes 97
tests.  An independent review reproduced the real collator type and found no analogous live-API
type blocker.  The corrected helper SHA-256 is
\begin{quote}
{\ttfamily\scriptsize 9e9c20f71fb5810ae39da45e29dddaa2c7b130e4b5b7740e819a9a2664cbb63c}.
\end{quote}
The runner remains byte-identical.  The replacement 4,238-file immutable snapshot is
\begin{quote}
\pathcode{/lustre-storage/checkpoints/ngocbh/hdn/diagnostics/rdiag_shared_dt_v1/}\\
\pathcode{source_snapshots/}\\
\pathcode{rdiag_shared_dt_gen0_repro_diag_}\\
\pathcode{f1cfd305cfe76039cf8a0d62881e84ec5faee6d7},
\end{quote}
independently revalidated against source-manifest SHA-256
\begin{quote}
{\ttfamily\scriptsize 524164d0e3b75e91b9cf3013ac581b0d0a12b8288aeec8ec33579efb6eb87a55}.
\end{quote}
Replacement two-node H200 job \pathcode{1835101} ran directly from that snapshot on
\pathcode{h200-081-214} and \pathcode{h200-199-061} from 2026-09-04 01:57:03 through 02:11:54
UTC and failed \pathcode{15:0}.  Unlike job \pathcode{1834859}, all three loader modes completed:
48 rank records prove that the two-field training-input projections of all six microbatches match
exactly between no-snapshot, snapshot-and-continue, and new-loader sidecar-restore paths on every
rank.  The L1 and L2 canonical loader-state hashes are also identical.  Thus neither calling
\pathcode{state_dict} nor restoring its captured state is the first observed loader fork in this
environment.

The first training arm then failed on all ranks before model construction or any update because
the shared-DT frontend import loads TileLang, whose TVM initialization intentionally prepends its
Python directory to the in-process \pathcode{PYTHONPATH}.  The harness validated the pristine
launch environment after that import and therefore rejected the expected library mutation.  No
training-arm record, summary, artifact inventory, completion marker, or scientific classification
exists.  The retained artifact root is
\begin{quote}
\pathcode{/lustre-storage/checkpoints/ngocbh/hdn/diagnostics/rdiag_shared_dt_v1/}\\
\pathcode{diagnostic_runs/gen0_repro_f1cfd305cfe7_1835101/artifacts}.
\end{quote}
Its scheduler log, stderr, invocation, and pre-run checkpoint-inventory SHA-256 values are
respectively
\begin{quote}
{\ttfamily\scriptsize 1d1770bc94f23f14a4d3d97083827130031eec773131a7512fd4ce2abe350c8f},\\
{\ttfamily\scriptsize ccb4628c9a94d73a0f324f748c106c38bc44b846b3e797eef809febb6d3ac9b9},\\
{\ttfamily\scriptsize b13029c37f098209ec5eb046e78ecd4652c38af6998fbe9e88fe8caf7be8b45c},\\
{\ttfamily\scriptsize 20bbc0c2a345eaa46671d670c0d865b6f9a73266c9310bd81a1f43243ed120fd}.
\end{quote}

The second correction moves the unchanged launch-environment validation before heavyweight model
imports; it does not weaken the check.  It is reviewed and pushed as
\begin{quote}
commit \pathcode{e2b5e317b750e88437d30e95eb6a89ac88cb11c2},\\
tree \pathcode{f625647a419c9ded1dd02a9b5367827981b00b18}.
\end{quote}
The helper and host-test SHA-256 values are now
\begin{quote}
{\ttfamily\scriptsize fc3c44e57af6d0994934a7bd0c7433580904a481428966773e6863b77adf2217},\\
{\ttfamily\scriptsize 4be33408b0e33743486be389d31968f3fd6337c83d58f823a9d1ce33860f803f}.
\end{quote}
The combined generation-zero/replay host suite passes 98 tests; an independent import-chain audit
reproduced the TileLang mutation and found no earlier mutating import.  A fresh 4,238-file
Git-object snapshot is
\begin{quote}
\pathcode{/lustre-storage/checkpoints/ngocbh/hdn/diagnostics/rdiag_shared_dt_v1/}\\
\pathcode{source_snapshots/rdiag_shared_dt_gen0_repro_diag_}\\
\pathcode{e2b5e317b750e88437d30e95eb6a89ac88cb11c2},
\end{quote}
verified against source-manifest SHA-256
\begin{quote}
{\ttfamily\scriptsize 797858b93cf5b7c1e98581e00bd3fa58b0db28d1de3c87e4b23b3b5a4775b968}.
\end{quote}
Replacement job \pathcode{1835588} ran on \pathcode{h200-046-019} and
\pathcode{h200-171-121} from 2026-09-04 02:49:35 through 03:07:27 UTC and failed
\pathcode{15:0}.  The environment-order correction worked: all three loader phases and the
three-update \pathcode{minimal_a1} arm completed on every rank.  Its 16 training records have
finite, rank-equal pre-clip norms at all three boundaries and use the same input projection as the
restored sidecar.  The subsequent cache-inventory step failed before the next arm because the
training path, unlike each loader probe, had not explicitly stopped its four persistent DataLoader
workers per rank.  A local reproduction showed that live workers create
\pathcode{TMPDIR/pymp-*/listener-*} AF\_UNIX sockets; the strict cache inventory correctly rejects
such a special entry.  This is post-measurement harness cleanup evidence, not a model failure, but
the job has no final summary, artifact inventory, completion marker, or admissible classification.
The retained artifact root is
\begin{quote}
\pathcode{/lustre-storage/checkpoints/ngocbh/hdn/diagnostics/rdiag_shared_dt_v1/}\\
\pathcode{diagnostic_runs/gen0_repro_e2b5e317b750_1835588/artifacts}.
\end{quote}
Its scheduler log, stderr, invocation, and rank-0 \pathcode{minimal_a1} record SHA-256 values are
respectively
\begin{quote}
{\ttfamily\scriptsize a064b744df142a24957ae3ad3915da441d127d7c37c84c2b6ed91018024f2fe3},\\
{\ttfamily\scriptsize f9f48e95fb92bebe27c4f0ddb5ec44f27dc560854f5c28da92ca3c696d812424},\\
{\ttfamily\scriptsize 356a404e648d383301ac25b79bd6680c9ce289dd49877538fc869bf7a62e9374},\\
{\ttfamily\scriptsize b358564f0dfb6f8a39f077258a620119d5030743ce49e4c2c10118a77cf2cf4c}.
\end{quote}

The third correction factors the already-used worker shutdown into one helper and invokes it from
both loader probes and training-arm cleanup.  Special cache entries remain fail-closed, and future
errors now include their relative path, decoded kind, and mode.  The reviewed, pushed identity is
\begin{quote}
commit \pathcode{0e9cda6875e8709fa250110dcb35772d46d22817},\\
tree \pathcode{dad58c03e607bf38304dd41ae7d6e8993d729720}.
\end{quote}
The helper and host-test SHA-256 values are respectively
\begin{quote}
{\ttfamily\scriptsize a37fdea21c59757246c6fabff0dac15d8ca93ac3cfdc784d0549413544e30f70},\\
{\ttfamily\scriptsize 92766a1a21c35be8a26314cf0c4e85b8d298eec88abc37467d6a5e69d632ffa4}.
\end{quote}
The combined focused suite passes 100 tests, including an AF\_UNIX-socket rejection test, and an
independent review confirmed that shutdown removes all four reproduced listener sockets without
changing measured training math.  Its fresh immutable snapshot is
\begin{quote}
\pathcode{/lustre-storage/checkpoints/ngocbh/hdn/diagnostics/rdiag_shared_dt_v1/}\\
\pathcode{source_snapshots/rdiag_shared_dt_gen0_repro_diag_}\\
\pathcode{0e9cda6875e8709fa250110dcb35772d46d22817},
\end{quote}
verified against source-manifest SHA-256
\begin{quote}
{\ttfamily\scriptsize b228c5e0b44a9888d96d2d381c53940d7f6eb7264c71367ecde06f26535cf0e6}.
\end{quote}
Replacement job \pathcode{1835824} ran directly from that snapshot on
\pathcode{h200-046-019} and \pathcode{h200-094-010} from 2026-09-04 03:42:39 through 04:13:57
UTC.  The job and every step \pathcode{.0}--\pathcode{.18} completed \pathcode{0:0}.  It produced
the exact preregistered closure: 112 rank JSON records across seven 16-rank phases, 214 mode-0400
files recorded by \pathcode{artifact-inventory.json}, and 215 final directory entries because the
inventory does not list itself.  Independent validation returned \pathcode{VALID}; the log has
exactly one non-authorizing completion marker and the scheduler stderr is empty.  The retained
artifact root is
\begin{quote}
\pathcode{/lustre-storage/checkpoints/ngocbh/hdn/diagnostics/rdiag_shared_dt_v1/}\
\pathcode{diagnostic_runs/gen0_repro_0e9cda6875e8_1835824/artifacts}.
\end{quote}
The summary, artifact inventory, invocation, scheduler log, and empty scheduler stderr SHA-256
values are respectively
\begin{quote}
{\ttfamily\scriptsize aa466253b9bcdfeebf695ead220357b6693189c93aa5d089c573a3e887a71df6},\\
{\ttfamily\scriptsize 7aef2fc9349c84876b0ea7d46ee114147a4dc52be2cd4b0c776c4d6245b08339},\\
{\ttfamily\scriptsize eba71feda5c6f650968a143f49c4881fa524ce22ac4dec31b713ca1e0bbc9534},\\
{\ttfamily\scriptsize 8567a91cdacd6df770a49cc61f13ffe5be0431e479f30bbeae21afb50bdc6f40},\\
{\ttfamily\scriptsize e3b0c44298fc1c149afbf4c8996fb92427ae41e4649b934ca495991b7852b855}.
\end{quote}

The result is valid diagnostic evidence but fails the scientific hard gate.  All loader and
training input matrices match the original sidecars; the reconstructed model, optimizer, loader,
first-RNG anchor, rank topology, and corresponding-rank mapping also match.  Nevertheless both
the minimal A1/A2 pair and instrumented B1/B2 pair classify as
\pathcode{ORDINARY_RUNTIME_NONDETERMINISM}, making the overall training classification the same.
The first classifier-stage difference is the boundary-1 local-gradient Merkle root after the first
two microbatches.  The earliest chronological scalar difference occurs even earlier at rank 1,
microbatch 1: minimal A1 records loss token \pathcode{0x1.50da340000000p+3}, whereas minimal A2
records \pathcode{0x1.50ddc80000000p+3}, despite joint input SHA-256
\pathcode{5dcab61b5c06aa429639e90ef660cba3caf85ed6161fd6cf36812696f369cf7a}
and pre-forward RNG SHA-256
\pathcode{7db887f05ff2b0af83e6e7c4b4c9d7ac1a71e2f172d903ea409ce54cf21bdabe}
being identical.  Because both minimal replicates diverge before clipping, this is neither a loader
restore failure nor instrumentation-only or pre-clip-norm-collective divergence.  The exact
operator remains unidentified; trigger capture stays blocked pending a focused runtime/kernel
repeatability diagnostic and a successful repeat of this hard gate.

That focused forward-determinism diagnostic was implemented in the isolated worktree
\begin{quote}
\pathcode{/lustre-storage/checkpoints/ngocbh/hdn/diagnostics/rdiag_shared_dt_v1/}\\
\pathcode{source_worktrees/rdiag_shared_dt_forward_determinism_20260904}
\end{quote}
on branch
\pathcode{codex/rdiag-shared-dt-forward-determinism-20260904}, based on commit
\begin{quote}
\pathcode{0e9cda6875e8709fa250110dcb35772d46d22817}.
\end{quote}
Its reviewed executable plan is
\pathcode{docs/plans/}\\
\pathcode{2026-09-04-rdiag-shared-dt-forward-determinism-plan.md}; its entry points are
\begin{itemize}[leftmargin=1.8em]
  \item \pathcode{scripts/analyses/reproduce_rdiag_shared_dt_forward_determinism.py};
  \item \pathcode{scripts/analyses/run_rdiag_shared_dt_forward_determinism.sbatch}; and
  \item \pathcode{scripts/tests/rdiag/test_rdiag_shared_dt_forward_determinism_host.py}.
\end{itemize}
The design compares two ordinary-autotune and two diagnostic-only RMSNorm-warp-8 phases.  Each
process retains one minimally observed cold forward through the normal Fabric/FSDP and health
path, stops after the exact ordinary cross-entropy reduction and before backward, then performs
three equal-RNG warm traces plus a post-trace replay of all six existing RMSNorm forward
configurations.  Cold and warm evidence are reported on separate axes; every configuration replay
must be stable across corresponding ranks before RMSNorm selection can be called strongly
supported.  This diagnostic cannot authorize trigger capture: any candidate execution control
must still pass a fresh copy of the unchanged generation-zero \pathcode{minimal=MATCHED} gate.
Independent plan and runtime reviews found and closed three pre-launch defects: an invalid
cross-scope FSDP weight-object identity requirement, conflation of the raw BF16 \pathcode{lm_head}
output with Fabric's distinct FP32 return/loss input, and a missing final fail-together exchange
after cleanup and rank-local result validation.  The focused suite now passes 57 tests and the
combined forward-diagnostic, generation-zero, and soak-790 host suites pass 157 tests.  The
reviewed and pushed diagnostic identity is
\begin{quote}
commit \pathcode{70165b017cd28df832a9a8797a72a1f8234284d0},\\
tree \pathcode{c58e220299a54525c443f4f0dc204e22e1fb19b8}.
\end{quote}
The plan, helper, runner, and host-test SHA-256 values are respectively
\begin{quote}
{\ttfamily\scriptsize 7aff0bea0a85887e7744f57816239dd598426c5bc4623ffb140a01f410a98f26},\\
{\ttfamily\scriptsize 52583f1d8cfdfdd571897a446ce00a69b598a72889e244be9e342498fd32010f},\\
{\ttfamily\scriptsize eeb1704bdb8bb30a376ba5483747e1ad6022426f96fe6cab3e8627cf642fe0fd},\\
{\ttfamily\scriptsize 71a20bbcfce4c1109766263bca07e5e8099312f7ec6196d4b2e276aba62a3a65}.
\end{quote}
The 4,242-file immutable source snapshot is
\begin{quote}
\pathcode{/lustre-storage/checkpoints/ngocbh/hdn/diagnostics/rdiag_shared_dt_v1/}\\
\pathcode{source_snapshots/rdiag_shared_dt_forward_determinism_}\\
\pathcode{70165b017cd28df832a9a8797a72a1f8234284d0},
\end{quote}
verified against source-manifest SHA-256
\begin{quote}
{\ttfamily\scriptsize e1129c0caa1832c5bc5ff3d2f978a8c6a18a38a1c55f7d0027b1792ddb77e665}.
\end{quote}
H200 diagnostic job \pathcode{1836996} ran from this snapshot on
\pathcode{h200-079-003} and \pathcode{h200-079-202} from 2026-09-04 06:17:47 through 06:37:39
UTC.  The allocation and every step \pathcode{.0}--\pathcode{.12} completed \pathcode{0:0}.
It produced the exact preregistered closure: 64 rank JSON records, 64 NCCL logs, eight node/phase
cache inventories, 154 mode-0400 files recorded by \pathcode{artifact-inventory.json}, and 155
final directory entries because the inventory does not list itself.  Independent validation
returned \pathcode{VALID}; the log has exactly one non-authorizing completion marker and scheduler
stderr is empty.  The retained artifact root is
\begin{quote}
\pathcode{/lustre-storage/checkpoints/ngocbh/hdn/diagnostics/rdiag_shared_dt_v1/}\\
\pathcode{diagnostic_runs/forward_determinism_70165b017cd2_1836996/artifacts}.
\end{quote}
The summary, artifact inventory, invocation, scheduler log, and empty scheduler stderr SHA-256
values are respectively
\begin{quote}
{\ttfamily\scriptsize 741e75362c6214b4a82644bbab2b12b61fc6ed827e80ff4ac4d3f86be94861ff},\\
{\ttfamily\scriptsize 463b68fc3d70f85682fe03137df857e5e577d9c9aae009e7f7c13659f591931f},\\
{\ttfamily\scriptsize 0179440c6f58ac1935626684916b4599bd2035820cde6ee4ccf1cc5ebcdbfde5},\\
{\ttfamily\scriptsize 45a7018391bd5ec572d48c2557d3a0239cadd679289cf465f8e241a444cfb29a},\\
{\ttfamily\scriptsize e3b0c44298fc1c149afbf4c8996fb92427ae41e4649b934ca495991b7852b855}.
\end{quote}
The pre/post generation-zero input inventories are byte-identical, both with SHA-256
\pathcode{b5bc763929a966c54adbfeb1d9a1d9d1c0ae891556860ad89449a333fec50cbe}.

The preregistered classification is
\pathcode{RMSNORM_AUTOTUNE_SELECTION_STRONGLY_SUPPORTED}.  Every initial-state and input
comparability axis matches, every within-process warm repeat matches its cold result, and both
warp-8 phases match exactly across processes.  By contrast, the two ordinary cold processes
diverge on ranks 3, 6, 12, and 14.  Those are exactly the ranks whose selected forward RMSNorm
warp count changes: respectively $8\mathbin{\to}4$, $2\mathbin{\to}4$,
$4\mathbin{\to}8$, and $4\mathbin{\to}1$.  Their first warm-trace difference is
\pathcode{block_00.norm_1}; their external logits, unreduced loss, and reduced loss then differ,
while z-loss remains equal.  Replaying each selected configuration reproduces its recorded warm
output, and all six forced configurations are stable when compared at equal input, weight, rank,
and configuration.  Thus timing-dependent RMSNorm forward autotune selection is strongly
supported as the source of the observed first-forward variance.  This remains diagnostic rather
than production authorization: it does not test backward repeatability, and a narrowly scoped
fixed-warp production control must pass a fresh unchanged generation-zero hard gate.

The corrective implementation is now isolated in worktree
\begin{quote}
\pathcode{/lustre-storage/checkpoints/ngocbh/hdn/diagnostics/rdiag_shared_dt_v1/}\\
\pathcode{source_worktrees/rdiag_shared_dt_rmsnorm_control_20260904}
\end{quote}
on branch \pathcode{codex/rdiag-shared-dt-rmsnorm-control-20260904}, based exactly on diagnostic
commit \pathcode{70165b017cd28df832a9a8797a72a1f8234284d0}.  Its executable plan is
\pathcode{docs/plans/}\\
\pathcode{2026-09-04-rdiag-shared-dt-rmsnorm-determinism-plan.md}.  The selected scope is the
narrow forward-only control justified by job \pathcode{1836996}: for the four resolved
shared-DT configurations only, require the cold RMSNorm forward autotuner's exact existing
six-candidate closure and install its unique existing warp-8 object as a singleton.  The backward
RMSNorm tuner, every FLA tuner, precision policy, recurrence, parameters, and generation-zero
helper/runner remain unchanged.  Activation keys on the resolved internal config identity rather
than the raw CLI name, because the generation-zero diagnostic deliberately uses a temporary CLI
alias whose copied payload retains the production internal name.  No environment selector,
config/recipe field, or runtime-fingerprint field is added: the historical runtime, data, model,
and sidecar identities remain fixed, while the new source-manifest SHA binds the control.  The
implementation must fail closed on warm, altered, or ambiguous tuner state, emit a canonical
policy digest with rank consensus before model construction, cover the direct evaluation and
gradient-gate constructors that bypass \pathcode{pretrain.py}, and prove that the backward tuner
is not mutated.  The process-global pin is intentionally irreversible; every authenticated
entry point runs one shared-DT model family per process, and a mixed-family caller must use a
separate process rather than attempting to restore Triton global state.

The initial policy implementation and smoke harness were reviewed, committed, and pushed as
\begin{quote}
commit \pathcode{28ea1c51a16825ae02d3d4c6520087e8b9b2ca42},\\
tree \pathcode{1f7dd60edbebe92b429e4a541fdea85e6ec3c182}.
\end{quote}
The plan, policy module, trainer integration, authenticated evaluation integration, and gradient
gate integration SHA-256 values are respectively
\begin{quote}
{\ttfamily\scriptsize 91cc8fd7580ff953f71cedd5594a489eacc126168bce68904d7b17ffebbcbbb2},\\
{\ttfamily\scriptsize 74898f2dfa3c3073e491616eb7e847acdaa695c4b8b70596af5534e5c6591433},\\
{\ttfamily\scriptsize ff9fa812d81266577c5d9b4404966e039e1bad205abf19009099ffc24651a584},\\
{\ttfamily\scriptsize b78e03c7996933e5c68a2e08a46ab074a8a235b5ac783e2588a523e14497c5f8},\\
{\ttfamily\scriptsize 755fa87cc6598b21391a06de20a9b25ce1a9286b383e8e714a26764c141015e6}.
\end{quote}
The corresponding implementation entry points are
\begin{itemize}[leftmargin=1.8em]
  \item \pathcode{lit_gpt/rdiag_shared_dt_rmsnorm.py} for applicability checks, fail-closed tuner
        inspection, installation, and the canonical policy record;
  \item \pathcode{pretrain.py}, \pathcode{scripts/eval/eval_lm_harness.py}, and
        \pathcode{scripts/tests/rdiag/rdiag_shared_dt_gradient_ab_run.py} for authenticated
        training, evaluation, and gradient-gate activation; and
  \item \pathcode{scripts/tests/rdiag/test_rdiag_shared_dt_rmsnorm_policy_host.py},
        \pathcode{scripts/tests/rdiag/test_rdiag_shared_dt_rmsnorm_integration_host.py}, and the
        policy smoke files below for host and H200 contracts.
\end{itemize}
The smoke helper and runner are
\pathcode{scripts/tests/kernels/test_rdiag_shared_dt_rmsnorm_policy.py} and
\pathcode{scripts/tests/kernels/run_rdiag_shared_dt_rmsnorm_policy.sbatch}, with SHA-256 values
\pathcode{aa58e4cfa0efb2f35e4cf4c38d86eb1492a1fbba359b644c7cf526c8c15e13cc} and
\pathcode{85f8ce9bf11cf61d006916ac44c3955fbe1f7b1c070ac17453f4afd89f850c84}.
It exercises the actual wrapped Triton forward at $M=16384$ and widths 1024, 2048, and 2304,
requires repeat and two-rank output/rstd byte equality, and benchmarks all six existing candidates
with 40 order-balanced samples each.  Warp 8 must remain within $1.25\times$ the empirical fastest
candidate by both median and lower-quartile latency.  Installation leaves the backward candidate
objects unchanged; a later backward smoke only proves that the untouched path still executes.
At the initial implementation revision, the focused policy/integration/smoke contract suite passed
49 tests.  The combined existing generation-zero, forward-diagnostic, runtime-identity,
gate-runner, and authorization suite passed 309 tests with three declared skips when pytest
temporary files were redirected to Lustre; an initial attempt exhausted the already-full system
temporary filesystem and is not a code failure.  After the later smoke-helper and launcher-only
revisions, the final focused runner/smoke-host suite passed nine tests.
Independent final review reports no Critical, Important, or Minor finding.  The generation-zero
helper and runner remain byte-identical to commit \pathcode{70165b0}, with SHA-256 values
\pathcode{a37fdea21c59757246c6fabff0dac15d8ca93ac3cfdc784d0549413544e30f70} and
\pathcode{f9d5c8a04f7384b360eb04435c32d0a2abeb177e6f6534e91f4910e5b7c1bb14}.

The 4,249-file immutable control snapshot is
\begin{quote}
\pathcode{/lustre-storage/checkpoints/ngocbh/hdn/diagnostics/rdiag_shared_dt_v1/}\\
\pathcode{source_snapshots/rdiag_shared_dt_rmsnorm_control_}\\
\pathcode{28ea1c51a16825ae02d3d4c6520087e8b9b2ca42},
\end{quote}
independently verified against source-manifest SHA-256
\begin{quote}
{\ttfamily\scriptsize 4207beb5e9eb0aeaaa11c7eb489d2295122ac257ecedfc77b5dfde21f35241c1}.
\end{quote}
Non-authorizing two-H200 smoke job \pathcode{1837294} was submitted from exactly this snapshot and
failed before Python after two seconds.  The batch shell and its \pathcode{srun} task steps receive
different node-local \pathcode{/tmp} namespaces on this cluster; the batch-created cache parent was
therefore absent in both GPU tasks.  No numerical test executed and no result JSON was published;
the empty artifact directory is
\begin{quote}
\pathcode{/lustre-storage/checkpoints/ngocbh/hdn/diagnostics/rdiag_shared_dt_v1/}\\
\pathcode{diagnostic_runs/rmsnorm_policy_28ea1c51_ZxLwOVLa/artifacts}.
\end{quote}
The runner correction creates and removes a distinct private \pathcode{/tmp} root inside each
\pathcode{srun} task, preserves the workload exit status through guarded cleanup, and has passed
its seven host contracts, shell syntax check, and independent review.  Because this changes the
runner bytes, the reviewed and pushed replacement identity is
\begin{quote}
commit \pathcode{edf5bb3aa95e2d64beabc8f4196c56029851b037},\\
tree \pathcode{fbcc46493836860a268d9fe56514ca8eadd38993}.
\end{quote}
The corrected plan, runner, and runner-host-test SHA-256 values are respectively
\begin{quote}
{\ttfamily\scriptsize cb77a67a4e06e5d31c93547e4560064eb10bd045261a37a8bc5b971e7338fdb7},\\
{\ttfamily\scriptsize 48463cff46aa3eae19a46514655dfac90b8886e02a8df08036a9df4ad12c66ba},\\
{\ttfamily\scriptsize 2b06e25a26341d7eb108f06c6abe7c087a7ea035d9ede29715eb36f801814b75}.
\end{quote}
Its independently verified 4,249-file snapshot is
\begin{quote}
\pathcode{/lustre-storage/checkpoints/ngocbh/hdn/diagnostics/rdiag_shared_dt_v1/}\\
\pathcode{source_snapshots/rdiag_shared_dt_rmsnorm_control_}\\
\pathcode{edf5bb3aa95e2d64beabc8f4196c56029851b037},
\end{quote}
with source-manifest SHA-256
\begin{quote}
{\ttfamily\scriptsize 2dc7fc4d3856c43f607e00df606a63c4c221fcbd5c5da5f461e6fab34fc59c98}.
\end{quote}
Replacement smoke job \pathcode{1837353} reached Python from this corrected snapshot but failed
on the first NCCL collective: task-local one-GPU visibility caused NCCL 2.27 to report
\pathcode{invalid device ordinal}.  The irrecoverable allocation was cancelled after 2:08 to
release its H200s.  It also published no result; its empty artifact target is
\begin{quote}
\pathcode{/lustre-storage/checkpoints/ngocbh/hdn/diagnostics/rdiag_shared_dt_v1/}\\
\pathcode{diagnostic_runs/rmsnorm_policy_edf5bb3a_MBXIibLe/artifacts}.
\end{quote}
The second launcher correction follows a completed project H200 reference: both allocated GPUs
are visible to both tasks via \pathcode{--gpus-per-node=2 --gpu-bind=none}, and ranks select CUDA
devices by \pathcode{LOCAL_RANK}.  It also prints rank-local phase tracebacks before recovery and
re-raises fatal CUDA/NCCL failures instead of attempting another collective on a poisoned CUDA
context.  The reviewed and pushed identity is now
\begin{quote}
commit \pathcode{fd9ca79e697b2fce14b8938459a7d2e63e434bea},\\
tree \pathcode{b72796b9106c8e06d2b7e59051e5cab094ec8b40}.
\end{quote}
The current plan, smoke helper, runner, and runner-host-test SHA-256 values are respectively
\begin{quote}
{\ttfamily\scriptsize 61277a3f3db22a8495f5cc2cfbdd50fd038cc53664c219c04bcacad8e8f3c572},\\
{\ttfamily\scriptsize af2f17b7b178930291a1225345373713c3fc1357165a3812afa9f237c494d31e},\\
{\ttfamily\scriptsize 87652ee3fdb1fd54cf17cfdbcb32e65877aa1fa27b8c146af96b860e0fe9d38f},\\
{\ttfamily\scriptsize cfb736f894504e8bc2f941074452579dbb87c825ce5fd7cb54c36e8674821d69}.
\end{quote}
Its independently verified 4,249-file snapshot is
\begin{quote}
\pathcode{/lustre-storage/checkpoints/ngocbh/hdn/diagnostics/rdiag_shared_dt_v1/}\\
\pathcode{source_snapshots/rdiag_shared_dt_rmsnorm_control_}\\
\pathcode{fd9ca79e697b2fce14b8938459a7d2e63e434bea},
\end{quote}
with source-manifest SHA-256
\begin{quote}
{\ttfamily\scriptsize 0d49e9471d3e6a3b48bc21838b6366cf26d7b8f9418639b70a5f012dec866f20}.
\end{quote}
Smoke job \pathcode{1837393} ran from this snapshot on \pathcode{h200-129-161} from
2026-09-04 07:51:39 through 07:52:08 UTC.  The allocation, batch, extern, and sole
\pathcode{srun} step all completed \pathcode{0:0}.  It published exactly one retained artifact,
\begin{quote}
\pathcode{/lustre-storage/checkpoints/ngocbh/hdn/diagnostics/rdiag_shared_dt_v1/}\\
\pathcode{diagnostic_runs/rmsnorm_policy_fd9ca79e_gz8Nu1Ly/artifacts/result.json}.
\end{quote}
That mode-0400, single-link, root-owned result is 17,835 bytes and has SHA-256
\begin{quote}
{\ttfamily\scriptsize f39137e97c07d7cafc9d930b4060dce798b352a439aba9a22201d106107aeb30}.
\end{quote}
The scheduler log SHA-256 is
\begin{quote}
{\ttfamily\scriptsize a6228b2518a21a27e627c69c2f0b61d96613396230833eaa7acd2c757b288490}.
\end{quote}
Independent validation returned \pathcode{VALID}.  The result is explicitly
\pathcode{authorizing=false} and \pathcode{production_gate=false}.  Both ranks agreed on policy
record SHA-256
\begin{quote}
{\ttfamily\scriptsize cd31a0f1bc60aa5421cb7074f5236a14ff4fb66d59632e587e0920f4e343ecac}.
\end{quote}
At widths 1024, 2048, and 2304, the first forward used the exact pre-existing warp-8 object, the
forward cache remained empty, repeat outputs and reciprocal standard deviations were byte-exact,
and the corresponding cross-rank values were byte-exact.  Across all rank/width pairs, the maximum
warp-8-to-fastest ratio was 1.04614 by median and 1.04703 by lower-quartile latency, well within the
preregistered $1.25\times$ bounds.

The deliberately untouched backward smoke also exposed the next concrete risk.  Rank 0 selected
warp 16 while rank 1 selected warp 8; equal forward outputs and weight-gradient hashes accompanied
different input-gradient hashes:
\begin{quote}
{\ttfamily\scriptsize f364b0deeac41662c4571b76f4ae57c9d45ac5923b47b0d0685f357ca7f83647},\\
{\ttfamily\scriptsize ae9fca3f6523a4a59d0670dd4af9bdbf041c6992f5df43eb37907baa1a9be3d3}.
\end{quote}
Backward cross-rank equality was not a preregistered requirement of this forward-control smoke, so
the bounded smoke validly passes; the observation nevertheless makes RMSNorm backward autotuning
the leading suspect if the forward-only control is insufficient.  The unchanged replacement
generation-zero run discriminates that sufficiency, not backward causality.  That gate, not
installation or this smoke, decides whether trigger capture can begin.  If it does not return
\pathcode{minimal=MATCHED}, focused backward-RMSNorm localization is the next priority, and a
separately reviewed narrow control is warranted only if that localization implicates it.  No
trigger or production authorization is presently claimed.

The unchanged replacement generation-zero job \pathcode{1837446} then ran from the same
\pathcode{fd9ca79e...} snapshot on \pathcode{h200-004-175} and \pathcode{h200-071-008} from
2026-09-04 07:57:40 through 08:27:18 UTC.  The allocation, batch, extern, and all steps
\pathcode{.0}--\pathcode{.18} completed \pathcode{0:0}.  Its retained artifact root is
\begin{quote}
\pathcode{/lustre-storage/checkpoints/ngocbh/hdn/diagnostics/rdiag_shared_dt_v1/}\\
\pathcode{diagnostic_runs/gen0_repro_fd9ca79e697b_1837446/artifacts}.
\end{quote}
Independent validation returned \pathcode{VALID}: the directory contains the exact 214-file
inventoried closure plus \pathcode{artifact-inventory.json}, for 215 mode-0400, single-link files.
The summary, inventory, invocation, scheduler log, and empty scheduler stderr SHA-256 values are
respectively
\begin{quote}
{\ttfamily\scriptsize c69adb84b265f7c20560ad1be54a7654fcda963ffe068cda252f398c942d9c54},\\
{\ttfamily\scriptsize 97a46cae0bc31162351ffa1d39acf6f55b048d4f8ccc22677b549082a3c64a14},\\
{\ttfamily\scriptsize 2bf614563fa4f82662220809cfaf2c1b9d23f512e11400e87fa11db40644f179},\\
{\ttfamily\scriptsize fb0cbcfc30f950f5b899095bb603caf1110cb369eb25eff6bdba682d396b1e90},\\
{\ttfamily\scriptsize e3b0c44298fc1c149afbf4c8996fb92427ae41e4649b934ca495991b7852b855}.
\end{quote}
The log has exactly one completion marker.  Its pre/post checkpoint-input inventories are
byte-identical, both with SHA-256
\begin{quote}
{\ttfamily\scriptsize 20bbc0c2a345eaa46671d670c0d865b6f9a73266c9310bd81a1f43243ed120fd}.
\end{quote}
Both node runtime-identity records retain SHA-256
\begin{quote}
{\ttfamily\scriptsize 135ea68e47d1e9728f1295e480dd22518566f818619d06b6853bed343b00639d}.
\end{quote}

The hard gate passes.  The reconstructed common anchor, rank topology, training inputs, and
sidecar inputs match; all three loader arms match over their first six rank-local microbatches.
Both preregistered replicate-pair classifications are \pathcode{MATCHED}, including the required
\pathcode{minimal=MATCHED}; global loss and pre-clip norm are exact at each of the first three
update boundaries.  The overall training classification is
\pathcode{MATCHED_RECONSTRUCTED_FIRST_THREE_UPDATES}.  This remains non-authorizing and does not
make the reconstructed process RNG historically comparable to the failed production run, whose
RNG was not retained.  It does establish that the forward-only control is sufficient for the
generation-zero repeatability gate.  The backward-smoke observation remains a risk note rather
than evidence warranting a second intervention, so no backward pin is introduced.  This result
released revision, review, and non-authorizing execution of the trigger-capture diagnostic; it
does not authorize production.

The gated follow-on capture work is isolated in worktree
\begin{quote}
\pathcode{/lustre-storage/checkpoints/ngocbh/hdn/diagnostics/rdiag_shared_dt_v1/}\\
\pathcode{source_worktrees/rdiag_shared_dt_trigger_capture_20260904}
\end{quote}
on branch \pathcode{codex/rdiag-shared-dt-trigger-capture-20260904}; its plan is
\pathcode{docs/plans/}\\
\pathcode{2026-09-04-rdiag-shared-dt-trigger-capture-plan.md}.  Its earlier plan-only identity was
\begin{quote}
commit \pathcode{8074daf35bbd3dac839f0fc2601361d7cf069d50},\\
tree \pathcode{b27a3f6ab3269df09daf02f08507f178a77cefe9}.
\end{quote}
The accepted control lineage was merged without rewriting either parent as
\begin{quote}
commit \pathcode{57afdd1f9fe497f3630f262f7febfcbbd56e4f7e},\\
tree \pathcode{ad12554140dac65bb315dd49a0c7fc003393acb1},
\end{quote}
whose parents are \pathcode{8074daf...} and \pathcode{fd9ca79e...}.  The plan was then revised,
independently reviewed for evidence integrity and runtime semantics, committed, and pushed as
\begin{quote}
commit \pathcode{dfd928cf7e861f5f828f6486730cf2836cdd3089},\\
tree \pathcode{1666a25b3ef188dca6c151a18209c11d57b470f9}.
\end{quote}
Its reviewed plan SHA-256 is
\begin{quote}
{\ttfamily\scriptsize 48f65aaf84d0c0ce403ed25bc9bd860c33b511c2119c28191b3e71123473e7d8}.
\end{quote}
The plan pins job \pathcode{1837446}'s summary, inventory, invocation, source, input, runtime, and
log identities at their actual nested classifier paths.  Its exact 1,281-byte root scheduler row
is retained separately, without modifying the closed generation-zero directory, at
\begin{quote}
\pathcode{/lustre-storage/checkpoints/ngocbh/hdn/diagnostics/rdiag_shared_dt_v1/}\\
\pathcode{diagnostic_runs/gen0_repro_fd9ca79e697b_1837446_scheduler_record/}\\
\pathcode{sacct-root-row.txt},
\end{quote}
with mode 0400, one link, and SHA-256
\begin{quote}
{\ttfamily\scriptsize 7b3df5231f751e1841c448cb55714c471522bcd74bbb973b577477785f07f587}.
\end{quote}
Review also closes process-RNG changes during loader read-ahead, rank-asymmetric capture failures,
unexpected loader EOF/worker errors, nonfinite-forward handling, live RMSNorm autotuner provenance,
and the two earlier launcher topology/cache failures.  The executable scope remains the plan plus
\pathcode{scripts/analyses/capture_rdiag_shared_dt_trigger.py},
\pathcode{scripts/analyses/run_rdiag_shared_dt_trigger_capture.sbatch}, and
\pathcode{scripts/tests/rdiag/test_rdiag_shared_dt_trigger_capture_host.py}.  The implementation
was independently reviewed, committed, and pushed as
\begin{quote}
commit \pathcode{ef57031fc2f342a0dfe9c9ef7af2d0e268e332cb},\\
tree \pathcode{db686e18708233ecbb049cb68cbc1ace9d81d4d6}.
\end{quote}
The initial helper, runner, and host-test SHA-256 values were respectively
\begin{quote}
{\ttfamily\scriptsize ec0db20eb6f0609425c5add313dab82c111b7619e8c2c0f94840ed45cc0e29d3},\\
{\ttfamily\scriptsize 11c92b07965f69e622fe7ec0af3ae00256942eaeb7c9192adda7d15cf6cde95e},\\
{\ttfamily\scriptsize 56eb9d0a04d053a0e1e19724155e56dc3f1d225ee150c4d18bbbf8c19127f8a8}.
\end{quote}
The initial focused suite passed 60 tests, the adjacent RDiag suite passed 170 tests, and the broad
compatibility sweep passed 495 tests with three skips.  Its immutable 4,253-file Git-object
snapshot, excluding only the paper gitlink, is
\begin{quote}
\pathcode{/lustre-storage/checkpoints/ngocbh/hdn/diagnostics/rdiag_shared_dt_v1/}\\
\pathcode{source_snapshots/rdiag_shared_dt_trigger_capture_}\\
\pathcode{ef57031fc2f342a0dfe9c9ef7af2d0e268e332cb},
\end{quote}
with source-manifest SHA-256
\begin{quote}
{\ttfamily\scriptsize 5d3edaa28885ee012fde8887f96af827f0fe320b84eba4ea5e0b5ff46cccb907}.
\end{quote}
Both the publisher's closed-tree verification and a fresh private-copy reconstruction from Git
objects passed.  The filesystem's unavailable no-replace rename path required the publisher's
fail-closed copy fallback; the final execution root is owner-held mode 0555.  Hidden staging
siblings retained by publication and by a concurrent collision check are evidence only and are
explicitly rejected as execution roots.

Non-authorizing two-node, 16-H200 trigger-capture job \pathcode{1838085} was submitted from that
snapshot at 2026-09-04 10:11:12 UTC with the exact failed-run data manifest.  It ran on
\pathcode{h200-205-088} and \pathcode{h200-231-137} from 10:30:02 through 10:33:51 UTC, then failed
closed before capture training.  Source-copy, topology, runtime-identity, and two-node exit-cleanup
steps all completed \pathcode{0:0}; the root and batch exited \pathcode{1:0}.  Its run root is
\begin{quote}
\pathcode{/lustre-storage/checkpoints/ngocbh/hdn/diagnostics/rdiag_shared_dt_v1/}\\
\pathcode{diagnostic_runs/trigger_capture_ef57031fc2f3_1838085},
\end{quote}
where \pathcode{anchor} and \pathcode{output} are empty and \pathcode{artifacts} contains only six
preflight files.  There are no capture results, NCCL logs, summary, inventory, checkpoint, or
completion marker.  The external root and full-step scheduler rows have SHA-256 values
\begin{quote}
{\ttfamily\scriptsize 2d9eae9ab2fc00dd7e407629b9214d97efb4629fb1753f5bbac9b5f0ebf9380f},\\
{\ttfamily\scriptsize 395f5ebda2a1ad3910b956209ead563a3de974fd80f9a72a1a27b6fec76105db},
\end{quote}
under \pathcode{diagnostic_runs/trigger_capture_ef57031fc2f3_1838085_scheduler_record}.  Scheduler
record filenames are \pathcode{root-row.txt} and \pathcode{all-steps.txt}.  Scheduler stdout and
stderr are
\begin{quote}
\pathcode{/lustre-storage/checkpoints/ngocbh/hdn/ckpts/logs/}\\
\pathcode{rdiag_sdt_trigger_1838085.log}, and\\
\pathcode{/lustre-storage/checkpoints/ngocbh/hdn/ckpts/logs/}\\
\pathcode{rdiag_sdt_trigger_1838085.err};
\end{quote}
their SHA-256 values are respectively
\begin{quote}
{\ttfamily\scriptsize 158b30dd9c16f898aed9b40df8272e059506c806384c8a75359ba919fe5a5447},\\
{\ttfamily\scriptsize 196a496ecaa4e47cba69b79bab5292c81d2006f65de9861d43c47701623cba87}.
\end{quote}

The defect was a strict framing mismatch rather than a scientific failure: the helper queried
\pathcode{sacct -P} (\pathcode{--parsable2}), whose 19 fields have 18 separators and no trailing
\pathcode{|}, but parsed it as \pathcode{-p} output with a trailing separator.  The retained
1,281-byte job-1837446 row and a fresh \pathcode{-P} query remain byte-identical with SHA-256
\begin{quote}
{\ttfamily\scriptsize 7b3df5231f751e1841c448cb55714c471522bcd74bbb973b577477785f07f587}.
\end{quote}
The minimal correction keeps \pathcode{-P}, requires exactly 19 fields, and uses a strict zip; it
does not alter any scientific path.  It was reviewed, committed, and
pushed as
\begin{quote}
commit \pathcode{4cbddbc1f715e5eb3a4177fe1557b29d7ab4f3d0},\\
tree \pathcode{869a5f036fe8cf2bcb085c8b4e8d38cb5e64f550}.
\end{quote}
The corrected helper, unchanged runner, and corrected host-test SHA-256 values are respectively
\begin{quote}
{\ttfamily\scriptsize 3d6879298bed98dfc6fa992d518d7d0ef0ae80267c4e8d5c7db3a134ba1abca6},\\
{\ttfamily\scriptsize 11c92b07965f69e622fe7ec0af3ae00256942eaeb7c9192adda7d15cf6cde95e},\\
{\ttfamily\scriptsize 8eacb854d8d3297c2e9f1eeb2a19a404f51521f900027e29af47f3ff3957602d}.
\end{quote}
The corrected focused suite passed 68 tests and the broad suite passed 503 tests with three skips;
the full real prerequisite command also passed against the pinned checkpoint, generation-zero
closure, data manifest, scheduler record, and initial snapshot.  The replacement immutable
snapshot is
\begin{quote}
\pathcode{/lustre-storage/checkpoints/ngocbh/hdn/diagnostics/rdiag_shared_dt_v1/}\\
\pathcode{source_snapshots/rdiag_shared_dt_trigger_capture_}\\
\pathcode{4cbddbc1f715e5eb3a4177fe1557b29d7ab4f3d0},
\end{quote}
with source-manifest SHA-256
\begin{quote}
{\ttfamily\scriptsize e27fb7c3a64465ddd3b4d92cbd9d3dc0e83d0e2f9d5a48e8b69540e7f9327b6d}.
\end{quote}
Both publication verification and a fresh private-copy reconstruction passed.  Replacement job
\pathcode{1838250} was submitted at 2026-09-04 10:48:12 UTC, started at 12:21:08 UTC on
\pathcode{h200-094-010} and \pathcode{h200-094-088}, and failed closed at 12:41:26 UTC.
The two-node source-copy, topology, and runtime-identity steps completed \pathcode{0:0}; exact
snapshot, generation-zero closure, data-manifest, and retained scheduler-evidence verification
all passed, so the corrected parsable2 boundary has been exercised successfully.  The intended
16-rank capture step started at 12:24:44 UTC, reconstructed the 621,284,352-parameter model,
resumed the exact generation-zero state, and trained normally with 60.4 GB peak allocation per
GPU.  Update 835 was the last printed completed update.  An eligible finite, rank-agreed trigger
with norm in $[3,10]$ subsequently became pending before update 840, but its exact boundary and
norm existed only in process memory and are not recoverable from this failed attempt.  During
anchor publication, all 16 ranks then failed together in \pathcode{anchor-rmsnorm-tuners} before
any anchor write: \pathcode{inspect_rmsnorm_tuners} passed the fully qualified display strings
\pathcode{lit_gpt.rmsnorm._layer_norm_fwd_1pass_kernel} and
\pathcode{lit_gpt.rmsnorm._layer_norm_bwd_kernel} directly to module \pathcode{getattr}, although
the module exposes their basename attributes.  Capture step \pathcode{.3} therefore exited
\pathcode{1:0} after 16:39; cache harvest/removal step \pathcode{.4} and private-source cleanup
step \pathcode{.5} completed \pathcode{0:0}.  This is an implementation failure after trigger
selection, not \pathcode{NO_TRIGGER_THROUGH_900} and not a scientific trigger result.  Its
deterministic run root is
\begin{quote}
\pathcode{/lustre-storage/checkpoints/ngocbh/hdn/diagnostics/rdiag_shared_dt_v1/}\\
\pathcode{diagnostic_runs/trigger_capture_4cbddbc1f715_1838250}.
\end{quote}
It contains 43 mode-0400 partial artifacts, an empty anchor root, and only the expected zero-byte
checkpoint placeholder under \pathcode{output}.  Rank capture results, \pathcode{capture-phase},
post-input inventory, invocation, summary, artifact inventory, loader verification, and completion
marker are all absent.  The scheduler stdout/stderr and capture stdout/stderr SHA-256 values are
respectively
\begin{quote}
{\ttfamily\scriptsize 2fae5175174f72c705af06fb0023a333785f447f9d83214f00e126d921a4aa1b},\\
{\ttfamily\scriptsize 028a4dacfc45125ba976c3990402c20a538e306929c352f6b8178d99351f5d8a},\\
{\ttfamily\scriptsize 1d39eebde75592b3a4fa202a37b0be5f122884f05cbfcba34b9a50b8c42d24c3},\\
{\ttfamily\scriptsize aa7dbceb04a817036fafc5192c1b32a7e921357e7348ce7ac2fce3c8f4567e38}.
\end{quote}
The externally retained root and nine-row full-step scheduler records are mode 0400, one-link files
under
\begin{quote}
\pathcode{diagnostic_runs/trigger_capture_4cbddbc1f715_1838250_scheduler_record},
\end{quote}
with SHA-256 values
\begin{quote}
{\ttfamily\scriptsize 3db33e5c49a94e7fcb7dfef72a19f9b4d6c67d0597a2108940e140eaf24a870e},\\
{\ttfamily\scriptsize 012202c4d97ce0ae5a7ccbcccc002e2bef202ca192cb182f8f0ea51bd1b96185}.
\end{quote}
The minimal repair resolves each qualified name's final component for Python lookup while retaining
the qualified string in canonical policy evidence.  A host regression now exercises the real
qualified-symbol/module-basename shape and proves inspection leaves both tuner states unchanged;
the focused suite passes 69 tests and the focused-plus-policy suite passes 91.  Independent review
found no blocking issue; Python compilation, shell syntax, and Git diff checks pass.  A deliberately
overbroad login-node sweep was not an acceptance test: its CUDA cases rejected the GPU-less host
and one unrelated source-publisher test observed the already documented Lustre staging fallback.
The reviewed and pushed repair is
\begin{quote}
commit \pathcode{015ba66e07f0ec1860b188ea74e638ce955384cc},\\
tree \pathcode{d1d59c7a08112e29ac942b0153f36552a7cc60ea}.
\end{quote}
The repaired helper, unchanged runner, and regression-test SHA-256 values are respectively
\begin{quote}
{\ttfamily\scriptsize c732721ae49d4798a2ff280300ffb09f7e4f83470ee415092e8955dd7f51ca5e},\\
{\ttfamily\scriptsize 11c92b07965f69e622fe7ec0af3ae00256942eaeb7c9192adda7d15cf6cde95e},\\
{\ttfamily\scriptsize 02ec5161fbb73dbe1db0a1ab197fc9cf908844fa04c1871a48ac8e9e1ad726ef}.
\end{quote}
Its immutable 4,253-file snapshot, again excluding only the paper gitlink, is
\begin{quote}
\pathcode{/lustre-storage/checkpoints/ngocbh/hdn/diagnostics/rdiag_shared_dt_v1/}\\
\pathcode{source_snapshots/rdiag_shared_dt_trigger_capture_}\\
\pathcode{015ba66e07f0ec1860b188ea74e638ce955384cc},
\end{quote}
with source-manifest SHA-256
\begin{quote}
{\ttfamily\scriptsize bf0b5288a873f5f7c67963492239226921150b2c0f7350de6a3cf57367987dca}.
\end{quote}
Both publication verification and an independent fresh private-copy reconstruction passed.
Replacement job \pathcode{1838750} was submitted from that snapshot at 2026-09-04 13:06:42 UTC
and ran on \pathcode{h200-094-010} and \pathcode{h200-094-088} from 13:09:11 through 13:45:29
UTC.  The root, batch, extern, and numbered steps \pathcode{.0}--\pathcode{.8} all completed
\pathcode{0:0}; the numbered task counts were exactly $2,2,2,16,2,16,2,2,2$.  Its sole completion
marker is present and scheduler stderr is empty.  The preregistered first eligible trigger was
absolute boundary 838, with rank-agreed norm 3.4058520793914795, canonical token
\pathcode{finite:0x1.b3f2f60000000p+1}, and zero health failures.  It retained the post-update
generation
\begin{quote}
\pathcode{step-000000000838-iter-000000001676-6ee7296afb63448f9863efab3454d49d}
\end{quote}
plus 16 RNG records and 16 future raw microbatches per rank.  All 16 independent loader-replay
records match every ordered joint and component-tensor hash.  The continuation later stopped
non-authorizingly at boundary 853 with finite norm 11.937856674194336; that postselection stop
does not invalidate the retained boundary-838 anchor.

The accepted run root is
\begin{quote}
\pathcode{/lustre-storage/checkpoints/ngocbh/hdn/diagnostics/rdiag_shared_dt_v1/}\\
\pathcode{diagnostic_runs/trigger_capture_015ba66e07f0_1838750}.
\end{quote}
Its summary, artifact inventory, invocation, byte-identical pre/post input inventory, capture
phase, derived anchor manifest, and checkpoint-generation manifest SHA-256 values are respectively
\begin{quote}
{\ttfamily\scriptsize 0123887d03356ca08a0e88aa267746ac1d4680b4c17a6fb63f98800175bb9b33},\\
{\ttfamily\scriptsize 0709972f30d9b4118e41865d8ecc03a55dda9f34bdc856d87d3ef23cbee48500},\\
{\ttfamily\scriptsize 04f0d067c7c52cac583f9a5ca528282a5c757c560e6135a8007eedb093586509},\\
{\ttfamily\scriptsize 729aee6bea81535c49a5c16092750d0ff603275b866fdfca662abf7ca23d3231},\\
{\ttfamily\scriptsize 05b5070b22064b56e1f3aba9403099824c45c338d3c5c6b676cb5918a89461f2},\\
{\ttfamily\scriptsize 1c01aa9798bcf3c1d75cc016bd103cac919ca7c8296cddffb4817bede8353bd8},\\
{\ttfamily\scriptsize 3b41e9f4590909f461b89aaec0041ff8b7f51347ae531bcd07faf18dba36d008}.
\end{quote}
The artifact directory has exactly 164 inventoried files and 165 physical files including the
inventory, all mode 0400 and single-link.  Its decision is
\pathcode{CAPTURE_RETAINED_NONAUTHORIZING} with
\pathcode{CAPTURED_LOADER_REPLAY_MATCHED}.  The external scheduler-record root is
\begin{quote}
\pathcode{/lustre-storage/checkpoints/ngocbh/hdn/diagnostics/rdiag_shared_dt_v1/}\\
\pathcode{diagnostic_runs/trigger_capture_015ba66e07f0_1838750_scheduler_record},
\end{quote}
whose root-row, full-step, completion-marker, and scheduler-log-hash records have SHA-256 values
\begin{quote}
{\ttfamily\scriptsize 6987ed8bf54cc5b80713503dce990fad621032a4c0d2a4d3c1116c1fd0828494},\\
{\ttfamily\scriptsize 3afc238b32b6631cba1bd64b15dd29f012c366c2f47b60789ab754fa57e8428a},\\
{\ttfamily\scriptsize f00be3ff84a1349b7bb29cb5a5460de217f3c932d7389d736e79809dc6d6b6ba},\\
{\ttfamily\scriptsize 7fb33026b6196a9d9321774c1650e1d41f302c4ff60a8dcc1e3f5a9b2fc25773}.
\end{quote}
Scheduler stdout is
{\ttfamily\scriptsize fb69191ebd7bbc42e9a20222629ae4f03525badc5df4544fb59871f02c5b2b50};
stderr has the empty-file digest.

A separate read-only validator, job \pathcode{1838966}, then ran on
\pathcode{h200-002-213} from 13:52:43 through 14:01:15 UTC and completed
\pathcode{0:0}.  It reconstructed the invocation-bound historical path
\pathcode{/tmp/rdiag_sdt_trigger_1838750/source}, verified the immutable 4,253-file source before
and after validation, revalidated source, generation zero, and anchor, returned
\pathcode{VALID} with \pathcode{artifact_count=164}, removed both owned node-local roots, emitted
one marker, and left empty stderr.  Its evidence root is
\begin{quote}
\pathcode{/lustre-storage/checkpoints/ngocbh/hdn/diagnostics/rdiag_shared_dt_v1/}\\
\pathcode{diagnostic_runs/trigger_capture_015ba66e07f0_1838750_independent_validation}.
\end{quote}
The retained validator runner, validation result, acceptance result, root-row, full-step,
completion-marker, scheduler-log-hash, and complete evidence-list SHA-256 values are respectively
\begin{quote}
{\ttfamily\scriptsize ce3ee70f2d5f80a9b98ab7fc692ad93d139d3d319115e4b3b889ff8478f2b600},\\
{\ttfamily\scriptsize 2f50dd577f9b5a507280c9ec913608b22f8f56a7028dc0f852a66848b544be42},\\
{\ttfamily\scriptsize 7e564d9f4131bbc820390cc798929739d8e291cb9535d1196e3a6e64d1511314},\\
{\ttfamily\scriptsize 254091248607cbbd8bb476ae5350029974fdb6be1d287fe889cf91c7dece84c9},\\
{\ttfamily\scriptsize acde647a47074436c368fcde4cce661b49ae4bb250a53786d87a0f00e73a0444},\\
{\ttfamily\scriptsize 0a2257b4e2b4f8c063059b52c7be11d353a2fa39a547c448d253bf6c4f247565},\\
{\ttfamily\scriptsize d455f9590ade96882b0cffe52132ef96977202c5e26bd559dbf5a084907647d3},\\
{\ttfamily\scriptsize 35e2822bc0a296487522dc3f6fe8a10b17cbe196085c5227b2e095425c393cbb}.
\end{quote}
This closes the trigger-capture prerequisite for implementation; it remains diagnostic and cannot
authorize production by itself.

The downstream matched-intervention design and implementation are isolated in
worktree
\begin{quote}
\pathcode{/lustre-storage/checkpoints/ngocbh/hdn/diagnostics/rdiag_shared_dt_v1/}\\
\pathcode{source_worktrees/rdiag_shared_dt_matched_intervention_20260904}
\end{quote}
on branch \pathcode{codex/rdiag-shared-dt-matched-intervention-20260904}; its plan is
\pathcode{docs/plans/}\\
\pathcode{2026-09-04-rdiag-shared-dt-matched-intervention-plan.md}.  The original reviewed and
pushed plan-only identity was
\begin{quote}
commit \pathcode{2b8066717dc05128a37404def1dcfa643fd0ef7d},\\
tree \pathcode{6aed24664104fa61833f95a7b5cc3626f1baea5d}.
\end{quote}
It requires one immutable post-trigger anchor and sequential fresh shared-A1, connected-zero
no-rotation, and shared-A2 arms over 16 retained microbatches per rank.  Exact A1/A2 equality is a
precondition for interpreting the intervention.  Review corrected the initial draft to compare
rank-local FSDP state only with the corresponding rank across arms, use replica-aware global
gradient diagnostics, limit first-boundary DT/alpha equality to the treated layer, stop only after
the eighth update is fully committed, capture nonfinite-loss exits artifact-locally, retain the
ordered 32-loss boundary estimand, and bind all 256 rank/batch joint hashes.  That original plan's
SHA-256 is
\begin{quote}
{\ttfamily\scriptsize 17a7d50a0349ed1038b4c703dfe25a12e791e5da40c194b5493a46ee087ca0db}.
\end{quote}
At that stage, a read-only merge preflight found no textual conflicts, and the branch deliberately
did not absorb the plan-only trigger tip; it remained gated on a successful final trigger capture
and repin before implementation.

Before observing the trigger outcome, the plan was tightened to pin accepted generation-zero job
\pathcode{1837446}, current trigger commit \pathcode{4cbddbc...} and manifest
\pathcode{e27fb7c3...}, and the still-pending retry without claiming its result.  It now fixes the
32-loss ordering and arithmetic, phase-dependent unavailable fields, value-exact zero semantics,
first-boundary frontend scope, early-stop estimands, operational-failure continuation, nonfinite
loss interception, and reconstruction of the trigger invocation's historical source path.  The
reviewed and pushed revision is
\begin{quote}
commit \pathcode{444946e6b1f588c4bcfffc083c54ac80ac1cca6d},\\
tree \pathcode{ffa896e55eebdb48f15c4b71ae2bb89da117c2e3},
\end{quote}
with plan SHA-256
\begin{quote}
{\ttfamily\scriptsize 43dc150759e293f39e810f42cc4d6dae11909d5ce28d10f4753a5fa6c2eefa51}.
\end{quote}
The acceptance clause explicitly requires both a structurally valid trigger closure and
\pathcode{CAPTURE_RETAINED_NONAUTHORIZING} with
\pathcode{CAPTURED_LOADER_REPLAY_MATCHED}; a valid no-anchor or loader-mismatch outcome cannot
release implementation.  Because revision \pathcode{444946e...} pins failed job
\pathcode{1838250}, it is now a historical blocked design record.  The reviewed plan was repinned
without implementation to pending job \pathcode{1838750}, trigger commit \pathcode{015ba66e...},
tree \pathcode{d1d59c7a...}, manifest \pathcode{bf0b5288...}, the exact historical source path
\pathcode{/tmp/rdiag_sdt_trigger_1838750/source}, and the corrected retained-anchor closure of 164
inventoried files (165 physical files including the inventory).  That reviewed and pushed plan is
\begin{quote}
commit \pathcode{41de0ee4836fa492e560f0d1640f6cac73c13382},\\
tree \pathcode{4ebdfc9335b9f4c2f11c6e2c86e98a48eca21587},
\end{quote}
with plan SHA-256
\begin{quote}
{\ttfamily\scriptsize e7a4df5999f17cae172d9de2f454299fca38879773e225690a14c2aed2c91ff8}.
\end{quote}
That repin was the last pre-outcome design state.  After both job \pathcode{1838750} and external
validator \pathcode{1838966} passed, the accepted trigger source was merged without conflict as
\begin{quote}
commit \pathcode{9d65a587be64f0a7dc9b2c3b3ed40ac0351c08ae},\\
tree \pathcode{53136608092429aad484efcaffed22cb021233cf}.
\end{quote}
The plan was then independently checked against every retained job, source, scheduler, artifact,
anchor, and validator identity and pushed as
\begin{quote}
commit \pathcode{285c36b9e75c897521e6529cd0e5766c2fcafbb7},\\
tree \pathcode{3787abb2b4480cf73cc581e718a3558706a6dd5d},
\end{quote}
with plan SHA-256
\begin{quote}
{\ttfamily\scriptsize 3b869d76eb17fafcca9ddb1364524927be4a924642f726d8a27649b076c98355}.
\end{quote}
The reviewed implementation is commit
\begin{quote}
\pathcode{49f06078d19f349ca9724e3bbfc08fa436e91e54},
\end{quote}
with tree \pathcode{ebb54a1719a6bf8490e9d6997ec471b5201423f3}.  Its implementation paths are
\pathcode{scripts/analyses/rdiag_shared_dt_matched_intervention.py},
\pathcode{scripts/analyses/run_rdiag_shared_dt_matched_intervention.sbatch}, and
\pathcode{scripts/tests/rdiag/test_rdiag_shared_dt_matched_intervention_host.py}.  The final plan,
helper, runner, and host-test SHA-256 values are respectively
\begin{quote}
{\ttfamily\scriptsize 13d5bf40a3c5f5974a64a91b4cf3ae72efd999e1bbf012f2e57fc3bac32643c3},\\
{\ttfamily\scriptsize 0a339701716fcb5d404deecd849d36a3cbd40eee654520a229801de90204b180},\\
{\ttfamily\scriptsize ceaae69aacba7b12325188c7b50d2e503977968f0022b777cfc79c331b40f935},\\
{\ttfamily\scriptsize 09002a176e6eceadf504746b774d0dcf50db08b02b616cc53c5e40c722964482}.
\end{quote}
All implementation steps before publication are complete.  The focused host suite passed 126
tests, and the matched plus trigger, generation-zero, and soak-replay sweep passed 295 tests;
Python compilation, shell syntax, diff hygiene, direct validation of all 16 retained rank-local
initial states, and live scheduler revalidation of jobs \pathcode{1838750} and
\pathcode{1838966} also passed.  Independent scientific, distributed-runtime, and
immutable-artifact reviews found no remaining code blocker after repairs to collective forwarding,
post-update evidence, rank-local failure coordination, exact anchor binding, nonfinite-loss
rederivation, contiguous-prefix estimands, and fixed-order arm continuation.

The immutable source snapshot is
\begin{quote}
\pathcode{/lustre-storage/checkpoints/ngocbh/hdn/diagnostics/rdiag_shared_dt_v1/}\\
\pathcode{source_snapshots/}\\
\pathcode{rdiag_shared_dt_matched_intervention_}\\
\pathcode{49f06078d19f349ca9724e3bbfc08fa436e91e54}.
\end{quote}
Its \pathcode{SOURCE_MANIFEST.json} SHA-256 is
\begin{quote}
{\ttfamily\scriptsize 45644d7de85703ceacfbd4789925245097b00332f1c9a9c5df8da334ce8ce3a5}.
\end{quote}
The manifest contains 4,257 regular Git files; the physical frozen tree contains 4,258 files
including the manifest and 409 directories, with directories and executables mode 0555 and other
files mode 0444.  The \pathcode{paper} gitlink is recorded and intentionally excluded.  The
publisher's built-in verification and a separate reconstruction into a private source tree both
matched the commit and manifest exactly.  Because Lustre used the supported no-replace fallback,
the complete read-only hidden staging tree is retained as publication evidence; it is not an
execution source.

Exactly one two-node/16-H200 matched-intervention allocation, job \pathcode{1839531}, was submitted
from that named snapshot at 2026-09-04 15:37 UTC.  It ran on
\pathcode{h200-071-128} and \pathcode{h200-074-040} from 15:59:07 through 17:06:49 UTC and
completed \pathcode{0:0}.  The root, batch, extern, and numbered steps \pathcode{.0}--\pathcode{.16}
all completed \pathcode{0:0}; the numbered task counts were exactly
$2,2,2,2,2,16,2,2,2,16,2,2,2,16,2,2,2$.  The run executed fresh shared-A1, connected-zero
no-rotation, and shared-A2 arms in that order, with no requeue and no production authority.  Its
run root is
\begin{quote}
\pathcode{/lustre-storage/checkpoints/ngocbh/hdn/diagnostics/rdiag_shared_dt_v1/}\\
\pathcode{diagnostic_runs/matched_intervention_49f06078d19f_1839531}.
\end{quote}
The in-job final validator returned \pathcode{VALID}, revalidating the matched source, accepted
trigger anchor, and all prerequisites.  The artifact inventory contains 118 closed files and the
directory contains 119 physical files including the inventory; every artifact is a mode-0400,
single-link regular file.  The summary and inventory SHA-256 values are respectively
\begin{quote}
{\ttfamily\scriptsize 140ca76485e6838c6ae3f6bfc086e0112aae5c48f809cfaa06adc92724efe83c},\\
{\ttfamily\scriptsize bdd78e22df0ac67242c36f559b9ede1c1ae3485b575bc3647d25285599d62765}.
\end{quote}

The validity checks passed exactly.  A1 and A2 are byte-identical under the complete common-arm
projection, with \pathcode{a1_a2_first_difference=null}; initial state, topology, all 256
rank/microbatch inputs, and pre-forward RNG records match.  All three arms committed eight updates
through boundary 846.  The no-rotation arm has exact-zero angle projection output, phase increment,
phase, and current angle gradient while the radial gradients remain live; layer-0 \pathcode{DT} and
\pathcode{alpha} match the normal arm exactly at boundary 839.  Consequently the summary reports
\pathcode{intervention_isolation=MATCHED}, no isolation failures, and causal estimands on every
boundary 839--846.

\begin{table}[ht]
  \centering
  \small
  \caption{Matched boundary-838 intervention, job 1839531.  ``On'' is the exactly repeated normal
  shared-DT arm and ``off'' is connected-zero rotation.  Only boundary 839 is an instantaneous
  counterfactual; later rows are paired longitudinal outcomes.}
  \label{tab:shared-dt-matched-1839531}
  \begin{tabular}{@{}rrrrrl@{}}
    \toprule
    Boundary & $g_{\rm on}$ & $g_{\rm off}$ & $g_{\rm on}-g_{\rm off}$ & Ratio & Limit-10 class \\
    \midrule
    839 & 1.423852 & 0.565017 & 0.858835 & 2.5200 & neither \\
    840 & 1.352574 & 0.451292 & 0.901282 & 2.9971 & neither \\
    841 & 1.605586 & 0.371823 & 1.233762 & 4.3181 & neither \\
    842 & 2.048217 & 0.338519 & 1.709698 & 6.0505 & neither \\
    843 & 1.848073 & 0.317264 & 1.530809 & 5.8250 & neither \\
    844 & 2.508924 & 0.304979 & 2.203945 & 8.2266 & neither \\
    845 & 4.601322 & 0.301506 & 4.299816 & 15.2611 & neither \\
    846 & 4.129866 & 0.277695 & 3.852171 & 14.8720 & neither \\
    \bottomrule
  \end{tabular}
\end{table}

At boundary 839, enabling rotation therefore raises the pre-clip norm by
$0.8588348031044006$, or $2.52001478149663\times$, at this one captured state/input.  The
A-minus-B global-loss contrast there is $-0.00017467141151428223$.  All eight norm differences are
positive, and the seven later signs agree with the first sign.  This is persistent longitudinal
separation in the observed window, but neither arm crosses 10 at any boundary.  It therefore does
not establish that rotation caused the historical limit crossing, does not explain the later
boundary-853 event, and does not justify a threshold change.  The summary remains explicitly
non-authorizing.

The retained scheduler record is
\pathcode{diagnostic_runs/matched_intervention_49f06078d19f_1839531_scheduler_record}.  Its
root-row, full-step, completion-marker, and scheduler-log-hash records have SHA-256 values
\begin{quote}
{\ttfamily\scriptsize 01713739b41fbc83a04e5a449cdbf702103c3b87568d8087adc1887c178c2053},\\
{\ttfamily\scriptsize 74ec9cb245a3086a1ad7afaa0da6479259492b730befaa42f951a1565b5ba672},\\
{\ttfamily\scriptsize 5346f5e235063d09bcf8e870fb98a19a05d2e099f99ab0ab8cd5a3735bcda2cd},\\
{\ttfamily\scriptsize 23b49df7cf3b72bb6de57273350ce165ce96676015f6e54cf30f70bad4905c38}.
\end{quote}
Scheduler stdout contains exactly one completion marker; its SHA-256 is
\begin{quote}
{\ttfamily\scriptsize 2760327de1c640203d5bc1552f0807fe298743b5f702997e9b022e51322d38ae}.
\end{quote}
Scheduler stderr is empty.  Fresh-node independent validator job \pathcode{1840069}, submitted from
retained runner SHA-256
\begin{quote}
{\ttfamily\scriptsize 9e602deac9319a7e12d1d71fbcba75bd6d52da5cdbfcfafd77d34661ff3e7c84},
\end{quote}
ran on
\pathcode{h200-021-171} from 18:30:14 through 18:39:40 UTC and completed \pathcode{0:0}.  It
reconstructed both invocation-bound source paths, validated all 4,257 matched-source and 4,253
trigger-source files before and after use, independently returned \pathcode{VALID}, rederived the
118-file inventory, and removed all three marker-owned node-local roots before its sole completion
marker.  Its evidence root is
\pathcode{diagnostic_runs/matched_intervention_49f06078d19f_1839531_independent_validation/}
\pathcode{evidence-job-1840069}.  The validation result, acceptance result, and complete 16-row
evidence-list SHA-256 values are respectively
\begin{quote}
{\ttfamily\scriptsize 55d0fed3dcabcbf9a60a444cce149108f4ec891a428946269a84fcb1330fcbaf},\\
{\ttfamily\scriptsize 364d899e1a7e0b999a68917780c562d00638226a6d5640653ecf1b53ad6c3ccf},\\
{\ttfamily\scriptsize 19d802d2e8a49f93d9fd8124078f307aa31f6cdf7bb5354bf478757672cb716f}.
\end{quote}
The retained validator scheduler root-row, all-step, marker, and log-record SHA-256 values are
\begin{quote}
{\ttfamily\scriptsize e186d5a86aa67c75ec0fa070dbd5fb6c78ab3bb947e1d8e51d3ade503bb1c6b1},\\
{\ttfamily\scriptsize 5f1e6be98f469a337e6915104b2bb3b02cc8bdb54f711e3108973cd1eb536fe3},\\
{\ttfamily\scriptsize 04165dfb52dfe1f9ced189fbb18ed94bdf6b518bc6e28631fa39d8453dbf9451},\\
{\ttfamily\scriptsize 974b7bdd010fb454f4063b6b13b705e9f8bdf929032a3a01ab5c4c57c1275cfb}.
\end{quote}
All retained validator files are mode-0400, single-link regular files and its scheduler stderr is
empty.

Remediation work is isolated from the main checkout in
\begin{quote}
\pathcode{/lustre-storage/checkpoints/ngocbh/hdn/diagnostics/rdiag_shared_dt_v1/}\\
\pathcode{source_worktrees/rdiag_shared_dt_late_angle_scale_20260904}
\end{quote}
on branch \pathcode{codex/rdiag-shared-dt-late-angle-scale-20260904}.  The independently reviewed,
non-authorizing design is
\begin{quote}
\pathcode{docs/plans/2026-09-04-rdiag-shared-dt-late-angle-scale-design-final.md}.
\end{quote}
It is committed and published with commit, tree, and file SHA-256 respectively
\begin{quote}
{\ttfamily\scriptsize c57c9617f0256189bcc5a733789699bf32731906},\\
{\ttfamily\scriptsize 1279bb20cac06d7ee8b5ba231a273cb955b8874e},\\
{\ttfamily\scriptsize 549a7592ffe10b3c516f98d023ea67cb00e64ad02c8529df13e2af9325f427ba}.
\end{quote}
It selects the checkpoint-compatible, rotation-preserving effective angle scale
$0.5H^{-1/2}=0.125$ at $H=16$, while retaining direct $H^{-1/2}=0.25$ and
connected-zero controls.  It does not alter \pathcode{shared_dt_v1}.  The executable plan is
\begin{quote}
\pathcode{docs/plans/}\\
\pathcode{2026-09-04-rdiag-shared-dt-late-angle-scale-implementation-plan-final.md}.
\end{quote}
It is published with commit, tree, and file SHA-256 respectively
\begin{quote}
{\ttfamily\scriptsize 17d86199a0eccf91a4ad8c4e0a256d063efc1d85},\\
{\ttfamily\scriptsize 65bd0d9d3c3f3d24b9eb51e4fd42ad327c60e942},\\
{\ttfamily\scriptsize 8497caef38ea1fc4811b629704ac97604e3daea8da46830a2865329365ce20e5}.
\end{quote}
Two formal adversarial plan-review rounds found and drove fixes for exact H=16 oracle coverage, a
distinct candidate-above-10 terminal, real-anchor FSDP preflight, complete later-boundary liveness,
fail-closed outcome mapping, immutable H200 provenance, prospective scheduler-step counts, and four
final CLI/snapshot/validator interfaces.  The final plan contains those last repairs and requires
their named red/green tests during implementation; the second formal review predated those repairs
and therefore remained \pathcode{NEEDS_REVISION}.  Implementation seed commit
\pathcode{68932d5aed62f717b91f8a5c46b38679b70805bb} added only the seven preregistered helper, host,
H200, preflight, runner, and validator paths.  The reversible H=16 angle-scale override was then
implemented and hardened through commits \pathcode{a14168a...}, \pathcode{5253faf...}, and
\pathcode{11e7d37...}.  Fresh specification and quality reviews accepted that slice after it made
rollback best-effort and fail-closed, detected mid-arm scale drift, validated exact per-rank
topology evidence, and preserved RNG and state-dict identity.  The matched-prerequisite slice was
then implemented at \pathcode{d31ba7d...}, hardened at \pathcode{31962f3...}, and completed by the
dependency bootstrap at \pathcode{9bb7d3c...}.  Adversarial review first demonstrated that an
unlisted \pathcode{sitecustomize.py} could forge the historical child result, JSON numeric aliases
could pass Python equality, and evidence could be exchanged between hashing and parsing.  The
accepted implementation now verifies the complete 4,257-file historical source closure before and
after the child, uses trusted-directory file descriptors and single-captured bytes for evidence,
requires canonical exact-type JSON, bounds and reaps child output, and retains Python \pathcode{-S}
while explicitly appending only the active interpreter's canonical, non-writable
\pathcode{purelib}/\pathcode{platlib} directories.  The latter was necessary because the real
historical validator imports PyTorch while checking the retained anchor; a real import smoke proves
that PyTorch loads without executing \pathcode{sitecustomize.py}.  Fresh specification and quality
re-reviews accepted the slice with no remaining findings.  Its five focused adversarial tests and
140 host tests pass; nine deliberately red assertions remain for later implementation slices.  The
extended one-chain prerequisite schema and external-preflight interfaces were then completed at
commit \pathcode{323ce5fc699e8611b5440cf368830a4d7098e839}.  This slice fully verifies both the
node-private execution tree and immutable source snapshot against the reviewed Git commit, tree,
and 4,257-file manifest; binds the exact numerical and actual-anchor preflight closures; checks the
root and every scheduler step, TRES, submission, working-directory, stdout, and stderr record; and
uses a sealed semantic projection on the retained rank path without repeating live queries or full
tree walks.  Source and scheduler sub-audits found no remaining scoped blocker.  A subsequent
quality review initially flagged a same-timestamp transient-edit probe, but its minimized trace
showed that it changed an already-consumed preflight runner which was neither reopened nor executed
while changed.  The reviewer therefore withdrew that finding and returned no Critical, Important,
or Minor issue and \pathcode{Ready: Yes}; the independent specification review returned
\pathcode{ACCEPTED}.  Eleven newly focused adversarial cases pass, the specification review's nine
focused functions pass as 15 parameterized cases, and the complete host file now has 152 passing
tests with the same nine deliberately red assertions reserved for later-boundary persistence,
final artifact closure, and runner postchecks.  The candidate finite-norm safety-stop half of the
next slice was implemented at \pathcode{838c236...}, hardened at \pathcode{156098f...}, and
accepted at \pathcode{d72e10e5a886300af42f3cf83d95ec5ada5380d6}.  It preserves exact scalar dtype
through the rank norm collective, continues at equality 10, stops only candidate arms for finite
norm strictly above 10, and retains the historical continuation of old-scale and no-rotation
controls.  Rank disagreement, finite negative values, and rank-agreed NaN, positive infinity, and
negative infinity have distinct terminal reasons.  Every no-step terminal retains complete
preclip and named-gradient evidence, proves the full real gradient inventory was cleared, skips
AdamW and counter commit, and binds unchanged model and optimizer roots to the authoritative prior
state while accounting only for the scheduled learning-rate change.  Review found and repaired a
float32 threshold collapse, a forgeable one-leaf cleared-gradient record, an invalid cross-boundary
optimizer comparison across LR scheduling, a missing no-rotation availability converse, and an
initially unbound equal optimizer-root pair.  Final specification and quality reviews report no
remaining finding, \pathcode{ACCEPTED}, and \pathcode{Ready: Yes}; 23 focused tests and 178 host
tests passed at that checkpoint.

The complete liveness half was implemented through commits \pathcode{d72fcd2...},
\pathcode{6db3ade...}, \pathcode{d2891f7...}, \pathcode{de6e919...},
\pathcode{47547ca...}, \pathcode{35a5f31...}, \pathcode{68846ec...},
\pathcode{d0a19b3...}, and accepted at
\pathcode{0d28a268295dde0dc534230b6f62d763b400001b}.  It records exact compact frontend
counts for all 32 layer--microbatch observations at boundaries 839--846, exact candidate
four-category gradient liveness, and replica-factor-two Adam \pathcode{exp_avg}/\pathcode{exp_avg_sq}
liveness over the real sparse FSDP ownership pattern.  Three coordinated optimizer phase gates
prevent pre-step, raw-step, and post-health collective skew; empty shards contribute zeros, while
nonempty shards require exact FP32 state, optimizer membership, and step equal to the current
boundary.  Both no-rotation repeats are checked across all ranks and observed boundaries, including
uncommitted terminal boundaries, with exact-zero rotational activations and gradients and live
radial gradients.  Review exposed and repaired numeric type aliases, incomplete first-boundary
metadata binding, impossible count tuples, conditional collective ordering, stale optimizer state,
unfaithful empty-shard tests, a single-repeat isolation path, and a nonfinite residual-gradient
false acceptance.  Final reviews returned \pathcode{ACCEPTED} and \pathcode{Ready: Yes}; 56
focused implementation tests, 74 independent focused regressions, and 221 full-host tests pass.
Eight deliberately red assertions remain for analytics/artifact closure and runner postchecks.
The repeatability and historical-reproduction checkpoint was then implemented at
\pathcode{759b752...} and accepted after the exact-rank-type repair
\pathcode{94f83dc5a94e0862fe12cad2dea79b731245f9e0}.  All three old-scale, candidate,
and no-rotation repeat pairs now compare the complete rank record while excluding only the
repeat-specific arm label.  The accepted job-1839531 controls are projected from their already
sealed bytes into an ordered $3\times16$ digest record of 3,472 canonical bytes with SHA-256
\pathcode{c65a7608683b10403ffae4a6eebe106fa2b16e3d403206b434d429f95d1a54da}; the
four current control arms make all 64 rankwise comparisons, while candidate arms have no
historical comparator.  Review found that Python equality had allowed rank zero to be represented
as \pathcode{False} or \pathcode{0.0}; exact integer typing and paired-candidate regressions close
that gap.  Independent reviews now report \pathcode{Ready: Yes}.  The focused repeat/history and
arm-result reviews pass 25 and 16 checks respectively, and the full host file has 231 passing
tests.  The same eight deliberately red assertions remain: seven belong to the final 226-file
artifact closure and one to runner source postchecking.  The boundary-839 scale and decision
checkpoint was implemented at
\pathcode{ac0c55eb5f5cc6788f4da9d296f3346e7cb4a8d7}, completed at
\pathcode{f7ed6c88941820ddcf4289be5a20c0cc688a5eba}, and accepted after the final
post-AdamW consensus repair
\pathcode{838b04566ec5347946b9e7be7f747f088af25a0f}.  Candidate arms now retain the
live layer-0 record of twice the raw angle at boundary 839; exact per-rank and per-microbatch tensor
records prove common DT/alpha inputs and the preregistered half-angle relation without comparing
layers 1--15 after their hidden states diverge.  Canonical A1 estimands preserve signed zero and
exact endpoint behavior, distinguish the sole instantaneous boundary from longitudinal outcomes,
and report the ordered eight-ratio median/p95 together with loss, gradient-category, angle-maximum,
and DT-tail descriptives.  The pure outcome classifier applies
\pathcode{INVALID_EVIDENCE} $>$ \pathcode{VALID_FAIL_SAFETY} $>$
\pathcode{VALID_FAIL_CRITERION} $>$ \pathcode{VALID_PASS}, while every summary remains
non-authorizing.  Review found and repaired partial-prefix indexing, finite-sum overflow, omitted
descriptives, inconsistent status/reason combinations, and rankwise-repeated but globally
inconsistent safety norms, including the post-AdamW path.  Final independent reviews report
\pathcode{Ready: Yes}; 247 host tests pass and the same eight deliberately red assertions remain.
The artifact-sealing and invocation-identity checkpoint was then completed at
\pathcode{79222cf0f750a71c3d06e4506b81426f1f42fb44}.  Structurally malformed or partial
evidence remains exceptional and cannot produce a trusted summary, whereas structurally complete
semantic contradictions seal only the non-authorizing status \pathcode{INVALID_EVIDENCE}.  The
main closure is now derived as 10 fixed files plus 36 files for each of six arms: 226 inventoried
files and 227 physical files including \pathcode{artifact-inventory.json}; the historical matched
closure remains pinned at 118/119.  Final validation reads sealed captures, rederives the scientific
status, revalidates the closure after rederivation, and reports that status separately from validator
success.  Invocation and summary identity now bind the design, matched validator, both preflight
closures, exact intervention constants, two-node/16-rank execution contract, retained runner, and
source manifest without reopening the runner between hashing and contract validation.  The three
named acceptance tests and seven focused sealing tests pass; the full host file reports 257 passes
and only the deliberately deferred runner source-postcheck assertion.  An initial independent
review found no remaining issue in this slice, but a stricter immutable review found two blockers:
complete, structurally valid source/anchor/RNG/scale contradictions still escaped as structural
exceptions instead of sealed \pathcode{INVALID_EVIDENCE}, and final validation bracketed reopened
files rather than rederiving from the exact bytes checked by the inventory.  Those repairs and
direct identity-mutation regressions were completed and independently accepted at
\pathcode{3f94176c94d8c383b0723f1e8d915ab0e10fc28c}, then integrated on top of the
runner work as \pathcode{bbe9868...}.  Every semantic input is now parsed from the first bounded,
inventory-verified capture; a second pass checks only post-derivation integrity.  Complete
source/prerequisite, anchor, first-forward RNG, scale, and CUDA-device contradictions become
deterministic invalid-evidence facts, while malformed shape/count/token/hook evidence still fails
externally.  The integrated full host suite passes all 269 tests.

The main-runner portion of that next task was completed and independently accepted at
\pathcode{0ce13479ef78c854b5060609f1e4b20deb66b2ea}.  It expands the diagnostic to the
frozen six-arm order and derives the exact 29-step task vector from the runner's actual successful
control flow; adversarial duplicate arm, prerequisite, and cleanup calls are rejected.  Scientific
candidate stops remain zero-exit evidence-bearing outcomes so every arm runs, whereas the first
operational failure is preserved through owned-root cleanup and suppresses summary, inventory, and
the sole completion marker.  Current, trigger, and matched sources are reconstructed on both nodes;
the full accepted prerequisite chain and both preflight closures are checked before every arm and
again in the source postcheck, while the in-job final validation deliberately uses retained mode.
The focused runner/parser set passes 13 tests, the full host file passes all 264 tests, and shell
syntax, Python compilation, and diff checks pass.  This commit does not implement or run either
preflight or the result-bearing six-arm GPU allocation.  Numerical, actual-anchor, and fresh-node
validator surfaces are therefore the current implementation task.  No numerical or result-bearing
GPU job has yet been accepted from this worktree.

The numerical-preflight host implementation was independently accepted at
\pathcode{417fe21644000c72f91e5f39946c0de5394e11a3} and integrated as
\pathcode{3617d40...}.  Its one-H200 runner is source-bound, network-scrubbed, uses private
node-local caches, and seals seven inventoried files plus the inventory.  Exactly three H200 tests
are collected at $H=16$, key/value width 128, and rotary width 64: a bitwise direct-$0.25$
historical replay through forward, VJP, parameters, and AdamW state; a candidate-$0.125$ comparison
to the independent FP32 literal recurrence; and an update/schema/liveness check for hidden, angle,
DT, and Adam moments.  Review found and repaired stale gradient accumulation between historical
branches.  The integrated host suite passes all 270 tests, the H200 module collects exactly three
nodes, and shell/Python checks pass.  These are host-side contracts only; the three numerical tests
had not yet run on an H200 at that implementation checkpoint; their subsequent S1 execution is
recorded below.

The actual-anchor preflight was independently accepted at
\pathcode{ae328418d406eaeb099e48a7a441399044d8aae8} and integrated as
\pathcode{25e63fb...}.  Its two fresh 16-rank process groups use a metadata-only anchor resolver
and sentinel loaders: no retained batch is opened, no loader is iterated, and no forward or
optimizer update occurs.  Each rank proves one strict \pathcode{Fabric.load}, exact accepted
checkpoint identity, rederived model/optimizer Merkle equality, ordered state topology before
installation and after restoration, all 16 layer overrides, and selected-device-only CUDA RNG
preservation.  Review added dynamic guards for all three file-open APIs, bound every row to the
candidate source and pre-inventory digest, and rejected forged Merkle roots/leaves.  The runner has
the exact 13-step topology and 82/83 closure.  Its standalone full host suite passes 266 tests; after
integration with the semantic, runner, and numerical slices, all 272 host tests pass.  The
first S1 actual-anchor allocation and its fail-closed result are recorded below.

The fresh-node validator was independently accepted at
\pathcode{357febcab395e07e72c6b393c28b98848a27ea05} and integrated as
\pathcode{7d92a8f214a081c5bfeb3b39c07fe3db60217929}.  It accepts the exact
12-argument source/data/target/preflight tuple, waits for the target job to become terminal,
reconstructs all three source roots, validates the target closure in
\pathcode{fresh-recursive} mode, and preserves every sealed scientific status rather than requiring
\pathcode{VALID_PASS}.  Its scheduler parser binds root and unfiltered queries, canonical task
strings, job names, two-node membership, per-step task counts, raw resource allocation, timestamps,
and the exact 29-step topology.  Review repaired final-marker cleanup ordering, NodeList and
cross-query gaps, and raw extern-step \pathcode{AllocTRES} comparison.  The validator's own
mode-0400 evidence closure is checked before its sole stdout marker.  Frozen host verification
passes 4/4 and 13/13 focused science batches, 4/4 runner contracts, and the complete 284-test host
suite; the four accepted regression suites pass 126, 69, 52, and 48 tests respectively.

The canonical branch was pushed and immutable snapshot S1 was materialized at
\pathcode{/lustre-storage/checkpoints/ngocbh/hdn/diagnostics/rdiag_shared_dt_v1/source_snapshots/}\\
\pathcode{rdiag_shared_dt_late_angle_scale_7d92a8f214a081c5bfeb3b39c07fe3db60217929}.
Its Git tree is \pathcode{d060cd04a60ad201cc75238f5de3a8fd3c20ef83}; its 4,266-file
\pathcode{SOURCE_MANIFEST.json} has SHA-256
\pathcode{1368bb987edcd3957e9b6d2cdc5a16f9268030da97f056a41829201d4672312d}.
Snapshot-local checks reproduce the live counts exactly, and a fresh private-copy verification
passes all 4,266 entries.  Numerical attempt \pathcode{1844567} passed all three H200 tests but is
not accepted because its scheduler \pathcode{SubmitLine} began with bare \pathcode{sbatch}, rather
than the exact frozen \pathcode{/usr/bin/sbatch} executable.  The byte-identical replacement job
\pathcode{1844573} completed \pathcode{0:0} on one H200 in 2 minutes 18 seconds: all three tests
passed with zero skips, scheduler stderr was empty, the sole completion marker was present, and all
seven inventoried/eight physical artifact hashes matched.  Its accepted closure is
\pathcode{/lustre-storage/checkpoints/ngocbh/hdn/diagnostics/rdiag_shared_dt_v1/diagnostic_runs/}\\
\pathcode{late_angle_scale_numerical_preflight_7d92a8f214a0_1844573_accepted_closure.json},
SHA-256 \pathcode{fe613ec2fb146cce7f52f906fdd42cb65278f7e24ed7be12f0032e51539950bb}.
Both live-query and retained-record validators pass that closure.  Actual-anchor preflight job
\pathcode{1844581} then ran from the same immutable snapshot on nodes
\pathcode{h200-115-025,h200-126-216}.  The root and batch records failed \pathcode{1:0} after
12 minutes 56 seconds; extern completed, preliminary steps \pathcode{.0}, \pathcode{.1}, and
\pathcode{.2} completed, and cleanup step \pathcode{.3} completed.  No anchor arm or rank record,
summary, inventory, or completion marker was published.  The fail-closed exception was
\pathcode{LateAngleScaleError: accepted matched prerequisite projection changed}.  Investigation
reproduced the mismatch on the host: the already validated matched source identity has six fields,
including integer \pathcode{file_count=4257}, while the strict downstream expected projection
listed only the other five.  This is an internal provenance-schema defect rather than a numerical,
checkpoint, kernel, or optimization observation.  The minimal remediation adds the pinned file
count to both expected matched-source projections and adds a realistic regression asserting exact
six-field equality; it does not change either preflight runner, the intervention, the optimizer,
the acceptance criteria, or the scheduler topology.  The red regression reproduced the exact
projection failure, affected focused tests pass, and the complete host file passes all 284 tests
with ten warnings.  Independent review found no scoped issue and returned \pathcode{READY}.  The
repair is published as commit \pathcode{17e82ca154f0c22543cb5424fe876d1d832284b8}, tree
\pathcode{394d9b82910c1bc830561f325b9540d45c10a5a8}.  Per the frozen plan's remediation
rule, this accepted source change makes S1 and its numerical closure ineligible for later gates.
The never-before-existing immutable S2 snapshot is
\pathcode{/lustre-storage/checkpoints/ngocbh/hdn/diagnostics/rdiag_shared_dt_v1/source_snapshots/}\
\pathcode{rdiag_shared_dt_late_angle_scale_17e82ca154f0c22543cb5424fe876d1d832284b8};
its 4,266-file manifest has SHA-256
\pathcode{5f5cf4ee079c69f37e1ecb01fed9f1a04d79677f215edaf1eeb612b5f94b4e8b}.
The canonical worktree and S2 independently reproduce focused counts 4 and 13, four runner-contract
passes, all 284 host tests with ten warnings, and the four accepted regression counts 126, 69, 52,
and 48.  Python compilation and all four shell syntax checks pass; the live source diff/scope checks
are clean, and a fresh private copy verifies every one of S2's 4,266 manifest entries.  Numerical
preflight job \pathcode{1844729} completed \pathcode{0:0} on \pathcode{h200-246-014} in 2 minutes
44 seconds.  All three tests passed with zero skips; root, batch, and extern are the only scheduler
rows and all completed \pathcode{0:0}; stderr is empty; the sole marker is present; and all seven
inventoried/eight physical artifact hashes match.  Its canonical mode-0400 accepted closure is
\pathcode{/lustre-storage/checkpoints/ngocbh/hdn/diagnostics/rdiag_shared_dt_v1/diagnostic_runs/}\
\pathcode{late_angle_scale_numerical_preflight_17e82ca154f0_1844729_accepted_closure.json},
SHA-256 \pathcode{4e22001f66832c1f4d150992732ea98670e84442ae48ea6c81426243d1bd47c2}.
Independent retained, live-query, scheduler-byte, live-log, mode, and artifact audits return
\pathcode{APPROVED}.  The first S2 actual-anchor attempt, job \pathcode{1844730}, correctly passed
the repaired prerequisite projection and began its first per-node recursive revalidation on
\pathcode{h200-132-171,h200-134-017}.  During monitoring, however, an unauthorized diagnostic
\pathcode{srun} added scheduler step \pathcode{.4}; because this makes the preregistered exact
13-step topology unattainable, the attempt was cancelled after 18 minutes 20 seconds and is
ineligible regardless of its payload.  It produced no accepted closure and contributes no
scientific observation.  Clean replacement job \pathcode{1844962} was submitted from the
byte-identical S2 tuple and ran on \pathcode{h200-152-154,h200-186-088}.  It passed immutable-source
reconstruction, runtime identity, the repaired initial prerequisite projection, and the old-scale
per-node recursive prerequisite recheck.  Its first 16-rank strict restore then failed collectively
before any forward, update, retained-batch read, or rank evidence publication with
\pathcode{LateAngleScaleError: shared-DT scale-layer count changed}; root and batch ended
\pathcode{1:0}, extern and cleanup completed, and no result, inventory, marker, or accepted closure
was emitted.  Strict \pathcode{Fabric.load} and the model/optimizer Merkle checks had already
succeeded.  The cause is deterministic wrapper aliasing: Lightning's \pathcode{_FabricModule}
registers the same GPT below both \pathcode{_forward_module} and \pathcode{_original_module}, while
the helper scanned \pathcode{named_modules(remove_duplicate=False)} over the wrapper.  It therefore
observed 32 paths for the 16 logical mixers.  Existing forward-determinism code already establishes
the correct contract: resolve the semantic GPT from \pathcode{fabric_model.module}, canonicalize
the exact FSDP wrapper components, and retain a bijection between 16 semantic paths and 16 object
identities.  A faithful host reproduction confirms 32 raw paths and 16 unique objects.  C2/S2 and
its otherwise valid numerical closure are now stale for authorization.  The test-first repair
resolves and binds the semantic \pathcode{fabric_model.module} root, strips only exact
\pathcode{_fsdp_wrapped_module} path components, and requires a bijection between the ordered
forward and semantic inventories.  Regressions cover the installed Lightning wrapper, nested FSDP
blocks, same-object aliases, canonical-name collisions, incomplete/mismatched/duplicated roots,
decoy owner graphs, and missing layers.  Two independent re-reviews returned \pathcode{READY}; 17
focused tests and all 288 host tests pass.  The change is published as commit
\pathcode{711155e9a685dfbdd5a4c11954017e4b6b583ee6}, tree
\pathcode{7e3dd18b8baa9a522e86a574a0ebf74fd625a59b}.  Immutable S3 is
\pathcode{/lustre-storage/checkpoints/ngocbh/hdn/diagnostics/rdiag_shared_dt_v1/source_snapshots/}\
\pathcode{rdiag_shared_dt_late_angle_scale_711155e9a685dfbdd5a4c11954017e4b6b583ee6};
its 4,266-file manifest SHA-256 is
\pathcode{8eafc98de7ed172ec0142a9a3160df6d076755f4f291bef140ca40b730130223}.
The canonical worktree and S3 reproduce focused counts 4 and 13, four runner tests, all 288 host
tests with ten warnings, and regression counts 126, 69, 52, and 48; compilation, shell syntax,
scope, and a fresh 4,266-file private-copy verification pass.  S3 numerical preflight
\pathcode{1845249} completed \pathcode{0:0} on \pathcode{h200-246-014} in 2 minutes 44 seconds:
all three tests passed with zero skips, only root/batch/extern were present, stderr was empty, and
all seven inventoried/eight physical artifacts matched.  Its canonical accepted closure is
\pathcode{/lustre-storage/checkpoints/ngocbh/hdn/diagnostics/rdiag_shared_dt_v1/diagnostic_runs/}\
\pathcode{late_angle_scale_numerical_preflight_711155e9a685_1845249_accepted_closure.json},
SHA-256 \pathcode{6d98106493f257e94ac8c466b79e7106c4a2cd70a1434ba0664eaf2ff291ebd5};
live, retained, and record validators pass, and an independent byte-level audit returns
\pathcode{APPROVED}.  S3 anchor preflight \pathcode{1845250} ran on
\pathcode{h200-005-176,h200-246-014} for 23 minutes 18 seconds.  Its source reconstruction,
two-node topology, runtime identity, initial prerequisite, old-scale recursive prerequisite, and
cache-preparation steps all completed, but the first 16-rank no-forward arm failed on every rank;
root and batch ended \pathcode{FAILED 1:0}, extern and both cleanup steps completed, and scheduler
stderr contained only the resulting task-exit reports.  No rank payload, result, inventory,
completion marker, or accepted closure exists.  The fail-closed exception was
\pathcode{LateAngleScaleError: anchor preflight state-dict topology is incomplete}.  Strict anchor
load, model/optimizer Merkle equality, and the repaired 16-owner angle-scale discovery had already
succeeded, so this remains a pre-forward infrastructure failure rather than a numerical result.

The exact mismatch is now independently evidenced and reproduced with Lightning's installed
\pathcode{\_FabricModule}.  Its inherited \pathcode{named\_parameters()} traverses the execution
graph and emits names such as
\pathcode{\_forward\_module.\_fsdp\_wrapped\_module.transformer.h.0.\_fsdp\_wrapped\_module.attn.rdiag\_angle\_proj.weight},
whereas its overridden \pathcode{state\_dict()} delegates to the original module and the accepted
checkpoint key is
\pathcode{transformer.h.0.attn.rdiag\_angle\_proj.weight}.  The topology recorder mixed those two
namespaces, leaving their name sets disjoint; its validator correctly rejected them.  The host
suite already contained a faithful Fabric/FSDP alias fixture, but exercised it only through
scale-owner discovery; the topology tests used either a bare layer or hand-built matching names.
The C4 repair resolves the same validated semantic root used by the state dictionary,
removes only path components exactly equal to \pathcode{\_fsdp\_wrapped\_module}, rejects every
canonical-name collision, and uses the resulting names consistently for parameters, buffers,
state dictionary, optimizer groups, and initialized optimizer slots.  The forward graph must bind
the semantic root exactly once and expose the identical ordered canonical object inventory;
non-deduplicating enumeration rejects parameter/buffer aliases, optimizer groups cannot duplicate
or introduce parameters, and optimizer state cannot survive outside a known group.  The generic
collector still permits a deliberate partial optimizer used by the numerical oracle, while the
anchor-record validator separately retains its exact full-partition requirement.  New regressions
exercise an installed Fabric wrapper with a persistent buffer, two AdamW groups and sparse live
slots, the actual nested-FSDP namespace difference, canonical collisions, aliases, missing and
duplicate group membership, orphan state, mismatched roots, and forward decoys.  Ten focused
topology/owner/numerical checks pass, and two independent adversarial re-reviews return
\pathcode{READY}.  Live-worktree verification reports 293 host tests with ten warnings and the
accepted regression counts 126, 69, 52, and 48; both changed paths pass the scope and diff checks.
The repair is published as commit
\pathcode{f869d5e6db5f9f8863c6ffaeb03bc4814bbf2bd1}, tree
\pathcode{b442e6570d9fb962875888a608394c96ff5cd565}.  Its never-before-used immutable S4 root is
\pathcode{/lustre-storage/checkpoints/ngocbh/hdn/diagnostics/rdiag_shared_dt_v1/source_snapshots/}\
\pathcode{rdiag_shared_dt_late_angle_scale_f869d5e6db5f9f8863c6ffaeb03bc4814bbf2bd1}; the
4,266-file source-manifest SHA-256 is
\pathcode{a14f398377ea72c3f9674887fe59c04c408d796262784f2a9bce773cfb7ca7df}.
Snapshot-local focused and static checks pass; its complete suites reproduce 293 host tests with
ten warnings and regression counts 126, 69, 52, and 48.  A fresh private execution tree rebuilt
from Git independently verifies all 4,266 manifest entries.  C4 invalidates S3 and its otherwise
valid numerical closure, so fresh S4 numerical job \pathcode{1845517} and actual-anchor job
\pathcode{1845518} were submitted through the frozen absolute \pathcode{/usr/bin/sbatch}
interfaces.  Numerical job \pathcode{1845517} completed \pathcode{0:0} on
\pathcode{h200-246-014} in 2 minutes 45 seconds: its only scheduler rows are root, batch, and
extern; all three tests passed with zero skips, seven warnings, one completion marker, empty
scheduler stderr, and an exact seven-file inventory/eight-file physical closure.  Its independently
requeried, mode-0400 accepted closure is
\pathcode{/lustre-storage/checkpoints/ngocbh/hdn/diagnostics/rdiag_shared_dt_v1/diagnostic_runs/}\
\pathcode{late_angle_scale_numerical_preflight_f869d5e6db5f_1845517_accepted_closure.json},
SHA-256 \pathcode{ae422425c386787c639adb37d608d6396418676fbe62cc52ba7b8fc06a88fb0c};
both retained and live-query validators pass.  Actual-anchor job \pathcode{1845518} then completed
\pathcode{0:0} on \pathcode{h200-005-176,h200-186-088} in 42 minutes 11 seconds.  Its root,
batch, extern, and all 13 numbered steps completed \pathcode{0:0}; the exact numbered-step task
vector is
\pathcode{2,2,2,2,2,16,2,2,2,16,2,2,2}.  Both 16-rank arms strictly loaded the accepted
checkpoint without opening a retained batch or executing a forward or update.  All 32 rank
records prove equal loaded/restored model and optimizer Merkle roots, stable canonical
parameter/buffer/optimizer topology, exact 16-layer scale installation and restoration, unchanged
CUDA RNG, and complete cleanup.  The closure contains 82 inventoried/83 physical artifacts,
exactly one completion marker, and empty scheduler stderr.  Its independently requeried mode-0400
accepted closure is
\pathcode{/lustre-storage/checkpoints/ngocbh/hdn/diagnostics/rdiag_shared_dt_v1/diagnostic_runs/}\
\pathcode{late_angle_scale_anchor_preflight_f869d5e6db5f_1845518_accepted_closure.json},
SHA-256 \pathcode{d381c5770cbe2f545b67cb2bb1432ecc537640d4ad3846f76c3c1a001edc9f79};
retained, live-query, and scheduler-record validators pass, and independent scientific and
operational reviews both return \pathcode{APPROVED}.  A second private-tree revalidation of S4
then reproduced all 4,266 source-manifest entries.  The frozen launch tuple is retained as
\pathcode{diagnostic_runs/late_angle_scale_launch_f869d5e6db5f_operator_record.json}, SHA-256
\pathcode{8f09aa03fcb23238cb10ac5b4f98e8153adfb56a9111701e41ef7e8183a6e93c}.
Exactly one result-bearing six-arm allocation, job \pathcode{1845784}, was submitted through the
frozen absolute \pathcode{/usr/bin/sbatch} interface from the clean canonical worktree.  It requests
two nodes and 16 H200 ranks.  It ran on \pathcode{h200-186-088,h200-193-042} for 1 hour
44 minutes 13 seconds.  All 29 numbered steps completed \pathcode{0:0} with exact task vector
\pathcode{2,2,2,2,2,16,2,2,2,16,2,2,2,16,2,2,2,16,2,2,2,16,2,2,2,16,2,2,2},
including all six 16-rank arms, their prerequisite rechecks and cache cleanups, the final two-node
source/prerequisite postcheck, and private-root cleanup.  All 96 rank records were emitted.  Both
candidate repeats completed all eight updates through boundary 846, were exactly repeat-equal,
and had finite live angle/DT/decay gradients and Adam moments.  Both no-rotation repeats were
exactly repeat-equal, with all rotational activations and gradients zero and radial gradients live.
Both old-scale repeats were exactly equal, and all 64 current control records exactly reproduced
the accepted historical semantic projections.

The root job nevertheless ended \pathcode{FAILED 1:0} during local post-step summarization, so
these payload observations are not an accepted result.  The producer writes the real full
18-field \pathcode{RDiagCheckpointIdentity.as_dict()} into each initial-state record, while the
structural summary validator and its synthetic fixture still required an obsolete five-field
subset.  Consequently rank zero of \pathcode{old_scale_a1} raised
\pathcode{LateAngleScaleError: initial anchor identity schema changed}.  An independent audit
confirms that all 96 full identities are identical, equal the authoritative trigger identity, and
that all 96 complete initial states equal the corresponding retained anchor state.  Fail-closed
handling suppressed \pathcode{summary.json}, \pathcode{artifact-inventory.json}, and the completion
marker; the run retained 225 partial top-level artifacts, and no fresh-node validator was
submitted.  This is classified \pathcode{INVALID_EVIDENCE} due to a diagnostic-validator schema
defect, not as an infrastructure or scientific failure.  Twice-queried scheduler root and
all-step records were byte-stable with SHA-256 values
\pathcode{f0df12e06291b684c3507389c7c6fd62b722c99f3b4937d8b16b4c94dd83ba67} and
\pathcode{81bec94ea18ae59a5f3e60292242b77257b7465e8d1ed7f53ba6e10b67286d01}.
Live stdout and stderr SHA-256 values are respectively
\pathcode{6f7722c57b87d6d2b856fb55cbb531034bbeb37034ce34b6f3c5aeef03836479} and
\pathcode{6baf9c9b5af503ee95eed862edc3d151b6eff79b8e5c1aacf9692db771c58b7d};
the pre/post prerequisite inventories are byte-identical at SHA-256
\pathcode{7f869aaf25307e73546f14e709f6bab4aab03de846d2adf81511e75cc76c5d53}.
No conventional accepted scheduler record was invented for this failed topology.

The test-first remediation replaces the stale fixture with the real 18-field schema,
adds a producer-to-validator round trip plus missing/extra/type/hash-alias regressions, and leaves
semantic equality to the later retained-anchor comparison.  Its RED test reproduced the exact
failure; the repaired focused set passes 29 tests and all 96 retained initial states validate.
Two independent exact-diff reviews return \pathcode{READY}; the final live-worktree suite has 299
passes with ten warnings, the four accepted regression suites pass 126, 69, 52, and 48 tests, and
the four runner contracts, Python compilation, shell syntax, diff, and scope checks pass.  The
repair is published as commit
\pathcode{9f5ec8c806f79731f2f1cd5339cde93eaf5fbf20}, tree
\pathcode{e9bcd490bd96e55e0b6ddc640c430844ffef8ec2}.  Immutable S5 is
\pathcode{/lustre-storage/checkpoints/ngocbh/hdn/diagnostics/rdiag_shared_dt_v1/source_snapshots/}\
\pathcode{rdiag_shared_dt_late_angle_scale_9f5ec8c806f79731f2f1cd5339cde93eaf5fbf20}; its
4,266-file manifest SHA-256 is
\pathcode{edcfd78431d8d4c152768f905834a5872bb0a8bf549d0839daa5681a7636b061}.
A private tree rebuilt from Git verifies every S5 manifest entry.  This source change requires
fresh numerical and actual-anchor preflights and a fresh candidate allocation; S4 and job
\pathcode{1845784} cannot authorize a later stage.  Fresh S5 numerical preflight
\pathcode{1846328} and actual-anchor preflight \pathcode{1846329} were submitted through the
frozen absolute \pathcode{/usr/bin/sbatch} interfaces.  Numerical job \pathcode{1846328}
completed \pathcode{0:0} on \pathcode{h200-112-045} in 2 minutes 35 seconds: all three H200
tests passed with zero skips, only root/batch/extern were present, stderr was empty, and all seven
inventoried/eight physical artifacts matched.  Its accepted closure is
\pathcode{/lustre-storage/checkpoints/ngocbh/hdn/diagnostics/rdiag_shared_dt_v1/diagnostic_runs/}\
\pathcode{late_angle_scale_numerical_preflight_9f5ec8c806f7_1846328_accepted_closure.json},
SHA-256 \pathcode{44390184226902be6a4fd7488f7752555e1fdefbf45d097a828ae2e798b20fc9};
live-query, offline, and retained-record validators pass.  Actual-anchor job
\pathcode{1846329} completed \pathcode{0:0} on
\pathcode{h200-101-019,h200-186-088} in 43 minutes 30 seconds.  Its root, batch, extern, and
all 13 numbered steps completed \pathcode{0:0} with exact task vector
\pathcode{2,2,2,2,2,16,2,2,2,16,2,2,2}.  Across both arms, all 32 ranks strictly restored the
accepted boundary-838 checkpoint without opening data, executing a forward, or applying an
update.  The candidate changed only the installed 16-layer multiplier and effective scale from
\pathcode{1}/\pathcode{0.25} to \pathcode{0.5}/\pathcode{0.125}; model and optimizer Merkle roots,
state topology, and device RNG were restored exactly, and cleanup completed.  The pre/post
prerequisite inventories are byte-identical at SHA-256
\pathcode{fa0943d2f98c6c0dede5b6e89e5f6c30a63fdb57dd996c899711a0123bb0d669}; the closure has
82 inventoried/83 physical mode-0400 single-link artifacts, exactly one completion marker, and
empty scheduler stderr.  Its accepted closure is
\pathcode{/lustre-storage/checkpoints/ngocbh/hdn/diagnostics/rdiag_shared_dt_v1/diagnostic_runs/}\
\pathcode{late_angle_scale_anchor_preflight_9f5ec8c806f7_1846329_accepted_closure.json}, SHA-256
\pathcode{63f4e61d8c531ad1b0702185bf49a1485bea9c24941f89a8622aaf477cc22970}; live-query,
offline, and retained-record validation all pass.  Independent scientific and operational
reviews return \pathcode{APPROVED}, including a replay of all 96 S4 rank records through the S5
repair.  A final private Git reconstruction again verified all 4,266 S5 manifest entries.  The
one-shot operator record is
\pathcode{diagnostic_runs/late_angle_scale_launch_9f5ec8c806f7_operator_record.json}, SHA-256
\pathcode{d3d0347bb47117e179a42d05692ed01e7e891ab26a4969e82d3ed22cb04e3dfd}.
Exactly one S5 six-arm allocation, job \pathcode{1846561}, was then submitted through its recorded
12-token absolute \pathcode{/usr/bin/sbatch} interface from the clean canonical worktree.  It
completed \pathcode{0:0} on \pathcode{h200-101-019,h200-186-088} in 1 hour 46 minutes 18 seconds.
Root, batch, extern, and all 29 numbered steps completed \pathcode{0:0} with exact task vector
\pathcode{2,2,2,2,2,16,2,2,2,16,2,2,2,16,2,2,2,16,2,2,2,16,2,2,2,16,2,2,2}.
All six arms completed eight updates through boundary 846; all three A1/A2 repeat projections are
exact, all four control records reproduce their accepted historical projections, intervention
isolation passes, and all candidate gradient and Adam-moment liveness observations are live.  The
pre-registered classifier returns non-authorizing \pathcode{VALID_PASS}.  At the sole primary
instantaneous boundary, candidate/no-rotation norm ratio is
\pathcode{1.2335551366597468}, below the locked \pathcode{1.25} maximum.  Candidate norm is below
old-scale norm at every boundary and never exceeds the unchanged limit of 10; the boundary-839
candidate and old-scale norms are respectively \pathcode{0.6969801187515259} and
\pathcode{1.4238522052764893}.  This remains one captured state/window: only boundary 839 is an
instantaneous counterfactual, while later boundaries compare paired longitudinal branches.
The exact artifact closure is 226 inventoried/227 physical mode-0400 single-link files.  Its
inventory and summary SHA-256 values are
\pathcode{7d9e81dc67e0f1e1610135fb8ab994697afdf289b0ebd48f12b4677490f83938} and
\pathcode{7ea9caa4d6b3ecf731df73cd8fe0aaaebdbf3c5cd9e43129fd4175ecbde0672b}; scheduler stderr is
empty and stdout SHA-256 is
\pathcode{0a5915dec9377b9745c9d75523d7380d52cea8f10dc976e4f7f4a24c3580f4d3}.
Fresh-node validator job \pathcode{1846728} was submitted with
\pathcode{h200-101-019,h200-186-088} explicitly excluded.  It correctly landed on the fresh node
\pathcode{h200-052-125}, but failed \pathcode{2:0} after one second before creating a validation
root or numbered step.  Its sole stderr line is
\pathcode{target run root does not bind commit and job ID}; stdout is empty.  The submitted target
path and job ID are correct.  The frozen runner instead placed \pathcode{'\_'} inside one outer
double-quoted expected-path expansion, so Bash retained the apostrophes literally and every valid
target path fails that comparison.  This is a validator control-plane defect, not contrary
scientific evidence, but it leaves job \pathcode{1846561} unaccepted.  The repair must add a
behavioral shell regression, replace that expression with braced adjacent expansions, freeze a
new source version, and repeat the source-bound preflights, six-arm allocation, and fresh-node
validator.  No production implementation or authorization is claimed.

The validator repair is now frozen as S6 on branch
\pathcode{codex/rdiag-shared-dt-late-angle-validator-fix-20260905} at commit
\pathcode{3d9c20f4e1019d6e1117604e94622ed46bc06084} (tree
\pathcode{bcba0c6c7f1780326cf018f70c1811c8ba8c5ee2}).  Relative to S5, the diff
contains only the one-line braced target-path expansion and a behavioral host regression that
extracts and executes the actual two-line guard.  The S5 expression rejects its valid target with
status 2; the repaired expression accepts that target while retaining rejection of wrong commit,
job, and suffix values.  Numerical-preflight, anchor-preflight, six-arm, helper, and H200-test
bytes are unchanged.  The immutable S6 snapshot is
\pathcode{/lustre-storage/checkpoints/ngocbh/hdn/diagnostics/rdiag_shared_dt_v1/}\
\pathcode{source_snapshots/rdiag_shared_dt_late_angle_scale_3d9c20f4e1019d6e1117604e94622ed46bc06084};
its 4,266-entry manifest has SHA-256
\pathcode{85d1a7f8de3de4274df8904cce2f119ee88b900f5063fa6bf14efff515d325c8}.
Independent private Git reconstructions before and after testing verified every entry.  The frozen
snapshot passed 595 host regressions, helper compilation, and all four runner syntax checks.
Independent scientific and operational launch reviews both returned \pathcode{APPROVED}; a
reviewer's discarded test attempt wrote temporary pytest state into the live worktree only because
the host temporary filesystems were full, then correctly failed the source-closure guard.  Those
artifacts were moved intact to a recoverable diagnostic quarantine, the canonical worktree was
restored clean, and the post-test private reconstruction passed.  The first S6 submission attempt
mistakenly added the operator-only \pathcode{--parsable} option to both \pathcode{sbatch}
commands.  Because the frozen scheduler validator requires exact equality with the registered
six-token numerical and eight-token anchor command lines, jobs \pathcode{1846846} and
\pathcode{1846847} were immediately classified ineligible and canceled; the former never
allocated, while the latter stopped after 81 seconds during preliminary step 2, before any result,
inventory, completion marker, or accepted closure.  The four retained preliminary anchor files do
not constitute result evidence.  The exact registered interfaces, with no extra option, were then
submitted as replacement numerical job \pathcode{1846871} and actual-anchor job
\pathcode{1846872}.  Their scheduler \pathcode{SubmitLine} fields exactly match the frozen
contracts.  Numerical job \pathcode{1846871} completed \pathcode{0:0} on
\pathcode{h200-067-061} in 2 minutes 34 seconds with the expected root/batch/extern-only topology,
three passing H200 tests, its sole completion marker, and empty stderr.  Its eight physical/seven
inventoried artifacts are mode-0400 single-link files, and the result is non-authorizing
\pathcode{VALID}.  Live-query, offline, and retained-record validation all accept its source-bound
closure
\pathcode{diagnostic_runs/late_angle_scale_numerical_preflight_3d9c20f4e101_1846871_accepted_closure.json},
SHA-256 \pathcode{004867fe85bbce178b861095dde71b9a68f5c8873214c1c095093ff227fc8093}.
Anchor job \pathcode{1846872} completed \pathcode{0:0} on
\pathcode{h200-067-061,h200-128-233} in 44 minutes 28 seconds.  Root, batch, extern, and all
13 numbered steps completed cleanly with exact task vector
\pathcode{2,2,2,2,2,16,2,2,2,16,2,2,2}.  Both 16-rank arms completed their strict
restore-only checks; the pre/post prerequisite inventories are byte-identical at SHA-256
\pathcode{f8801e1163bf859bf420bb9dc25f0565e4a74e034fd92998300db03781e2eccd}.
Its 82 inventoried/83 physical artifacts are all mode-0400 single-link files, its sole completion
marker is present, and scheduler stderr is empty.  Live-query, offline, and retained-record
validation all accept its source-bound closure
\pathcode{diagnostic_runs/late_angle_scale_anchor_preflight_3d9c20f4e101_1846872_accepted_closure.json},
SHA-256 \pathcode{46d5ef721929457f3ccfbb3ad814461d06d8b3d24350442430498018cd4bb387}.
No S5 closure was reused.  Final scientific and operational Action89 reviews both returned
\pathcode{APPROVED}.  The mode-0400, single-link one-shot operator record is
\pathcode{diagnostic_runs/late_angle_scale_launch_3d9c20f4e101_operator_record.json}, SHA-256
\pathcode{4c68809091d51cecf35344e32b5a31c0946b5c8166eca1b963da0b2c33285275}.
A final private Git reconstruction again verified all 4,266 S6 manifest entries.  The exact
12-token main interface, with no \pathcode{sbatch} option, was then submitted once as job
\pathcode{1847039}; its scheduler-recorded command matches the operator record.  It completed
\pathcode{0:0} on \pathcode{h200-129-161,h200-130-169} in 1 hour 47 minutes 8 seconds.  Root,
batch, extern, and exactly numbered steps 0--28 completed cleanly with task vector
\pathcode{2,2,2,2,2,16,2,2,2,16,2,2,2,16,2,2,2,16,2,2,2,16,2,2,2,16,2,2,2}.
The retained non-authorizing summary reports \pathcode{VALID_PASS}: all six arms completed eight
updates, A1/A2 repeatability matched, topology and boundary-839 intervention isolation passed,
all locked criteria are true, causal interpretation is allowed, and every failure-reason list is
empty.  At the instantaneous boundary 839, candidate, historical-scale, and no-rotation norms are
respectively \pathcode{0.6969801187515259}, \pathcode{1.4238522052764893}, and
\pathcode{0.5650174021720886}; the candidate/no-rotation ratio is
\pathcode{1.2335551366597468}, below its preregistered 1.25 ceiling.  The candidate is below the
historical scale and below the norm-10 ceiling at all boundaries 839--846.  Its 226 inventoried
members plus the inventory are exact mode-0400, single-link files; inventory, summary, invocation,
and unchanged pre/post input-inventory SHA-256 values are respectively
\pathcode{373f6c6d20d50e5e616403d4963433a6d6f7f1630fdeca291fa5de0fd2dd6adf},
\pathcode{99e35928cb6e7d510b55c431747ad445241aad59793d1774f9a42a9d5bf7de00},
\pathcode{3f085cbd42ec6125a4fbb2f76b1f45a3311fe7e38534690cb0f06e3fec0c9582}, and
\pathcode{e28296916695c549b2ecb75f94fd1bf491e6be380338a530023c8378df6b063b}.
The completion marker occurs once and scheduler stderr is empty.  Fresh-node validator job
\pathcode{1847464} was then submitted with both main-job nodes explicitly excluded and ran on
\pathcode{h200-017-170}.  Its recursive scientific validation completed: the retained
\pathcode{validation-result.json} reports \pathcode{VALID} and revalidates all 226 artifacts,
prerequisites, anchor, and source; the subsequently written \pathcode{acceptance-result.json}
reports the target's \pathcode{VALID_PASS}.  Their SHA-256 values are respectively
\pathcode{84ff545fb479a089f547e8fec9c654b66c975ba6d2627a8685cfeed88e6c8c83} and
\pathcode{71ee18f943f3c121dbe9f0d2a899839a8be48c4c2bf69ad0218434e92baae6c6}.
The validator nevertheless ended \pathcode{FAILED 1:0} after 11 minutes 35 seconds during final
sealing.  The frozen runner invoked \pathcode{/usr/bin/printf -v}; \pathcode{-v} is a Bash-builtin
option, not a coreutils option.  External printf therefore returned zero while writing the two
bytes \pathcode{-v} to stdout and a warning to stderr, did not assign
\pathcode{COMPLETION_MARKER}, and the next expansion failed as an unbound variable.  The partial
root contains six of the required nine files and lacks its completion marker, scheduler-log hash
record, and evidence hash closure.  It is consequently \pathcode{INVALID_EVIDENCE}; the two
pre-seal JSON files are forensic corroboration only, not an accepted closure.  A focused
executable regression reproduces the exact warning and unbound-variable failure on S6.  The S7
repair replaces only that command with \pathcode{builtin printf -v} and adds the regression; the
focused test, surrounding validator-contract tests, Python compilation, and Bash syntax check pass.
The repair is committed on branch
\pathcode{codex/rdiag-shared-dt-validator-seal-fix-20260905} as
\pathcode{0a87fe50a47228c90a72d233615bee422fd2fe7a} (tree
\pathcode{a911270714350b23987e13ffb44070fbe45d1e8d}) and is pushed at that exact remote SHA.
The corresponding clean worktree is
\pathcode{/lustre-storage/checkpoints/ngocbh/hdn/diagnostics/rdiag_shared_dt_v1/source_worktrees/rdiag_shared_dt_validator_seal_fix_20260905}.
Its immutable 4,266-file snapshot is
\pathcode{source_snapshots/rdiag_shared_dt_late_angle_scale_0a87fe50a47228c90a72d233615bee422fd2fe7a};
the source-manifest SHA-256 is
\pathcode{065f6b6094b064f66975acdcb41745a4d94a8c58ee92a5ae51aca0be941ea02d}.
An independent private Git reconstruction verified the entire snapshot, and both the live worktree
and snapshot passed the frozen 596-test host suite plus helper compilation and all four runner
syntax checks.  Independent specification, code-quality, scientific, and operational reviews found
no remaining issue in the frozen S7 diagnostic/validator contract and approved the strict full S7
repetition.  No S6 closure is reused.  The exact
registered commands, without extra \pathcode{sbatch} options, were submitted from the clean S7
worktree as actual-anchor preflight job \pathcode{1847596} and numerical preflight job
\pathcode{1847597}; their scheduler \pathcode{SubmitLine} values bind the S7 snapshot, commit,
manifest, worktree, and, for the anchor, the unchanged data manifest.  Anchor job
\pathcode{1847596} completed \pathcode{0:0} on
\pathcode{h200-112-045,h200-128-233} in 42 minutes 57 seconds.  Root, batch, extern, and all 13
numbered steps completed with exact task vector
\pathcode{2,2,2,2,2,16,2,2,2,16,2,2,2}; both 16-rank arms proved strict boundary-838 restore,
zero data/forward/update activity, unchanged model and optimizer Merkle roots, preserved RNG, and
exact historical/candidate scale behavior $0.25\rightarrow0.25$ and
$0.25\rightarrow0.125\rightarrow0.25$.  Its 82 inventoried/83 physical artifacts are sealed, its
result is \pathcode{VALID}, and its accepted closure is
\pathcode{diagnostic_runs/late_angle_scale_anchor_preflight_0a87fe50a472_1847596_accepted_closure.json},
SHA-256 \pathcode{171d985217f5e119f7037a6742531838f608e65a1d502a4a7a40db46e39f4fc6}.
Numerical job \pathcode{1847597} completed \pathcode{0:0} on \pathcode{h200-128-233} in 2 minutes
19 seconds; all three H=16 historical-identity, candidate forward/VJP, and one-step live-moment
tests passed.  Its seven inventoried/eight physical artifacts are sealed, its result is
\pathcode{VALID}, and its accepted closure is
\pathcode{diagnostic_runs/late_angle_scale_numerical_preflight_0a87fe50a472_1847597_accepted_closure.json},
SHA-256 \pathcode{d4e33bdc0f5093fbe893323be31669470f32b358d0f3f9002a2a1dd0fc9844ec}.
Both closures pass live-query, offline, and retained-record validation with identical projections.
They were sealed serially with the reviewed operator-side tool
\pathcode{operator_tools/rdiag_preflight_closure_b6aeea289545a0d8d7b4fd27d447ae9dfc4b96001748b34fbb94ee882bebf11a.py},
whose SHA-256 is encoded in its filename; the same frozen bytes first reproduced both accepted S6
closure digests and then reproduced each S7 closure without \pathcode{--seal} after publication.
Independent scientific and operational Action-89 reviews returned \pathcode{APPROVED}.  A fresh
4,266-file private reconstruction then reverified the S7 snapshot.  The mode-0400, single-link
operator record is
\pathcode{diagnostic_runs/late_angle_scale_launch_0a87fe50a472_operator_record.json}, SHA-256
\pathcode{1aa59e8e0b0487c49636d2271f9abca567746219a594e6db014678bd41087957}.
The two accepted preflight closures and this operator record all explicitly carry
\pathcode{authorizing=false}: they authorize only the frozen diagnostic repetition, never a
production-training launch.  The record's exact 12-token
main command, with no \pathcode{sbatch} option, was submitted once as job \pathcode{1847784}; the
scheduler-recorded \pathcode{SubmitLine}, worktree, source, data, and both closure pins match the
operator record.  It ran from 00:40:07 UTC for 1 hour 43 minutes 44 seconds on
\pathcode{h200-053-020,h200-140-109} and completed \pathcode{0:0}.  Root, batch, extern, and all
29 numbered steps completed with exact task vector
\pathcode{2,2,2,2,2,16,2,2,2,16,2,2,2,16,2,2,2,16,2,2,2,16,2,2,2,16,2,2,2}; the completion
marker occurs exactly once and scheduler stderr is empty.  Its run root is
\pathcode{/lustre-storage/checkpoints/ngocbh/hdn/diagnostics/rdiag_shared_dt_v1/diagnostic_runs/late_angle_scale_0a87fe50a472_1847784}.
All 96 rank records report eight committed updates at boundaries 839--846, terminal completion,
and no global non-finite loss decision.  The non-authorizing summary reports
\pathcode{VALID_PASS}, exact A1/A2 repeatability, valid boundary-839 scale isolation, matched
topology, causal interpretation allowed, and empty criterion/evidence failure lists.  Its SHA-256
is \pathcode{ba980ce6045b0ec84ab31ca4b37a4ecc56414240b82388397afb19584668d438}; the 226-member
artifact inventory SHA-256 is
\pathcode{774fb297aec553de0e0ab3e335d5f835626a99c636c387c31c71357f3501ccfc}, and the pre/post
input inventories are byte-identical at SHA-256
\pathcode{62ce767d73abd2b8528444db5f11215a8eb1bda5c4c3188e5e181a4db6457f42}.
After an independent terminal audit recomputed every inventory member, fresh-node validator job
\pathcode{1848162} was submitted with the main nodes as its sole scheduler exclusion.  It completed
\pathcode{0:0} in 11 minutes 46 seconds on the disjoint node \pathcode{h200-178-069}.  The exact
nine-file validation root is
\pathcode{/lustre-storage/checkpoints/ngocbh/hdn/diagnostics/rdiag_shared_dt_v1/diagnostic_runs/late_angle_scale_0a87fe50a472_1847784_independent_validation};
all members are mode-0400, single-link files, and its evidence list recomputes exactly (list-file
SHA-256 \pathcode{779d58728154d0e12de6315bd6196be64e697d1d5fc712858f813d647b16afc7}).
The validation result is \pathcode{VALID} at SHA-256
\pathcode{6cc0a46b2f801bf40eaa315fad49f9e3788911bf50b0b6b1125cdcc38a96c036}; the acceptance result
is \pathcode{ACCEPTED} at SHA-256
\pathcode{d33a913fa67915812e5f548b790fe20d24ad672fe29bf4fef422971fbdc87542}.  Validator stdout is
exactly the sole completion-marker line and stderr is empty.  Independent replay therefore accepts
S7 commit \pathcode{0a87fe50a47228c90a72d233615bee422fd2fe7a} as the implementation base.  The closure remains
explicitly \pathcode{authorizing=false}: it permits v2 implementation work, not production training.

A broader preparatory 539-test host baseline also exposed one pre-existing, platform-sensitive
test-harness defect: 531 tests pass and seven skip, while
\pathcode{test_retention_journal_resumes_partially_deleted_generation} fails on Lustre.  The test
globally intercepts every directory-relative deletion of \pathcode{model.pth}; because Lustre
rejects \pathcode{renameat2(RENAME_NOREPLACE)} with \pathcode{EINVAL}, the injection fires during
the claim-and-copy staging cleanup before pointer commit rather than during the intended
post-commit retention cleanup.  The resulting ordinary checkpoint error is therefore correct;
the future v2 branch must scope the injection to the retained victim directory's device/inode and
assert the new pointer before the resume check.  This host-test repair is separate from the frozen
S7 diagnostic source and does not alter the completed and externally validated experiment.

The completed prerequisites and diagnostics are the forward localization in job \pathcode{1836996}, the reviewed
shared-DT-only warp-8 RMSNorm-forward control, its passing bounded smoke in job
\pathcode{1837393}, and the passing replacement generation-zero gate in job
\pathcode{1837446}, followed by accepted trigger capture \pathcode{1838750} and independent
validator \pathcode{1838966}, the valid matched intervention \pathcode{1839531}, and its fresh-node
validator \pathcode{1840069}.  The remaining
implementation surfaces for the additive version are
\pathcode{lit_gpt/rdiag_shared_dt_kla.py}, \pathcode{lit_gpt/config.py},
\pathcode{lit_gpt/model.py}, \pathcode{lit_gpt/rdiag_shared_dt_rmsnorm.py},
\pathcode{lit_gpt/rdiag_shared_dt_training_health.py}, \pathcode{lit_gpt/rdiag_checkpoint.py},
\pathcode{pretrain.py}, \pathcode{scripts/train/rdiag_kla/rdiag_recipe.py}, and the adjacent
campaign, source-gate, qualification, authorization, job-control, run-guard, supervisor, and
evaluation modules.  These changes were made in a new worktree based on the validator-accepted
S7 commit; the frozen S7 worktree above remains untouched.  That worktree was created
on branch \pathcode{codex/rdiag-shared-dt-v2-20260906} at
\pathcode{/lustre-storage/checkpoints/ngocbh/hdn/diagnostics/rdiag_shared_dt_v2/source_worktrees/rdiag_shared_dt_v2_implementation_20260906},
with initial HEAD exactly \pathcode{0a87fe50a47228c90a72d233615bee422fd2fe7a}.
The implementation plan is
\pathcode{docs/plans/2026-09-06-rdiag-shared-dt-v2-production-plan.md} in that worktree.
The live namespace-aware execution ledger is
\pathcode{/lustre-storage/checkpoints/ngocbh/hdn/diagnostics/rdiag_shared_dt_v2/releases/v2_b37b3c4d4205ccb04ade7000cb7f3b16c42ca547/OPERATOR_STATUS.md};
its empirically corrected allocation count supersedes the static estimate frozen in the source plan.
All implementation and host-verification slices are frozen in reviewed source commit
\pathcode{b37b3c4d4205ccb04ade7000cb7f3b16c42ca547}, tree
\pathcode{4c2b2bd8cb645d342c0e605a926ccb474c961024}, and the exact pushed branch above.  The immutable
qualification snapshot is
\pathcode{/lustre-storage/checkpoints/ngocbh/hdn/diagnostics/rdiag_shared_dt_v2/releases/v2_b37b3c4d4205ccb04ade7000cb7f3b16c42ca547/qualification_source_parent/source_b37b3c4d4205ccb04ade7000cb7f3b16c42ca547};
its 4290-file source manifest has SHA-256
\pathcode{94207969a5281d42064006e7b9b177438bfba6589eb27eb4841e78a243628211}.  Independent reopening
revalidated the commit, manifest digest, and complete file closure.  Lustre used the documented
retained-complete-read-only publication fallback, so this qualification snapshot is immutable but
is not itself a production snapshot or authorization.  The earlier Lustre-only checkpoint-test
failure was repaired by forcing the
\pathcode{renameat2(RENAME_NOREPLACE)} fallback, scoping the injected unlink failure to the
retained victim generation's device/inode, and proving that the new pointer resolves before
retention reconciliation.  The complete checkpoint file then reported 81 passes and three skips;
after the versioned identity integration, the combined checkpoint, generation, and pretrain suite
reported 128 passes and three skips.  The additive layer accepts only the exact
\pathcode{shared_dt_v1} and \pathcode{shared_dt_v2} strings, retains the v1 scale
$H^{-1/2}$, and assigns v2 the sole changed scalar $0.5H^{-1/2}$ at the existing raw-angle
multiplication site before the unchanged FP32 \pathcode{tanh}.  Four field-identical v2 configs,
apart from name and clock mode, dispatch to the same layer; tests establish identical parameters,
buffers, RNG consumption, optimizer topology, and parameter counts.  The RMSNorm policy,
first-update health, checkpoint-run, and first-update-marker identities now have explicit v2-family
schemas.  Crossed model-name/mode, schema/mode, v1/v2 resume, and string-subclass aliases fail
closed before Lightning payload loading, while pinned v1 canonical checkpoint and health bytes
remain unchanged.  Independent core/health and checkpoint/pretrain reviews both returned
\pathcode{APPROVED}; no file below \pathcode{lit_gpt/kla_ops/} changed.  Parallel v2
recipe/catalog identities, separate H200 numerical and train-smoke harnesses, all four qualification
gates, six-seal aggregation, exact authorization, run guard, supervisor, dedicated v2 submitter,
and full seven-mode evaluation path are implemented.  Exact family maps choose one v1 or v2
runner and reject crossed schemas, model names, checkpoints, manifests, seals, and string
subclasses before side effects.  The existing v1 submitter and runner remain byte-identical to the
accepted base.

The first full broad run on Lustre reported 1125 passes, seven skips, and three failures.  Root-cause
tracing found test assumptions rather than production-code defects: one source-snapshot assertion
assumed atomic no-replace rename instead of the documented retained-read-only Lustre fallback, a
data-manifest test requested a 1 ms mtime mutation on a whole-second-resolution mount, and a
pretraining source assertion still expected the retired cross-family recipe lookup.  The tests now
validate both snapshot publication success dispositions, force and verify a distinct filesystem
timestamp, and require manifest-schema-to-family-to-recipe lookup.  The merged regressions pass
3/3 on Lustre; the clean full rerun reports 1132 passes and three skips.  All 55 changed or new
Python files compile, the new JSON parses, eight shell or sbatch files pass syntax checks,
\pathcode{git diff --check} is clean, and \pathcode{lit_gpt/kla_ops/} still has no diff.  An
independent whole-diff review found no critical or important issue and returned \pathcode{READY}.
The 266.7 GiB corpus/tokenizer closure was first hashed inside the Codex UID namespace into
data-manifest SHA-256
\pathcode{adc7ad730a1d4dfcf16721733bcfca93d70c8108d7a7ec7f6eabafd3d66e3466}.  Production-H200
campaign \pathcode{b76e18cce1d842618b910ed164dfa9b2}, execution digest
\pathcode{3cb5cab2c3b671e25322222fb3f50de4bbff445ae55aa35a93c6e9352aafa90f}, was submitted as
one-H200 job \pathcode{1849816}; it failed closed in 31 seconds before the scientific test.
Diagnostic job \pathcode{1849835} established the exact namespace crossing: the manifest recorded
data and tokenizer owner UID/GID 0, whereas the compute node correctly observed UID/GID 741838.
No data bytes or source identity differed.  Compute-side manifest regeneration job
\pathcode{1849849} completed in 44m35s and produced corrected manifest SHA-256
\pathcode{47e8650e20025ce8dc3d2f261b25f1493b991b34f466e644b7eb7795659641e7}.  Compute-side
preparer \pathcode{1850489} then completed and created fresh campaign
\pathcode{4234b79ab0a8440f9e1532ae0baf59c8}, execution digest
\pathcode{0fad65d66b54fc5d618378a1ff6c686f564f684cc32eac5683826c98639a2e53}.  Its exact emitted
submission became workload \pathcode{1850491}, which completed in 29m09s on
\pathcode{h200-046-019}; both numerical cases passed, and all source, runtime, cleanup, and result
checks returned zero.  Fresh external attestor \pathcode{1850765}, submitted with that workload node
excluded, completed in 46m26s on the disjoint node \pathcode{h200-042-061}.  It rehashed the full
corpus and published the read-only production-H200 seal with SHA-256
\pathcode{65fbe8abc741b2f2f08e4f3d5536e28178b161bd219c6a29b6e284a1529cc463}.  Independent
third-node control job \pathcode{1851164} then reopened the seal on \pathcode{h200-004-175} and
returned \pathcode{VALID_PASS}.  The production-H200 gate is therefore accepted.  Compute-side
preparers \pathcode{1851167}--\pathcode{1851171} all completed \pathcode{0:0} and emitted immutable
execution manifests for the same campaign.  Their exact submission vectors became train-smoke,
checkpoint-fault, authenticated-resume, gradient-A/B, and two-node 3300-update-soak workloads
\pathcode{1851172}--\pathcode{1851176}.  Train-smoke job \pathcode{1851172} and checkpoint-fault job
\pathcode{1851173} completed \pathcode{0:0} in 8m42s and 4m52s, respectively, and each published a
mode-0400 passing result.  Their node-excluding external attestors
\pathcode{1851298} and \pathcode{1851291} completed \pathcode{0:0} on
\pathcode{h200-094-088} and published the train-smoke and checkpoint-fault seals with SHA-256
\pathcode{fa8812b932666aa3c654c5e724236c31ff057c844d7c71b497034156842ad830} and
\pathcode{3e15913d33b459ae293f1dbcd0542ede550b778deb211b2fcc3e755da60eada4}, respectively.
The two-node soak \pathcode{1851176} ran on
\pathcode{h200-053-020,h200-124-110}; it subsequently failed closed \pathcode{1:0} after 22m33s.
Its authenticated journal contains 904 sequential finite updates through iteration 1808 with zero
health failures; the next update was rejected because its pre-clip global gradient exceeded the
declared limit 10.  The sole retained scientific artifact is that mode-0400 journal, with SHA-256
\pathcode{ddefb5c6d8d15eab251ed84523f021dc19a6a5b1c7e4198eee6e5657ea0fcd0e}.
Cleanup completed, but no soak result or seal was published.  This is a scientific
gate rejection, not a scheduler failure, and the attempt will neither be attested nor reused.  Resume
\pathcode{1851174} and gradient-A/B \pathcode{1851175} started together at
\pathcode{2026-09-06T09:23:02Z} on \pathcode{h200-094-088} and
\pathcode{h200-052-125}, respectively; both first reverified the immutable b37 source snapshot.
Gradient-A/B completed \pathcode{0:0} in 6m21s with passing result SHA-256
\pathcode{68734fcfcf78dfcbecc74b179a35d0fc3a79b5cbdc5edb835107c8bc00fd007c}; its disjoint-node
attestor \pathcode{1851656} then completed \pathcode{0:0} on \pathcode{h200-053-020} and published
seal SHA-256 \pathcode{85ea2ced80e2e43bedff943edd2f40e5160ff8676b07bc8304c484be65a7eac9}.
Authenticated resume completed \pathcode{0:0} in 12m10s after both its fresh and authenticated
resumed segments passed, publishing result SHA-256
\pathcode{05ae105d9e81e13439cf77b2d1fdacf9a8c8f7e571e580dd01317fc1247e4f86}; its node-excluding
attestor \pathcode{1851744} completed \pathcode{0:0} on \pathcode{h200-048-224} and published seal
SHA-256 \pathcode{6b57da131dcc81a89f45ddcf34ac46d3e6ec930a14088bb88a475b5da54160aa}.
Thus five of six external gate seals are accepted; only a scientifically acceptable replacement soak
lacks a seal.  No production training job is submitted; the state remains non-authorizing.

One production-launch integration risk is now explicit.  Final \pathcode{prepare-fresh} must run in
the namespace that observes the compute-authored numeric UID/GID data witness, while its frozen
source verifier also performs an authenticated \pathcode{git ls-remote}.  That Git subprocess
intentionally removes every proxy variable.  A read-only reproduction of its exact sanitized query
on the login side timed out after 35 seconds, whereas the earlier proxy-routed remote query passed.
Compute-side exact-verifier probe \pathcode{1851318} then failed \pathcode{1:0} after 25 seconds on
\pathcode{h200-094-088} at the sanitized authenticated \pathcode{git ls-remote}, with empty stdout.
Together with the login-side timeout, this closes the environment search: no tested namespace can
both reobserve the compute-owned data witness and satisfy the frozen network check.  The
qualification workloads remain valid evidence, but production launch now requires a reviewed source
correction and complete requalification; no verifier bypass is allowed.
A separate correction worktree developed a short-lived Ed25519 SSHSIG witness
that binds the authenticated remote ref to the exact commit, tree, source manifest, campaign,
aggregate, execution manifest, and launch mode, with an immutable one-use nonce claim.  Independent
review first found two pre-submission races: freshness after nonce burn and binding the held claim
directory to the same control-root inode through submission.  The first repair passed 64 focused
and 240 adjacent host tests, and an explicit-config, proxy-scoped read-only query against the exact
b37 branch succeeded in pueue task \pathcode{4531}.  A second independent review correctly rejected
that revision because public-repository \pathcode{git ls-remote} need not invoke the credential
helper, the validated GitHub configuration pathname remained mutable before child execution, and
freshness was not yet the last check at the actual \pathcode{sbatch} spawn boundary.  A direct
one-call GitHub GraphQL prototype in pueue task \pathcode{4533} did prove authenticated viewer
\pathcode{ngocbh} and returned the exact b37 branch, commit, and tree.  The current repair therefore
replaces \pathcode{ls-remote} with that authenticated query, supplies it from a private stable
snapshot of \pathcode{hosts.yml}, isolates GitHub CLI state/cache writes in the same transient
directory, and moves freshness into the final pre-spawn callback.  Final independent review returned
\pathcode{READY}: 69 focused and 286 focused-plus-adjacent tests pass, compilation and diff checks
are clean, and post-fix authenticated pueue task \pathcode{4540} returned the exact pushed branch,
commit, and tree without leaving a caller-directory state artifact.  The correction is committed
and pushed on branch \pathcode{codex/rdiag-shared-dt-v2-remote-witness-20260906} at commit
\pathcode{eea823e3dcdcb051ff256299a196a888f39b6ea0}, tree
\pathcode{e5a75ea5f2fa93442795f340bc9251fc35bb1e1a}.  A second reviewed commit then closed the
pre-freeze and launch-control gaps: canonical per-gate submit evidence, exact argv/comment/digest
rederivation, all-user exact-name scheduler checks, active-lease-bound v2 adoption, protected path
lineages, aggregate-output race closure, and planner/submitter parity.  That 12-file change is
committed and remotely verified at
\pathcode{c721817db9fbf8f7dff4c48850c15be34ef45549}, tree
\pathcode{c6cbf2620f6124ebf130a256befa6325ff001eaf}; 122 focused/bounded/remote tests and 147
legacy-adjacent tests pass, with one expected skip.  This is not a production authorization: the
new source still requires an evidence-supported stability correction and a completely fresh
qualification campaign rather than an amendment to b37 evidence.

The non-authorizing update-905 capture implementation is independently reviewed and frozen on
branch \pathcode{codex/rdiag-shared-dt-v2-soak905-diagnostic-20260906} at commit
\pathcode{b1b46f491f274a349728cf77255cce9a9028054e}, tree
\pathcode{263c9dab0e0fc4cccc9000abfba34a536801a779}.  Its four diagnostic-only files passed 22
focused and 166 focused-plus-adjacent host tests; production model, pretraining, launcher, and soak
gate files have zero diff from b37.  The immutable 4,294-file snapshot is
\pathcode{/lustre-storage/checkpoints/ngocbh/hdn/diagnostics/rdiag_shared_dt_v2/}\
\pathcode{diagnostic_source_snapshots/soak905_b1b46f491f274a349728cf77255cce9a9028054e},
with source-manifest SHA-256
\pathcode{c97683c8c23c980d4fd6bfa1a22d7c3629a8c38ecc8dbd6a99bcca50da99816d}.
It retains every boundary, the exact crossing microbatches and their pre-forward RNG states,
gradient categories and top parameters, a pre-update anchor, and eight committed successors while
preserving clipping at 1.0.  Two-node 16-H200 job \pathcode{1852006} completed \pathcode{0:0} in
20m16s on \pathcode{h200-130-169,h200-231-137}, with immutable-source verification on both nodes
before and after capture.  A field-by-field comparison of boundary, iteration, global pre-clip
norm, health-failure count, and alpha counts found diagnostic boundaries 1--904 bit-for-bit
identical to all 904 retained rows of failed production soak \pathcode{1851176}.  Boundary 905
then reproduced the finite pre-clip norm \pathcode{10.604082107543945}.  Clipping kept every
parameter and health observation finite, but the eight-successor window rules out a harmless
one-batch excursion: boundaries 908--913 crossed again and boundary 913 reached
\pathcode{2458.377197265625}; seven of boundaries 905--913 would fail the declared limit 10.  At
the first crossing, reconstructed category norms were 6.36818 for the rotational-angle weights,
4.34821 for DT weights, and 7.27711 for all other parameters, with the layer-9 angle projection the
largest retained parameter contribution.  This is reproducible escalating training instability,
not scheduler failure; whether its trigger is data-specific or v2-scale-specific is not established
by this run alone.

The completed non-authorizing closure contains 141 mode-0400 inventoried files totaling
15,555,765,453 bytes.  Its journal SHA-256 is
\pathcode{37572adf7e3594238ab87fde27318d2ca20f011dd14b7cc399e4678adec875aa}, its anchor-manifest
SHA-256 is \pathcode{e30b4432d0e8e1b53fa41b83048be173c29dec4e5b1c2b5d8dcd2ec0bc040abf}, its summary
SHA-256 is \pathcode{6e882b60495cfd1131365a280022156a43e8faf178fafe68fc481eca27ab1a41}, and its closed
inventory SHA-256 is \pathcode{e3276e7489e0798509077e87c4822b89e260ef372a08efbfa1d7c727575faae3}.
The summary is explicitly \pathcode{authorizing:false} with status
\pathcode{completed_crossing_window}; it cannot issue a production result or seal.  Exact
anchor-restored \pathcode{v2-A1}, \pathcode{v1-scale}, and \pathcode{v2-A2} replay was therefore
run to isolate the effect of the rotational scalar.

That replay is now implemented, independently reviewed, committed, and pushed on the same
diagnostic branch at commit \pathcode{016a5ab9018774fae6ddd303c946b832e2a188a6}, tree
\pathcode{01b74f1ad623c3c58762e5f9f6066d2acbc4c638}.  Relative to capture commit
\pathcode{b1b46f491f274a349728cf77255cce9a9028054e}, it adds exactly a replay plan, helper, isolated
Slurm runner, and host test; no production source changes.  The 33 combined capture/replay host
tests pass, and independent review returned \pathcode{READY} after an initial capture-only writer
schema mismatch was repaired.  Its immutable 4,298-file snapshot is
\pathcode{/lustre-storage/checkpoints/ngocbh/hdn/diagnostics/rdiag_shared_dt_v2/}\
\pathcode{diagnostic_source_snapshots/soak905_replay_016a5ab9018774fae6ddd303c946b832e2a188a6},
with source-manifest SHA-256
\pathcode{0abe22cf6c395902dcbb41bedff97486d9d36e6babd6366a3d03e400498db6b4}.  External-bootstrap
source verification passed, as did all 31 snapshot-applicable tests; the two deselected tests
require a live Git worktree.  Two-node, 16-H200 replay job \pathcode{1852750} completed
\pathcode{0:0} in 5m22s on \pathcode{h200-153-168,h200-252-145}, with both capture nodes excluded.
Both v2 references were bit-exact to each other and to every rank of the captured boundary-905
backward, and no optimizer step was executed.  At effective angle scale 0.125, the reconstructed
global norm was \pathcode{10.604082390132142}; changing only to the older 0.25 scale raised it to
\pathcode{12.423034238808185}, a factor of 1.17153.  The corresponding angle, decay, DT-bias,
DT-weight, and other category norms changed from 6.36818, 0.08549, 0.14930, 4.34821, and 7.27711
to 6.55957, 0.10176, 0.21774, 5.37835, and 9.07301.  Thus the larger v1 scale causally worsens this
boundary, but the present v2 scale still violates the production limit and is not acceptable.
The non-authorizing summary, validation, and closed-inventory SHA-256 values are respectively
\pathcode{82c22cee330756b2527bbc77e6fa39a2ad40f417042ec548a16039a2cc8914ed},
\pathcode{1a575a7c565bbae0896402be7e6d9affb70040f9f0d6bc8de0782b4d538dd7ee}, and
\pathcode{5ae2c8adaff6d531c4a74919235c9402d46ca46b78d60f8c3e13a7e7da892ce5}.

The next non-authorizing discriminator is implemented, reviewed, committed, and pushed, but is not
a source correction.  It is isolated in worktree
\pathcode{rdiag_shared_dt_v2_stability_screen_20260906} on branch
\pathcode{codex/rdiag-shared-dt-v2-stability-screen-20260906}, based exactly on replay commit
\pathcode{016a5ab9018774fae6ddd303c946b832e2a188a6}.  Its four-file diagnostic-only delta is frozen
at commit \pathcode{470d3af6126ea45181c5e8b6e4bb1caa98e1ec34}, tree
\pathcode{e4b201f25748d9d81a5647c1d1f096d1f33ca22b}; independent review returned \pathcode{READY}
with no Critical or Important findings.  The final adjacent suite passed 43 tests with one known
base-inherent deselection.  The immutable 4,302-file source snapshot is
\pathcode{diagnostic_source_snapshots/stability_screen_470d3af6126ea45181c5e8b6e4bb1caa98e1ec34},
with source-manifest SHA-256
\pathcode{377efb243daf0fc86610654009a75d998833067f62b2926670f4e56b4b86546d}; an independent
commit-derived verifier and its 42 snapshot-applicable tests passed, with only two tests requiring
a live Git worktree deselected.  All 16 loader snapshots and 144 retained pre-forward RNG states
reopen successfully.  After an all-user exact-name duplicate check returned no rows, two-node,
16-H200 job \pathcode{1853781} was submitted once at 12:28 UTC under \pathcode{h200_dev}, with
both capture hosts excluded independently in the submission and frozen runner.  It started on
\pathcode{h200-046-153,h200-094-010} and failed \pathcode{15:0} after 8m05s.  All 16 ranks first
reproduced boundary 905 exactly at reconstructed global norm
\pathcode{10.604082390132142}, but that boundary uses the retained tensor bundle.  On the first
live successor fetch, boundary 906, global rank 1 reported that DataLoader child PID 815552 died
with \pathcode{SIGSEGV}; \pathcode{--kill-on-bad-exit=1} then terminated the other ranks.  Each A1
journal therefore contains exactly one valid boundary-905 row, while no boundary-906 row, rank
result, summary, output validation, closed inventory, or terminal marker exists.  Job 1853781 is
partial non-authorizing infrastructure-failure evidence and cannot select an intervention.

The failure is localized to the deep four-worker
\pathcode{StatefulDataLoader}/HF-Parquet state restore, not teardown, FSDP, the captured RNG, or
the boundary-905 computation.  The capture used one fresh loader continuously and the exact replay
used retained tensors with no loader; only the screen called \pathcode{load\_state\_dict} on the
saved 1,810-yield multiworker cursor.  The production source itself documents that this restore
path can segfault a worker.  Retrying the same code or merely excluding the host is therefore
rejected.

The repaired diagnostic is complete, pushed, and independently accepted at commit
\pathcode{c0aa93389e4d28ee9021bf3e0dff573817d39afe}, tree
\pathcode{bd2d32c32e65e7cc7e059fdd869835538a3589ae}, in the same worktree and branch.  Its focused
suite passed 19 tests; the combined scoped/adjacent suite passed 51 tests with one historical
cumulative-branch-only test deselected, and both specification and quality review reported no
findings.  Before CUDA/NCCL, all 16 ranks now construct fresh pinned four-worker loaders, discard
the first 1,808 batches, materialize batches 1,808--2,415 (boundaries 905--1208), require all 18
batches for boundaries 905--913 to match the retained capture byte-for-byte, synchronously close
every worker, and seal one rank-local read-only tape.  The FSDP phase independently rehashes that
tape and creates no loader.  Retained RNG is restored through boundary 913, while every later shared
arm boundary carries cross-arm batch and RNG witnesses.  Exact rank/node mapping, isolation on
entry and exit, cache separation, cleanup-before-completion, and the 126-file/125-record terminal
closure are all fail-closed.  Its immutable 4,302-file snapshot is
\pathcode{diagnostic_source_snapshots/stability_screen_c0aa93389e4d28ee9021bf3e0dff573817d39afe},
with source-manifest SHA-256
\pathcode{fd03d23ae4b4650f11ee9db194af937ed5088b49fa154bb2e7fc2fe18f177b7a}; a separate
commit-derived private rematerialization verified every file and mode.  The first probe allocation,
\pathcode{1854865}, was cancelled before start after its unnecessary GPU reservation was noticed;
it produced no artifacts.  CPU-only H200-node replacement \pathcode{1855110} started on
\pathcode{h200-060-199}, constructed the fresh four-worker stream, skipped 1,808 batches, and fetched
the 18-batch comparison prefix, but failed closed during worker teardown before executing the
retained-byte comparison.  Torchdata 0.11 waits five seconds per worker, sends \pathcode{SIGTERM} to
survivors, and does not reap them afterward.  Two PyArrow workers emitted ``terminate called
without an active exception'' and remained in kernel \pathcode{do\_coredump}; because this cluster
pipes dumps to \pathcode{systemd-coredump}, the job was killed after 4m59 to release the node.  It has
only the capture-validation artifact and cannot pass the probe.  The first reviewed cleanup repair
was committed and pushed as
\pathcode{b5bfef19bf00f8a6602addc03b152cacdd327337}, tree
\pathcode{02207725432936cda53370311300d744342a3f86}.  Its immutable 4,302-file snapshot is
\pathcode{diagnostic_source_snapshots/stability_screen_b5bfef19bf00f8a6602addc03b152cacdd327337},
with source-manifest SHA-256
\pathcode{3f4e661d94496f5da4f38554bfd34b3792bc476149a3a78bf8e5552e2f1700fc}; an independent
private rematerialization was byte-identical and 25 snapshot-applicable tests passed.  Probe
\pathcode{1855634} then failed in one second before creating a run root because its ad-hoc
\pathcode{sbatch --wrap} used Bash \pathcode{pipefail} under \pathcode{/bin/sh}.  Corrected wrapper
probe \pathcode{1855635} ran on \pathcode{h200-128-202}, rebuilt the four-worker stream, and
reached teardown, but exposed an earlier failure than the timeout fallback: torchdata first sends a
queue sentinel, and unwinding the live PyArrow parquet readers made three workers abort before
\pathcode{terminate()} was reached.  The probe failed closed in 1m57s, retained only the exact
capture-validation document (SHA-256
\pathcode{116b4cb5adf6cb20f2f6181ae118cdf701a8437caaf68557b35cd64952d28bb2}), and produced no
probe result or marker.  The final repair removes the iterator from PyTorch's SIGCHLD registry while
all workers are live, sends SIGKILL to and reaps all four, and only then lets the parent
cancel/close its index and result queues directly.  It never invokes torchdata's graceful routine
or touches the process-shared done event after a forced kill, avoiding the documented dead-child
event-lock hazard.  Independent review returned READY after 27 focused tests, a 59-test scoped
suite (one historical cumulative-branch test deselected), registered-child stress, and explicit
ineffective-kill fallback.  The reviewed source is pushed at commit
\pathcode{7ea1735b3fed99b6a24e7ac07753771defdec944}, tree
\pathcode{2a9b111657f0d8b68a6839fbeea7a2c2a593e6b6}.  Its immutable 4,302-file snapshot is
\pathcode{diagnostic_source_snapshots/stability_screen_7ea1735b3fed99b6a24e7ac07753771defdec944},
with source-manifest SHA-256
\pathcode{4a5caf3db17da8be95660f78074f972765df624f49a8c1dd31b53c74fe8eb769}; private
rematerialization was byte-identical and 26 snapshot-applicable tests passed.  Fresh CPU-only H200
probe \pathcode{1855869} completed \pathcode{0:0} in 1m49s on \pathcode{h200-126-216}.  Its 18
retained batches matched, CUDA/distributed remained uninitialized, and all four worker PIDs were
reaped with exit code $-9$.  Its probe-record SHA-256 is
\pathcode{0a481a5313de98c2f48715d69b4f887997b2e1335751e4de804835e9df5cebe8}; stderr is empty.
After an all-user exact-name conflict check, exactly one full two-node/16-H200 non-authorizing
screen, job \pathcode{1855985}, was submitted from this snapshot.  It completed
\pathcode{0:0} in 16m54s on \pathcode{h200-017-170,h200-046-153}; every root and step row is
\pathcode{COMPLETED}.  References A1 and A2 are exact to each other and to retained capture
boundaries 905--913, reaching the same maximum norm \pathcode{2458.377197265625}.  The
quarter-head arm ($0.25/\sqrt{H}$, effective scale $0.0625$) completed boundaries 905--1200 with
zero crossings, maximum \pathcode{2.5316383838653564}, and nearest-rank p99
\pathcode{0.5541900992393494}; it is the sole direct-promotion-eligible arm.  Connected-zero
rotation was also stable through 1200 (maximum \pathcode{0.4557827115058899}, p99
\pathcode{0.3324354290962219}) but remains mechanistic-only by declaration.  Half-LR crossed again
at boundary 911, reached \pathcode{28.53986930847168}, and closed its successor window at 919, so
it is rejected.  The read-only closure has exactly 126 files and 125 inventory rows.  Summary,
validation, and inventory SHA-256 values are respectively
\pathcode{19f0962e04cc45b5df8d9df6e39025cecf54b48423d6e0290cd60a6180d95ec5},
\pathcode{4683c0609bccbde088b9c41ad1ea32f509346826fbc255ce0c87229048699776}, and
\pathcode{4c9180ab8b49d5b206a5dabe56010fcfd09ac5d2f3fb97ff12270ca4bf3ebe83}.
Independent re-execution of the immutable source, output, and inventory validators passed; all 125
inventoried payloads rehashed exactly, the terminal marker occurs once, and cleanup precedes that
marker.  The retained 98,660-byte screen stderr contains only FSDP/NCCL deprecation and process-
group teardown warnings, with no traceback or error.  Cleanup proof is the fail-closed runner
ordering plus completed scheduler steps and marker, rather than a standalone per-node certificate.
The screen remains explicitly \pathcode{authorizing=false}: it selects quarter-head for a new
production-candidate source and fresh soak, not for direct launch.

This failure also exposes a separate production-readiness gap.  The intended 50B job is budgeted
for 70 training hours inside a 72-hour allocation and, at the measured 31.08K tokens/s/GPU on 16
GPUs, nominally needs about 27.9 hours.  Nevertheless the launcher explicitly enables requeue, the
production loader uses four persistent workers, and the first ordinary checkpoint cursor is iter
2,000, deeper than the failed 1,810-yield cursor.  The accepted resume job \pathcode{1851174}
covered only a 220M, two-GPU, one-worker restore from iter 16 to iter 24; it does not test this
failure class.  The bounded control plane forbids \pathcode{HDN\_SKIP\_DATA\_RESUME}, and a worker
SIGSEGV becomes terminal rather than safely requeued.  Production therefore remains blocked until
the loader-resume mechanism is fixed or replaced and the mandatory resume gate is upgraded to the
real two-node/16-rank, four-worker HF-Parquet path at iter at least 2,000, with byte-exact successor
batches against an uninterrupted control and successful post-restore optimizer progress.  The
diagnostic batch tape fixes D1 only; it is not a production-resume correction.  The selected
production design is instead a snapshot-free deterministic replay cursor: set
\pathcode{snapshot\_every\_n\_steps=0}, prohibit managed-production loader
\pathcode{state\_dict}/\pathcode{load\_state\_dict}, checkpoint the yielded-batch count and a rolling
SHA-256 chain over canonical rank-local raw batches plus the exact loader/data/runtime identities,
and on resume fresh-build the same four-worker stream, replay and rehash the saved prefix, restore
the saved training RNG, and continue from that iterator.  Opaque v1 sidecars must fail closed rather
than migrate silently.  The replacement gate will use the exact 2-by-8, 750M, four-worker geometry,
checkpoint at iter 4,000/update 2,000, compare the next two batches and the resulting update against
an uninterrupted control in a separate allocation, deliberately fail if either loader restore API
is called, and require a projected full-horizon replay time of at most eight hours.  A sealed
rank-local token tape remains the fallback only if deterministic replay misses that throughput gate.

Implementation is isolated from both the diagnostic branch and dirty main checkout in worktree
\pathcode{/lustre-storage/checkpoints/ngocbh/hdn/diagnostics/rdiag_shared_dt_v2/source_worktrees/rdiag_shared_dt_v2_replay_cursor_20260906},
branch \pathcode{codex/rdiag-shared-dt-v2-replay-cursor-20260906}, based on reviewed launch-control
tip \pathcode{c721817db9fbf8f7dff4c48850c15be34ef45549}.  The cursor/data core is committed and
pushed as \pathcode{4a2002b}; v2-only trainer integration is committed and pushed as
\pathcode{cc18784afe16ee760af596de3c24461e283a34e5}, tree
\pathcode{87181e1d47edbfc27ad294bbeb659aeb8c82d1c8}.  It uses one iterator for prefix replay and live
continuation, validated cursor/RNG sidecars, no v2 loader state API, coordinated asymmetric failure
handling, and four-worker hard-stop cleanup before any terminal checkpoint or PASS evidence.
Independent review returned READY after 573 tests (three skips), a two-rank asymmetric-collective
probe, and a real four-child cleanup test.  A separate reviewed hardening commit
\pathcode{7b5808a177d03dd1526d1e0c2dc7a83a0f816217}, tree
\pathcode{ba0a884d48aedba42ce673bc1c7a0c089a316401}, installs an opt-in, spawn-picklable worker
initializer that verifies \pathcode{RLIMIT_CORE=(0,0)} and Linux dumpability zero before any
PyArrow read; 158 focused/adjacent tests and a real subprocess check pass, while legacy loaders are
unchanged.  This is a reviewed implementation slice, not a launch
authorization: the v2 control-plane/evidence schemas, uninterrupted-control gate, compute-node
StatefulDataLoader/PyArrow validation, and final executable plan remain open.  Its immutable
4,303-file snapshot is
\pathcode{diagnostic_source_snapshots/replay_cursor_7b5808a177d03dd1526d1e0c2dc7a83a0f816217},
with source-manifest SHA-256
\pathcode{e657f2617e98b25b4c413120bd92b3bf59b0d636e1888f3ec85cc63a449bf445}; private
commit-derived rematerialization and 101 snapshot-local tests passed.

The implementation audit also identified work that is still genuinely prospective at this tip.
The RNG payload currently contains Python tuples and a NumPy array, and its sidecar reader still
uses \pathcode{torch.load(...,weights_only=False)}; the current v2 resume evidence class also
retains the old pending-loader-state attestation.  The final implementation must replace these with
the specified primitive/tensor-only RNG representation, restricted weights-only loading, and the
new replay evidence object before the source can be frozen as a production candidate.

A separately reviewed, non-authorizing compute probe then exercised the actual production
StatefulDataLoader/HF-Parquet path.  CPU-only H200 job \pathcode{1856719} completed
\pathcode{0:0} in 3m19s on the previously unused node \pathcode{h200-171-121}.  Two wholly fresh
logical-rank-1-of-16 loaders, each with four persistent workers, consumed 1,826 batches and agreed
on the boundary-1,808 cursor digest and every one of the following 18 cursor digests.  All eight
distinct workers reported both core limit zero and \pathcode{CoreDumping: 0}, were reaped at exit
code $-9$, and left no visible \pathcode{/proc} entry; CUDA and distributed remained uninitialized,
post-run data/tokenizer identity validation passed, and stderr is empty.  The slower transcript took
58,317,465,752 ns, giving the integer-ceiling full-190,734-batch projection
6,091,524,377,187 ns, or 1.6921 hours, below the proposed eight-hour gate.  The mode-0400 result has
SHA-256 \pathcode{ab308edc368c7920cd8f338efa200819609a1e4d1887ba46b20931500dd3502c} and
remains explicitly \pathcode{authorizes_production=false}; it validates one-rank loader
determinism, throughput, and teardown, not the required two-node training/resume equivalence.
Independent post-run audit returned \pathcode{VALID} after rehashing all 4,303 snapshot files,
the scheduler rows and artifact closure, both transcript chains, all worker records, and the
integer projection.

Plan revisions v1--v3 are retained in the worktree.  Revision v4 was also rejected by independent
review: it selected the invalid login-namespace data witness for compute work, put full-data
attestation on the login node, omitted node-local source/cache setup for the second node, had
contradictory timing cardinalities, tried to construct all four family manifests despite singleton
750M authorization, under-specified canonical recipe and detached-state mechanics, and merely
printed rather than machine-asserted the final Slurm lifecycle.  Revision v5 then exposed remaining
compute-namespace, staged-publication, descriptor-bound loading, and cleanup-handshake gaps.
Revision v6 is complete at
\pathcode{docs/plans/2026-09-06-rdiag-shared-dt-v2-replay-cursor-plan-v6.md}, SHA-256
\pathcode{9f4fec8f3c5e208a0b33af05656283975c0c44dc103c56f8f83040391e675094}.  Independent
review returned \pathcode{NEEDS\_REVISION}: its nominally allowed U07a--U45d tranche still rewired
active v2 sidecars, evidence, runners, and seals; the v2 retirement was prose rather than an
executable guard; dependencies placed publication tests before staged APIs; final verification was
blocked; and the prepared-final handle, bounded cleanup acknowledgement, signal conversion, and
mapped-reader lifetime were incomplete.  No v6 implementation began.  Revision v7 was then frozen
at SHA-256
\pathcode{2c2851f6e2db2ae4a84aaee0e09dccf965894156fd4b5c6a3913b14b4f7fccbe}, but independent
review again returned \pathcode{NEEDS\_REVISION}: retirement missed qualification preparers and
occurred after some CLI side effects; it bound the legacy final-model path instead of the actual
\pathcode{rdiag-final-committed.json} pointer; post-rename failures were modeled as rollbackable;
cleanup and detached publication lacked exact collective APIs; and several contracts/dependencies
remained underspecified.  No v7 implementation began.  Revision v8 was frozen at SHA-256
\pathcode{dcf4ea6294c901b3d4dbcb2a6362af303274b1ff034c2f134c278de95c4b0105}.  It expands
retirement to 25 supported v2 construction, publication, attestation, witness-generation, and
submission boundaries; preserves historical read-only validation; specifies separate stable
directory and closed-file identities; uses the actual
\pathcode{rdiag-final-committed.json} pointer; models irreversible and uncertain commit outcomes;
and defines family-neutral restricted-load, Merkle-evidence, collective cleanup, and staged
publication APIs.  Independent review nevertheless returned \pathcode{NEEDS\_REVISION}: guarded
SubmitLine builders were still called by historical read-only validators; detached publication
could not represent post-hardlink failure; its staged-handle digest was circular and unverifiable
after publication; detached/final rank transport lacked exact primitive codecs; attestor edit
rules contradicted the terminal audit; two v2 evidence/authorization paths remained live; and the
cleanup contract disagreed on one versus two join passes.  No v8 implementation began; revision v9
was frozen at SHA-256
\pathcode{5eb6bb60cc06ee415241b9e7c66814cfa9076b09461d0abce3b4396a8510d92c}, and revision v10
was frozen at SHA-256
\pathcode{ef6e1bc095ca40182d146fa411b7fafafbf3a2b64422c825ba93d193c8dda152}.
Both are immutable superseded planning evidence; neither produced code, a compute job, or a
production authorization.  Revision v11 is frozen read-only at SHA-256
\pathcode{b0e681bd0fe01559ae8f8ad5797b2b5ebf21a3804fd7a639fbc392132b82f0f8}.
Independent review returned \pathcode{NEEDS\_REVISION}: the generic manifest guard names the
execution-manifest schema although its input is a recipe; the nominal 33-boundary audit scans more
host modules than its literal disposition inventory; the lifecycle cannot supply the worker and
queue identities required by its cleanup-certificate validator; it publishes the split before
continuing the same iterator; the reconstructed-SubmitLine runtime type check is outside its
declared edit surface; the evaluation-seal test path is wrong; and lifecycle/dependency and
cross-result bindings remain incomplete.  Revisions v12--v18 were subsequently frozen and each
rejected by independent review; their SHA-256 values are, in order,
\pathcode{5dbf3d58321ea72da34bc762038ea2878e93bec457982a92c7cc6956163f7cfb},
\pathcode{82845d70a67ace5152b47252cfe35400af0af32ffad9d90354d646640805b790},
\pathcode{648dcdbb54e78d07ba6e2fc57c8d0bcd4914f59122242eb3ae9b8baf4cd5cbe5},
\pathcode{1ce145b2e9f74400efe26e8caa428819b4c6818982a4ba483b51708f9910983a},
\pathcode{caa082f5c15c6caae00e97adb1ede1c024be456725b9d8cae8156eab89284a27},
\pathcode{c35c1e824dc6d007931cd6c2c4e2a0c4c6a69c5c39609fd026ef239e574fa265},
and \pathcode{feac1b553db2842c4d8b4e032cfd13e5f96ac676e9f2f8806e89dc42fee4659f}.
Revision v19 is frozen at SHA-256
\pathcode{41e0ca7c0f3f91c79929724a5fe8cad2b1f01a94ee7835d87ca1a52a811d9467}
and remains under independent review; it is not yet accepted.

The T00 fail-closed retirement tranche is committed locally as
\pathcode{9d92d77941e877bf0b99d55df0f051929bb2bcff}, tree
\pathcode{b5c0e19c75a8813df73163e85a07c8f0c87036fd}, and received an independent
\pathcode{APPROVED}.  It guards exactly 33 supported v2 boundaries and pins 463 literal audit
fixtures.  Integrated pueue task \pathcode{4607} passed 653 tests with four skips, compared with
baseline task \pathcode{4594} at 525 passes and four skips.  This local T00 commit has not been
pushed, submitted as a compute job, or used to grant production authorization; no later replay
substrate tranche is implemented by that result.

The first additive \pathcode{shared\_dt\_v3} quarter-head plan was reviewed at SHA-256
\pathcode{aee51f621f5e66aa15b94a369b2cf61b3a79f4b28bcdc8d3f1769ee399ba398e} and rejected by
independent review.  Its ignored worktree file was subsequently edited in place during v8 planning,
contrary to the declared freeze, and now has SHA-256
\pathcode{8a429606e563189e7700ea0ae906403f68e247a40c57ce38f8ec19e4fa47af3c}; the original bytes
could not be recovered and neither version is executable evidence.  The reviewed version retained
a stale eleven-boundary retirement check, disagreed with v8's
sidecar API/schema, omitted the required 48-rank-timing authorization calculation, placed the final
pointer directly under the output root instead of the manifest's \pathcode{out\_dir}, allowed a
TERM-before-KILL cleanup path, omitted exact attestor provenance and a v3 nonce-claim contract, and
contained dependency inversions.  Revision v2 was frozen read-only at SHA-256
\pathcode{83a2c07f502824ac982e4efb0c43141f2571aa60e82c7fc0e567448f789f70e3}, but remains a
non-executable review response pending an accepted replay-substrate successor and further review.
The eventual accepted revision must preserve v1/v2 identities, add exactly the singleton
\pathcode{rdiag\_kla\_shared\_dt\_v3\_750M}, and own every new model/schema/runner/seal binding
before any production operation.

The completed matched screen crossed the exact warmup peak at boundary 954 with one model resident
at a time and satisfied every declared comparability invariant before interpretation.  Under the
predeclared policy, quarter-head is therefore selected over half-LR as the production-source
candidate; the radial-only result strengthens the causal attribution to the rotational path but is
not itself eligible.  This selection is still only a design decision.  The new versioned clock must
pass a fresh generation-zero 3,300-update soak and disjoint-node attestation; the diagnostic cannot
replace that qualification gate.

The remaining sequence is:
\begin{enumerate}[leftmargin=1.8em]
  \item preserve failed infrastructure run \pathcode{1853781}, failed cleanup probes
        \pathcode{1855110}/\pathcode{1855635}, and pre-Python wrapper failure \pathcode{1855634};
        preserve accepted non-authorizing screen \pathcode{1855985}, and implement its sole
        promotion-eligible quarter-head correction in a new versioned source without weakening or
        blindly retrying the gate;
  \item complete independent review and acceptance of frozen replay/retirement plan v19; only then
        implement its dormant substrate beyond the already approved local T00 retirement at
        \pathcode{9d92d77941e877bf0b99d55df0f051929bb2bcff}, then run the replacement authenticated
        gate at exact production topology through iter 4,000, proving byte-exact
        successor batches, a committed post-resume update against an uninterrupted control, and
        the eight-hour maximum projected full-horizon replay time;
  \item combine only an evidence-supported soak and resume correction with reviewed launch-control tip
        \pathcode{c721817db9fbf8f7dff4c48850c15be34ef45549}, freeze one
        final production candidate, and do not aggregate the superseded b37 seals;
  \item reuse only the recursively revalidated compute-authored data manifest, then repeat all six
        preparations, workloads, and disjoint-node attestors for the final source, run the independent
        derived-gates control, build a fresh six-seal aggregate, and independently reopen and
        recursively validate that aggregate while keeping every existing v1 canonical byte sequence
        and pin unchanged; and
  \item generate a short-lived signed remote-ref witness for that exact aggregate and execution
        manifest, create the separate required production snapshot, and invoke the \pathcode{prepare-fresh}
        mode of a new exact-family v3 submitter; the existing v1 and retired v2 submitters must not
        be reused.  Submit only the additive singleton
        \pathcode{rdiag_kla_shared_dt_v3_750M} with train configuration
        \pathcode{tsz128x4k_50B}; retain the supervisor-enforced \pathcode{standard},
        \pathcode{niah_single_1}, \pathcode{niah_single_1_8k}, \pathcode{niah_single_2},
        \pathcode{niah_single_2_8k}, \pathcode{niah_single_3}, and
        \pathcode{niah_multikey_1} evaluation modes.
\end{enumerate}
Job \pathcode{1835824}'s scientific failure invalidated the prior nine-workload lower bound;
jobs \pathcode{1836996}, \pathcode{1837446}, \pathcode{1838750}, \pathcode{1839531}, and
\pathcode{1847039} completed the localization, replacement generation-zero, trigger-capture,
matched-intervention, and first source-matched replacement-candidate workloads; S7 job
\pathcode{1847784} has now completed their strictly repeated six-arm diagnostic.
Jobs \pathcode{1846561} and \pathcode{1847039} returned provisional valid passes, but validators
\pathcode{1846728} and \pathcode{1847464} respectively failed before validation and after
scientific validation but before evidence sealing; neither cycle is externally accepted.  From
this point, S7 sealing and the v2 production-H200 gate are complete, and the matched screen has
selected quarter-head as the correction candidate.  It still requires implementation, a new
acceptable soak, and then a complete fresh-source campaign before production launch; all
four short b37 qualification workloads and their external attestors have already completed.
Production training would follow the new campaign.  Because data identities include
numeric UID/GID, the final-source hard successful-path minimum is 21 compute allocations before
production: six gate preparers, six scientific workloads, six fresh external attestors, one
independent prepared/derived-gates control, one aggregate builder, and one independent recursive
aggregate validator.  The existing compute-authored data manifest is reusable only after full
compute-side verification folded into the first H200 preparation; regeneration would add an
allocation.  A compute-side guarded-submit controller and production itself add two more, while
retaining the earlier optional third-node H200-seal control adds one.  On the historical b37 path,
seventeen of its original 21 outcomes have completed: the manifest builder,
all six preparers, the H200 workload and attestor, and the train-smoke and checkpoint-fault
workloads and their attestors, plus the gradient-A/B workload and attestor and the resume workload.
Its attestor is also accepted.  The failed soak still leaves four original hard-minimum outcomes: its
successful replacement and attestor, the aggregate builder, and the aggregate validator.  At least one new soak preparation is
additional to that lower bound.
Including the new soak preparation and preserving the earlier independent pre-submit-control
convention adds two allocations, leaving at least six outcomes before complete b37 scientific
closure; retries and diagnostic allocations are additional.  More importantly, the required
remote-ref correction changes source identity, so the b37 seals cannot authorize that source: after
review, all six preparations, all six workloads, all six attestors, the prepared-gates control,
aggregate construction, and independent aggregate validation must be rerun in a fresh campaign.
The compute-authored data
manifest remains reusable subject to recursive revalidation.  The six
immutable runners contain 13 distinct \pathcode{srun} call sites and execute 15 steps on a normal
success path because the soak reuses its source-copy and cache actions during cleanup.  One separate
production \pathcode{sbatch} follows only after aggregate acceptance.  A held-out replication would
increase that count.
Failed attempts \pathcode{1838085} and
\pathcode{1838250} did not satisfy the trigger-capture stage.  This count excludes retries, additional
discriminator arms, implementation-time GPU checks, held-out causal replication, and production
requeues.
Until that sequence closes, neither a production authorization nor a replacement training result
is claimed.

\appendix
\section{Reproducibility audit}
\label{app:sources}

\paragraph{RDiag implementation baseline.}
The sealed E2 baseline implementation statements in this note were audited against branch
\texttt{rdiag} at Git commit
\texttt{c3346feaecae1486cdbc71f82b4edeacc4198950} (tree
\texttt{88e423d2926d8ae81945c08393a59861d285907c}).  This pin identifies the sealed implementation
snapshot; the present documentation-only working-tree update is intentionally outside that
snapshot.  Relevant regions are:
\begin{itemize}[leftmargin=1.8em]
  \item \pathcode{lit_gpt/rdiag_kla.py}, construction and buffers around lines 137--193, phase
        construction around 222--293, and the fused layer route around 409--516;
  \item \pathcode{lit_gpt/kla_ops/rdiag_clock.py}, anchor and eager oracle around 72--119 and
        190--270, exact fused clock around 283--390, and native-order custom backward beginning
        around line 515;
  \item \pathcode{lit_gpt/kla_ops/rdiag_direct.py}, the unchanged direct-pair forward around line
        162, guarded direct-pair backward around line 338, precise clock materializer around line
        300, and radial wrapper around line 920;
  \item \pathcode{lit_gpt/kla_ops/rdiag_phase.py}, negative scan contract and reference around
        1--19 and 294--328;
  \item \pathcode{lit_gpt/kla_ops/rdiag_bounded_phase.py}, bounded fail-visible phase contract
        around 1--11 and 41--172;
  \item \pathcode{lit_gpt/kla_ops/rdiag_norm.py}, public fused normalization/rotation contract
        around 10--90;
  \item \pathcode{lit_gpt/kla_ops/rdiag_norm_reference.py}, normalization, BF16 boundary, and
        adjacent-pair rotation around 24--96;
  \item \pathcode{lit_gpt/kla_ops/diag_naive.py}, KLA equations around 14--32, 144--199, and
        451--491;
  \item \pathcode{lit_gpt/diag_kla.py}, process-noise and isotropic-output initialization around
        280--300;
  \item \pathcode{lit_gpt/config.py}, bounded RDiag recipe fields around 349--369 and 742--763;
  \item \pathcode{lit_gpt/rdiag_training_health.py}, health thresholds and observed counters around
        24--37 and 121--255, and fail-closed decision logic around 579--612;
  \item \pathcode{lit_gpt/fsdp_strategy.py}, the FP32-buffer mixed-precision override beginning
        around line 131;
  \item \pathcode{scripts/tests/rdiag/test_rdiag_bounded_h200.py}, exact radial, actual-logit,
        exceptional-value, full-layer parity, and A/B performance gates;
  \item \pathcode{sanboxes/kernels/rotational_diag_kla_v1/DESIGN.md}, moving-frame semantics and
        the pair-tied literal model; and
  \item \pathcode{sanboxes/kernels/rotational_diag_kla_v1/RESULTS.md}, the correctness and v1
        performance evidence around lines 128--166.
\end{itemize}
The whole tree is sealed by the snapshot manifest identified below, which is the authoritative
byte-level inventory.  Three central files in the final parity implementation and gate have
SHA-256 values
\begin{itemize}[leftmargin=1.8em]
  \item \pathcode{lit_gpt/kla_ops/rdiag_clock.py},\\[-2pt]
        {\ttfamily\scriptsize d02a3f0d0994067afc174f1a77f7ebd8fcecd73c3a18956e72ab79ab49afc966};
  \item \pathcode{lit_gpt/kla_ops/rdiag_direct.py},\\[-2pt]
        {\ttfamily\scriptsize 9fa3e29102bfa512ecbc57fc04ce6453e3a33e7aed61e7049026b71fe81faace};
  \item \pathcode{scripts/tests/rdiag/test_rdiag_bounded_h200.py},\\[-2pt]
        {\ttfamily\scriptsize 68d7db9b6793a33ef2574c7a39527306c3b3b500a03e1807790ce59889be8517}.
\end{itemize}

The final result followed three deliberately narrow diagnoses.  Commit
\texttt{fe1b1258} made the radial clock mean bitwise exact but left
three full-layer angle-weight mismatches in job 1777082.  Job 1777229 then proved that the first
remaining forward divergence occurred at sigmoid output despite a bitwise-identical input logit.
Commit \texttt{4ce6e19} matched the CUDA-eager sigmoid and angle VJP, and commit
\texttt{c3346fe} added the missing-delta compile-time guard that makes the detached radial path
immune to nonfinite discarded angles.

\paragraph{Installed Mamba implementation.}
The source audit used \texttt{mamba-ssm 2.3.2.post1} installed under the active \texttt{hdn}
environment; those bytes match upstream tag \pathcode{v2.3.2.post1} at commit
\pathcode{a14b1dff0454a3bc27d9eb31355dc01e4b2490ec}.  Paths and SHA-256 identities are:
\begin{itemize}[leftmargin=1.8em]
  \item \pathcode{mamba_ssm/modules/mamba2.py},\\[-2pt]
        {\ttfamily\footnotesize 605e4439ff0baec8d8acaf4a191d9f0570eea9900065a065909124c472b08707};
        inspect roughly lines 95--136, 210--260, and 309--318;
  \item \pathcode{mamba_ssm/ops/triton/ssd_combined.py},\\[-2pt]
        {\ttfamily\footnotesize 39c19c1e4c8e36982847079bc12b4f07ed3af03ac68f6a6653200ef80df5515d};
        this is the fused Mamba-2 SSD scan;
  \item \pathcode{mamba_ssm/ops/triton/ssd_chunk_state.py},\\[-2pt]
        {\ttfamily\footnotesize 6d82771b2f62a7bf84c381c8ef8e4b579a95e0352021d117507a1072e4d71dcd};
        inspect roughly lines 40--86 and 102--174;
  \item \pathcode{mamba_ssm/modules/ssd_minimal.py},\\[-2pt]
        {\ttfamily\footnotesize 6442a9fe25a6792258df657392ab1fab4dd873ee777d9dd5783cc76030218e07};
        this is the readable Mamba-2 SSD reference;
  \item \pathcode{mamba_ssm/modules/mamba3.py},\\[-2pt]
        {\ttfamily\footnotesize de1b104865c941eb6eeac63c3a75d5913fc2e7a119d52991590ef19c381b05ac};
        inspect roughly lines 76--87, 150--178, and 220--237;
  \item \pathcode{mamba_ssm/ops/triton/mamba3/angle_dt.py},\\[-2pt]
        {\ttfamily\footnotesize 7d99f8cff37371cfbfe9d4dd36f17a16135fb07f2c6e14add3ad4dd5c0034e2e};
        inspect roughly lines 83--118 and 298--337;
  \item \pathcode{mamba_ssm/ops/triton/mamba3/mamba3_siso_combined.py},\\[-2pt]
        {\ttfamily\footnotesize 410ad57555b5841fba7bd76fdafdc3ddf82741176489c3e3493f5c6f3c0c6012};
        inspect roughly lines 90--105 and 247--259;
  \item \pathcode{mamba_ssm/ops/triton/mamba3/mamba3_siso_fwd.py},\\[-2pt]
        {\ttfamily\footnotesize 68d8d0dba0e3913f2e996862603289fd0d7157e7dc57c872e04234be8bb82695};
        inspect roughly lines 256--332, 348--352, and 386--425; and
  \item \pathcode{mamba_ssm/ops/triton/mamba3/mamba3_siso_bwd.py},\\[-2pt]
        {\ttfamily\footnotesize aa1b0aa8f8f7950503adf55220d3feebfea2fe006dc418b5a4e75e1e6842eadc};
        inspect roughly lines 978--1051 and 1475--1550.
\end{itemize}
The hashes, package version, paths, and upstream commit are the durable identity; line numbers are
navigation hints.  Repository Mamba-3 SISO dispatch goes through the FLA wrapper.  Its executable
baseline uses the patched \pathcode{hdn_mamba} environment; the local compatibility/decay patches
do not alter the angle scan or adjacent-pair rotation formulas.  The unpatched \pathcode{hdn} FLA
wrapper is not the runnable decay-correct baseline and was not used to infer those formulas.

\paragraph{Qualification evidence and present status.}
The lower-tail rejection is recorded in
\pathcode{/checkpoints/ngocbh/hdn/diagnostics/rdiag_bounded_clock/probe_v1_1774423/probe_trace.json}.
The redesigned E2 decision is recorded in
\pathcode{/checkpoints/ngocbh/hdn/diagnostics/rdiag_bounded_clock/qualification_bounded_de_20260828_1145/D_E_DECISION.json}.
The precision experiment is recorded in
\pathcode{/lustre-storage/checkpoints/ngocbh/hdn/diagnostics/rdiag_bounded_clock/radial_clock_precision_1776573/result.json};
focused VJP evidence is in
\pathcode{/checkpoints/ngocbh/hdn/ckpts/logs/rdiag_vjp_diag_k64_1776509.log}, and the five-case
composed check is in
\pathcode{/checkpoints/ngocbh/hdn/ckpts/logs/rdiag_angle_order_1776788.log}.  These artifacts support
the specific debugging observations in Section~\ref{sec:lessons}.

The sealed implementation snapshot used by the one-H200 gate is
\pathcode{/lustre-storage/checkpoints/ngocbh/hdn/diagnostics/rdiag_bounded_clock/snapshots/production_e2_20260828_2344}.
Its \pathcode{SOURCE_MANIFEST.json} SHA-256 is
\begin{quote}
\ttfamily\footnotesize
05db76db18b8f4fe170a102e2e5df28bb54a4e36ca7ebe4565a1d813b44530a3.
\end{quote}
The authoritative qualification log is
\pathcode{/checkpoints/ngocbh/hdn/ckpts/logs/rdiag_h200_1777777.log}.  SLURM job 1777777 completed
with status \texttt{COMPLETED}, exit code \texttt{0:0}, 117 passing core tests, and 54 passing E2
tests.  This closes the one-H200 kernel/layer gate.  Distributed training-system checks and the
multi-node soak/attestation remain before submission; training completion and the full
post-training evaluation are subsequent empirical deliverables.

\end{document}
